% =====================================================================
% kl_analytical_reference.tex
% A clean, self-contained re-derivation of ALL analytical results in
% Kekre & Lenel (2024, AER), "The Flight to Safety and International
% Risk Sharing" — Section II of the paper and Appendix B of the online
% appendix. Faithful re-derivation (not an extension, not a critique).
%
% Produced for the big_cy project (task T1b / Assignment 02) as a durable
% reference for the KL extension (T1c) and the quantitative work.
%
% Build:  pdflatex kl_analytical_reference  (x2)
% Sources used:
%   - Main paper:  Kekre_Lenel_2024_AER_...pdf  (Section I model, Section II results)
%   - Online appendix (proofs):  scratch/kl_sources/kl_online_appendix.pdf
% =====================================================================
\documentclass[11pt]{article}

\usepackage[margin=1in]{geometry}
\usepackage{amsmath,amssymb,amsthm,mathtools}
\usepackage{bm}
\usepackage{booktabs}
\usepackage{longtable}
\usepackage{enumitem}
\usepackage{microtype}
\usepackage[colorlinks=true,linkcolor=blue,citecolor=blue]{hyperref}

% =====================================================================
% macros.tex — notation for the KL analytical-results reference document
% Notation tracks Kekre & Lenel (2024, AER) and its online appendix
% exactly. The symbol table in the front matter is the authoritative map;
% these macros enforce consistency across sections.
% =====================================================================

% ---- expectation, deviation, difference operators ----
\newcommand{\E}{\mathbb{E}}                 % expectation
\newcommand{\Et}{\mathbb{E}_t}              % time-t expectation
\newcommand{\Eone}{\mathbb{E}_1}            % period-1 expectation (B.1 timing)
\newcommand{\Ezero}{\mathbb{E}_0}           % period-0 expectation
\renewcommand{\hat}[1]{\widehat{#1}}        % log/level deviation from det. SS
\newcommand{\dev}[1]{\widehat{#1}}          % alias for a deviation
\newcommand{\D}{\Delta}                     % first difference
\newcommand{\Cov}{\operatorname{Cov}}
\newcommand{\Var}{\operatorname{Var}}
\newcommand{\diag}{\operatorname{diag}}
\newcommand{\sign}{\operatorname{sign}}

% ---- country / starred convention (asterisk = Foreign / RoW) ----
\newcommand{\fr}[1]{#1^{*}}                  % Foreign counterpart of #1
\newcommand{\zs}{\zeta^{*}}                  % Foreign measure of households

% ---- preference primitives ----
\newcommand{\hbias}{\varsigma}              % home-bias parameter (KL: \varsigma)
\newcommand{\tel}{\sigma}                    % trade elasticity
\newcommand{\rra}{\gamma}                     % Home risk aversion
\newcommand{\rraF}{\gamma^{*}}               % Foreign risk aversion
\newcommand{\ies}{\psi}                       % IES / labor-comp. parameter
\newcommand{\disc}{\beta}                     % Home discount factor
\newcommand{\discF}{\beta^{*}}               % Foreign discount factor
\newcommand{\frisch}{\nu}                     % Frisch parameter
\newcommand{\labdis}{\bar{\nu}}              % labor-disutility scale
\newcommand{\Lphi}{\Phi}                      % labor-disutility function \Phi(\ell)
\newcommand{\bondu}{\Omega}                   % bond-in-utility function \Omega(B/P)

% ---- production / supply side ----
\newcommand{\cshare}{\alpha}                 % capital share
\newcommand{\dep}{\delta}                     % depreciation
\newcommand{\lel}{\epsilon}                   % labor-variety elasticity
\newcommand{\wadj}{\chi^{W}}                 % wage adjustment-cost scale
\newcommand{\kadj}{\chi^{x}}                 % capital adjustment-cost scale
\newcommand{\taylor}{\phi}                    % Taylor-rule inflation coefficient

% ---- quantities & prices (levels) ----
\newcommand{\val}{v}                          % recursive value
\newcommand{\ce}{ce}                          % certainty equivalent
\newcommand{\pk}{m}                           % real pricing kernel m_{t,t+1}
\newcommand{\tot}{s}                          % terms of trade s_t = E P_H / P*_F
\newcommand{\rer}{q}                          % real exchange rate q_t = E P / P*
\newcommand{\qk}{q^{k}}                       % real price of capital Q^k/P
\newcommand{\divr}{\pi}                       % real dividend \Pi/P
\newcommand{\rw}{w}                           % real wage W/P

% ---- bonds, capital, convenience yield ----
\newcommand{\cy}{\omega}                      % equilibrium convenience yield \omega_t
\newcommand{\cyd}{\omega^{d}}                % latent demand shock \omega^d_t
\newcommand{\edem}{\epsilon^{d}}             % elasticity of bond demand to \Omega'
\newcommand{\inom}{i}                          % nominal rate on safe dollar bonds
\newcommand{\ioth}{\iota}                      % nominal rate on other dollar bonds

% ---- returns ----
\newcommand{\rr}{r}                            % real return on dollar bonds
\newcommand{\rk}{r^{k}}                        % real return on capital
% Foreign real bond return is \fr{r}; capital-in-Foreign return r^{k*} via \fr{\rk}

% ---- aggregates ----
\newcommand{\sav}{a}                           % real aggregate saving a_t
\newcommand{\nfa}{\mathit{nfa}}               % real net foreign assets
\newcommand{\wsh}{\theta}                      % Home financial wealth share
\newcommand{\out}{y}                           % output

% ---- rescaled (stationarized) variables: tilde ----
\newcommand{\rs}[1]{\widetilde{#1}}           % rescaled by z_t

% ---- composite coefficients that recur in the proofs ----
\newcommand{\atil}{\tilde{\alpha}}            % \tilde\alpha (eq. 75-78)
\newcommand{\lwedge}{\tau}                     % steady-state labor wedge
\newcommand{\Gam}{\Gamma}                      % \Gamma (eq. 99)
\newcommand{\Gamp}{\Gamma'}                    % \Gamma' (Props 4-5 proof)

% ---- helper: a flagged reconstruction / gap note ----
\newcommand{\gapflag}[1]{\medskip\noindent\textbf{[Reconstruction flag]}\,\emph{#1}\medskip}


% ---- theorem-like environments ----
\theoremstyle{plain}
\newtheorem{proposition}{Proposition}
\newtheorem{lemma}{Lemma}
\newtheorem{corollary}{Corollary}
\theoremstyle{definition}
\newtheorem{definition}{Definition}
\newtheorem{remark}{Remark}
\newtheorem{assumption}{Assumption}

\allowdisplaybreaks
\numberwithin{equation}{section}

\title{\textbf{The Flight to Safety and International Risk Sharing:\\
A Complete Re-Derivation of the Analytical Results}\\[4pt]
\large A faithful, self-contained reference for Kekre \& Lenel (2024, AER),\\
Section II and online Appendix B}
\author{big\_cy project --- task T1b (theory reference)}
\date{\today}

\begin{document}
\maketitle

\begin{abstract}
\noindent This document re-derives, with complete algebra and full proofs, every
analytical result in Kekre and Lenel (2024, AER), ``The Flight to Safety and
International Risk Sharing.'' We restate the simplified environment (their
Definition~1), the household problems, first-order conditions, market clearing,
the recursive (re-scaled) equilibrium, and the Devereux--Sutherland (2011)
second-order portfolio method; we then prove Propositions~1--5, the
efficient-risk-sharing / law-of-one-price relation, and the ``specialness of safe
dollar bonds'' arguments. Each result is followed by a discussion of its economic
intuition, the load-bearing assumptions, and its connection to the paper's three
target empirical facts (the dollar's negative beta; the United States as the
world's insurer; the reserve-currency paradox). The aim is twofold: a durable
reference for the big\_cy extension, and a complete account of \emph{exactly} what
drives each result, with nothing taken on faith. Steps where the published
appendix is terse and the derivation had to be reconstructed are flagged
explicitly with \textbf{[Reconstruction flag]}.
\end{abstract}

\tableofcontents
\newpage

% ---- front matter: symbol table + paper-to-document map ----
% =====================================================================
% 00_notation.tex — symbol table + paper-to-document map (WP-A)
% =====================================================================
\section{Notation and paper-to-document map}
\label{sec:notation}

This section fixes notation for the entire document. We track Kekre and Lenel
(2024, AER) --- hereafter KL --- and its online appendix exactly; where we
introduce shorthand the table below records the correspondence. Throughout,
``their eq.~(\(n\))'' refers to the equation numbered \(n\) in the published
paper, and ``their app.~eq.~(\(n\))'' to the equation numbered \(n\) in the
online appendix. The macros in \texttt{macros.tex} enforce the symbol choices
recorded here.

\subsection{Conventions}

\begin{itemize}[leftmargin=1.4em,itemsep=2pt]
  \item \textbf{Asterisks denote Foreign / Rest-of-World (RoW).} Any symbol
  \(x^{*}\) is the Foreign counterpart of the Home object \(x\). Home is the
  United States and its nominal unit of account is the dollar.
  \item \textbf{Lower case denotes real variables.} With the sole exceptions of
  the nominal interest rates \(i_t,i_t^{*}\) (and the ``other-bond'' rate
  \(\ioth_t\)), lower-case symbols are real, i.e. denominated in (or deflated
  by) the relevant consumption-good price index. Thus \(\qk_t\equiv Q^k_t/P_t\),
  \(\divr_t\equiv\Pi_t/P_t\), \(\rw_t\equiv W_t/P_t\), \(b_{Ht,s}\equiv
  B_{Ht,s}/P_t\), etc. (their app.\ A.5.1).
  \item \textbf{Hats denote first-order deviations from the deterministic steady
  state.} \(\dev{x}\) is the log deviation for a strictly positive variable and
  the level deviation (scaled as in KL) otherwise; all propositions are stated
  to first order. The macro \(\backslash\)\texttt{dev} and the redefined
  \(\backslash\)\texttt{hat} both produce \(\widehat{\,\cdot\,}\).
  \item \textbf{Tildes denote stationarized (re-scaled) variables.} A tilde
  \(\rs{x}\) divides the corresponding real (or real-deflated nominal) variable
  by the unit-root global productivity level \(z_t\), \emph{labor excepted}
  (their app.\ A.6). Lagged stocks carry a double tilde
  \(\widetilde{\widetilde{x}}_{t-1}\) when a second deflation by the
  productivity \emph{growth} term \(\exp(\sigma^z\epsilon^z_t+\varphi_t)\) is
  applied, again following KL.
  \item \textbf{Timing.} Variables dated \(t\) are end-of-period
  (their app.\ A.4, fn.~53). In the impulse-response analysis (their app.\ B.1)
  the economy is in steady state in \emph{period \(0\)}, all shocks hit in
  \emph{period \(1\)}, and there are no shocks from period \(2\) onward; hence
  \(\Eone[\cdot]\) is the period-\(1\) (post-shock) conditional expectation and
  \(\Ezero\) the steady-state expectation.
  \item \textbf{Difference operator.} \(\D x_{t+1}\equiv x_{t+1}-x_t\) (the
  macro \(\backslash\)\texttt{D} produces \(\Delta\)); e.g.\
  \(\D\dev{P}_t\) is the (log) inflation deviation.
\end{itemize}

\subsection{Symbol table}

The categories below cover every symbol used downstream. The final column gives
the KL reference (main-paper equation, online-appendix equation, or
section/page).

\begin{longtable}{@{}p{2.55cm} p{2.0cm} p{7.2cm} p{2.4cm}@{}}
\toprule
\textbf{Category} & \textbf{Symbol} & \textbf{Meaning} & \textbf{KL ref.} \\
\midrule
\endfirsthead
\toprule
\textbf{Category} & \textbf{Symbol} & \textbf{Meaning} & \textbf{KL ref.} \\
\midrule
\endhead
\midrule \multicolumn{4}{r}{\emph{continued on next page}}\\
\endfoot
\bottomrule
\endlastfoot
% --- countries / measures ---
Countries        & Home, Foreign            & Two countries; Home $=$ US, unit of account $=$ dollar                         & Sec.\ I \\
                 & $\zs=\zeta^{*}$          & Measure of Foreign households (Home has measure one)                            & Sec.\ I \\
% --- preference primitives ---
Preferences      & $\val_t=v_t$             & Recursive (Epstein--Zin) value at Home                                          & eq.\ (1) \\
                 & $c_t$                    & Home consumption (CES aggregate)                                               & eqs.\ (1)--(2) \\
                 & $c_{Ht},c_{Ft}$          & Home consumption of Home- and Foreign-produced goods                            & eq.\ (2) \\
                 & $\ell_t$                 & Home labor (continuum of varieties $j$, $\ell_t=\int_0^1\ell_t(j)\,dj$)         & p.\ 1656 \\
                 & $\disc=\beta$            & Home discount factor                                                            & eq.\ (1) \\
                 & $\ies=\psi$              & IES; also governs consumption--labor complementarity                            & eq.\ (1) \\
                 & $\rra=\gamma$            & Home coefficient of relative risk aversion                                      & eq.\ (1) \\
                 & $\rraF=\gamma^{*}$       & Foreign coefficient of relative risk aversion                                   & p.\ 1656 \\
                 & $\Lphi(\ell_t)=\Phi(\ell_t)$ & Labor-disutility multiplier (Shimer 2010; Trabandt--Uhlig 2011)            & eq.\ (3) \\
                 & $\bondu(B_{Ht,s}/P_t)=\Omega(\cdot)$ & Bond-in-utility (convenience) function                            & eqs.\ (1),(15)--(16) \\
                 & $\frisch=\nu$            & Frisch-elasticity parameter                                                     & eq.\ (3) \\
                 & $\labdis=\bar\nu$        & Scale of labor disutility (Foreign: $\bar\nu^{*}$)                              & eq.\ (3) \\
% --- trade / production parameters ---
Trade/prod.\     & $\hbias=\varsigma$       & Home-bias parameter                                                             & eq.\ (2) \\
                 & $\tel=\sigma$            & Trade elasticity (CES across Home/Foreign goods)                                & eq.\ (2) \\
                 & $\cshare=\alpha$         & Capital share in production                                                      & eq.\ (7) \\
                 & $\dep=\delta$            & Depreciation rate of capital                                                    & p.\ 1656 \\
                 & $\lel=\epsilon$          & Elasticity of substitution across labor varieties                               & eq.\ (6) \\
                 & $\wadj=\chi^{W}$         & Rotemberg wage-adjustment-cost scale                                            & eq.\ (5) \\
                 & $\kadj=\chi^{x}$         & Capital-adjustment-cost scale (global capital producer)                         & eqs.\ (8)--(9) \\
% --- assets ---
Assets (level)   & $B_{Ht,s}$               & Home holding of one-period \emph{safe} dollar bonds (rate $i_t$)                & eq.\ (4) \\
                 & $B_{Ht,o}$               & Home holding of one-period \emph{other} dollar bonds (rate $\ioth_t$)           & eq.\ (4) \\
                 & $B_{Ft}$                 & Home holding of one-period Foreign nominal bonds (rate $i_t^{*}$)               & eq.\ (4) \\
                 & $k_t$                    & Home holding of capital (price $Q^k_t$, dividend $\Pi_{t+1}$, deprec.\ $\delta$)& eq.\ (4) \\
                 & $B^g_{Ht,s}$             & Government supply of safe dollar bonds ($-B^g_{Ht,s}$ is net borrowing)         & eqs.\ (12)--(13),(16) \\
                 & $B_{Ht},B_{Ht}^{*}$      & Total Home/Foreign net dollar-bond positions, return $1+\ioth_t$               & p.\ 1661 \\
% --- convenience yield ---
Convenience      & $\cy_t=\omega_t$         & Equilibrium convenience yield on safe dollar bonds                              & eqs.\ (14)--(17) \\
                 & $\cyd_t=\omega^d_t$      & Latent demand (safety) shock driving $\omega_t$                                 & eqs.\ (16)--(17),(20) \\
                 & $\edem=\epsilon^d$       & Elasticity of bond demand to the nonpecuniary value $\Omega'$                   & eqs.\ (16)--(17) \\
% --- nominal rates / prices ---
Nominal rates    & $\inom_t=i_t$            & Nominal rate on safe dollar bonds (Foreign: $i_t^{*}$)                          & eqs.\ (4),(10) \\
                 & $\ioth_t=\iota_t$        & Nominal rate on other dollar bonds                                              & eqs.\ (4),(14) \\
Prices (level)   & $P_{Ht},P^{*}_{Ft}$      & Prices of Home-/Foreign-produced goods in their own units of account           & eq.\ (4) \\
                 & $E_t$                    & Nominal exchange rate (Foreign units of account per dollar)                     & eq.\ (4) \\
                 & $P_t,P^{*}_t$            & Home / Foreign ideal CPI                                                        & eq.\ (11) \\
                 & $Q^k_t$                  & Nominal price of capital                                                        & eqs.\ (4),(9) \\
                 & $\Pi_t$                  & Nominal dividend per unit capital                                               & eqs.\ (4),(7) \\
                 & $W_t$                    & Nominal wage (variety $j$: $W_t(j)$)                                            & eqs.\ (4)--(6) \\
                 & $\tot=s_t$               & Terms of trade, $s_t\equiv E_tP_{Ht}/P^{*}_{Ft}$                                & app.\ A.5.1 \\
                 & $\rer=q_t$               & Real exchange rate, $q_t\equiv E_tP_t/P^{*}_t$                                  & app.\ A.4 \\
% --- real (lower-case) variables ---
Real variables   & $b_{Ht,s}$               & Real safe dollar bonds, $B_{Ht,s}/P_t$ (sim.\ $b_{Ht,o},b^g_{Ht,s},b_{Ft}$)    & app.\ A.5.1 \\
                 & $\qk_t=q^k_t$            & Real price of capital, $Q^k_t/P_t$                                              & app.\ A.5.1 \\
                 & $\divr_t=\pi_t$          & Real dividend, $\Pi_t/P_t$                                                      & app.\ A.5.1,A.5.3 \\
                 & $\rw_t=w_t$              & Real wage, $W_t/P_t$                                                            & app.\ A.5.1 \\
                 & $\kappa_t$               & Capital \emph{used} in Home production (rented internationally)                 & eqs.\ (7),(63) \\
% --- returns / kernel ---
Returns          & $\rr_{t+1}=r_{t+1}$      & Real return on dollar (CPI-deflated) bonds, $1+r_{t+1}=(1+i_t)P_t/P_{t+1}$      & app.\ A.4,(67) \\
                 & $\fr{r}_{t+1}=r^{*}_{t+1}$ & Real return on Foreign bonds, $1+r^{*}_{t+1}=(1+i^{*}_t)P^{*}_t/P^{*}_{t+1}$  & app.\ A.4,(68) \\
                 & $\rk_{t+1}=r^k_{t+1}$    & Real return on capital, in Home consumption units                              & app.\ A.4,(69) \\
                 & $\pk_{t,t+1}=m_{t,t+1}$  & Home real (EZ) pricing kernel                                                   & app.\ A.5.1 \\
                 & $\ce_t=ce_t$            & Home certainty equivalent of continuation value                                & app.\ A.5.1,(34) \\
% --- aggregates ---
Aggregates       & $\sav_t=a_t$            & Real value of Home aggregate saving                                            & app.\ A.4 \\
                 & $\nfa_t$                 & Real Home net foreign assets                                                    & app.\ A.4 \\
                 & $\wsh_t=\theta_t$        & Home global financial-wealth share (incl.\ seigniorage)                         & app.\ A.6,(51) \\
                 & $\out_t=y_t$             & Home output, $y_t=(z_t\ell_t)^{1-\alpha}\kappa_t^{\alpha}$                      & app.\ A.4 \\
% --- rescaled ---
Rescaled         & $\rs{x}$                 & Stationarized variable, $x/z_t$ (labor excepted)                               & app.\ A.6,(33)--(72) \\
                 & $\widetilde{\widetilde{x}}_{t-1}$ & Lagged stock, further deflated by growth $\exp(\sigma^z\epsilon^z_t+\varphi_t)$ & app.\ A.6 \\
% --- driving forces ---
Driving forces   & $z_t$                    & Global productivity (unit root with disasters), $\log z_t=\log z_{t-1}+\sigma^z\epsilon^z_t+\varphi_t$ & eq.\ (18) \\
                 & $z_{Ft}$                 & Foreign relative productivity, AR(1) in logs                                   & eq.\ (21) \\
                 & $p_t$                    & Disaster probability (AR(1) in logs); disaster size $\varphi_t$                 & eqs.\ (18)--(19) \\
                 & $\varphi_t$              & Realized log disaster ($0$ w.p.\ $1-p_t$, $\underline\varphi<0$ w.p.\ $p_t$)    & eqs.\ (18)--(19) \\
                 & $\cyd_t=\omega^d_t$      & Safety (convenience-demand) shock, $\omega^d_t=\Delta^\omega+\bar\omega^d_t$    & eq.\ (20) \\
% --- operators ---
Operators        & $\dev{x},\widehat{x}$    & First-order deviation from deterministic steady state                           & Sec.\ II \\
                 & $\D x$                   & First difference, $\D x_{t+1}=x_{t+1}-x_t$                                       & Sec.\ II \\
                 & $\Et,\Eone,\Ezero$       & Time-$t$ / period-$1$ / steady-state conditional expectation                    & Sec.\ II; app.\ B.1 \\
% --- composite proof coefficients ---
Proof coeff.\    & $\atil=\bar\alpha$       & Composite coefficient in the period-$\ge 2$ dynamics                            & app.\ eqs.\ (73)--(78) \\
                 & $\lwedge=\tau$           & Steady-state labor wedge                                                        & Prop.\ 4--5 (app.\ B) \\
                 & $\Gam=\Gamma$            & Composite coefficient in the portfolio/risk-premium algebra                     & app.\ B.2 \\
                 & $\Gamp=\Gamma'$          & Auxiliary composite coefficient (Props.\ 4--5 proof)                            & app.\ B.2 \\
                 & $\Lambda$                & Recurring denominator $\tfrac{\alpha}{1-\alpha}+\sigma+(1-\sigma)(\varsigma\tfrac{1+\zeta^*}{\zeta^*})^2$ (our shorthand; KL do not name it) & defined at eq.~in \S\ref{sec:engine} \\
                 & $\Delta$                 & Symmetric-portfolio normalization ($=1$ at symmetric portfolios; Prop.\ 4 proof) & app.\ B.3.2 \\
                 & $\Delta^{\omega}$        & Mean-shift parameter normalizing $\E\omega_t=0$                                 & eq.\ (20) \\
\end{longtable}

\medskip
\noindent\textbf{A note on $\bar\alpha$ vs.\ $\bar\nu$ macros.} The macro
\(\backslash\)\texttt{atil} renders \(\tilde\alpha\) to match KL's online
appendix (their app.\ eqs.\ (73)--(78), where the composite is written
\(\bar\alpha\) in some editions and \(\tilde\alpha\) in the macro). We use
\(\bar\alpha\) in the prose and \(\tilde\alpha\) when the macro is invoked; they
denote the same object. \gapflag{KL's published appendix writes the composite as
$\tilde\alpha$ in their eq.\ (74)'s defining display but refers to it as
$\bar\alpha$ in the text of B.1.1. We treat the two as identical and flag the
notational drift so the critic can confirm no second coefficient is intended.}


% ---- 1. simplified environment (Definition 1) ----
% =====================================================================
% 01_environment.tex — the simplified environment (Definition 1) (WP-A)
% =====================================================================
\section{The simplified environment}
\label{sec:environment}

The analytical results of KL's Section~II are derived in a \emph{simplified
environment}: a set of parametric restrictions on the full model of Section~I
that admits closed-form first-order results while preserving the two essential
ingredients --- a time-varying convenience yield on safe dollar bonds and
cross-country heterogeneity in risk-bearing capacity. We state these
restrictions exactly as KL's Definition~1 (their p.~1660), and after each item
note which step of the analysis it buys, following KL's own discussion on
p.~1660.

\begin{definition}[Simplified environment; KL Definition~1]
\label{def:simplified}
The simplified environment features:
\begin{enumerate}[label=(\roman*),leftmargin=2.0em,itemsep=3pt]
  \item \textbf{Flexible wages, or wages set one period in advance.} In place of
  the Rotemberg (1982) wage-adjustment cost of the full model (their eq.~(5),
  scale $\wadj$), wages are either fully flexible or predetermined by one
  period. \\[2pt]
  \emph{What it buys:} this replaces the forward-looking wage Phillips curve
  (their app.\ eq.~(52)) by a static or one-period-lagged wage equation,
  simplifying the dynamics so that the impulse responses can be solved in closed
  form. The two cases (flexible vs.\ one-period-in-advance) correspond exactly
  to the two bullet groups in Propositions~1--2.
  \item \textbf{Fixed global capital stock} $\bigl(\kadj\!\to\!\infty,\
  \dep\!\to\!0\bigr)$. Infinite capital-adjustment costs freeze the global
  capital stock $\bar k_t$, and zero depreciation removes investment dynamics. \\[2pt]
  \emph{What it buys:} together with (i), this eliminates the endogenous
  capital-accumulation state and the investment Euler equation as a source of
  dynamics, so the only ``real'' margins left are the international allocation of
  a fixed capital stock and the portfolio. It also makes capital a pure
  store-of-value claim whose quantity is fixed.
  \item \textbf{Unitary IES} $(\ies=1)$, \textbf{complete home bias}
  $\bigl(\hbias\!\to\!\zs/(1+\zs)\bigr)$, and \textbf{infinite Frisch
  elasticity} $(\frisch\!\to\!0)$. \\[2pt]
  \emph{What it buys:} $\ies=1$ collapses the Epstein--Zin time aggregator to the
  Cobb--Douglas / log case, so that intertemporal substitution drops out of the
  consumption Euler equation's drift and the value-function recursion
  log-linearizes cleanly. Complete home bias
  $\hbias\!\to\!\zs/(1+\zs)$ sends the CES weight on the foreign good in
  consumption to zero (each country consumes only its own good in the
  deterministic steady state), which decouples the steady-state real exchange
  rate from the consumption shares and is what makes the role of home bias
  \emph{transparent} in the final results. Infinite Frisch ($\frisch\!\to\!0$ in
  the parameterization of their eq.~(3)) makes the labor-disutility multiplier
  $\Phi(\ell)$ linear in the relevant sense, simplifying the supply block.
  \item \textbf{No disaster risk} $(p=0,\ \sigma^p=0)$, \textbf{constant relative
  productivity} $(\sigma^F=0)$, and \textbf{transitory safety}
  $(\rho^\omega=0)$. \\[2pt]
  \emph{What it buys:} switching off disasters and Foreign-productivity shocks
  leaves \emph{two} aggregate shocks --- global productivity $\epsilon^z$ and the
  safety shock $\epsilon^\omega$ --- which is exactly the count needed for the
  ``locally complete'' three-asset structure used in the portfolio analysis
  (Proposition~4). Transitory safety ($\rho^\omega=0$) means a safety shock dies
  out after one period, which is what licenses the period-$1$/period-$\ge 2$
  decomposition of the impulse responses (their app.\ B.1).
  \item \textbf{Identical per-capita wealth across countries in the
  deterministic steady state.} \\[2pt]
  \emph{What it buys:} this pins down a symmetric deterministic steady state
  around which the perturbation is taken, ensuring the steady state is
  well-defined and that the linearized system has the symmetric structure used
  throughout the proofs.
  \item \textbf{Identical discount factors} $(\disc=\discF)$. \\[2pt]
  \emph{What it buys:} with $\disc=\discF$ the two countries share the same
  steady-state interest rate and (given (v)) the same steady-state wealth share
  $\theta=1/(1+\zs)$, removing a level asymmetry that would otherwise complicate
  the steady state. (In the calibrated model KL allow $\disc\neq\discF$ to match
  the level of net foreign assets; here it is shut off.)
\end{enumerate}
\end{definition}

\begin{remark}[General home bias in the proofs; KL app.\ B.1]
\label{rem:general-sigma}
A crucial nuance for what follows: although Definition~\ref{def:simplified}
imposes \emph{complete} home bias $\hbias\!\to\!\zs/(1+\zs)$, KL carry out the
impulse-response and portfolio proofs of Appendix~B for a \emph{general}
home-bias parameter $\hbias$, and only specialize to complete home bias
$\hbias\!\to\!\zs/(1+\zs)$ at the very end, ``so that the role of home bias is
clear'' (their app.\ B.1, opening paragraph). Accordingly, the re-scaled
equilibrium system (Section~\ref{sec:setup} below) and the period-$1$
log-linearization (introduced in the impulse-response engine, WP-B) retain
$\hbias$ as a free parameter; the limit $\hbias\!\to\!\zs/(1+\zs)$ is taken only
where a proposition's stated form requires it. We preserve this convention
throughout this reference: \emph{$\hbias$ is general until explicitly
specialized}.
\end{remark}

\begin{remark}[Perturbation, not the global solution]
The simplified environment is studied with a perturbation (first-order for the
impulse responses, second-order for portfolios) around the deterministic steady
state guaranteed by (v)--(vi). KL emphasize that \emph{none} of the assumptions
in Definition~\ref{def:simplified} is imposed in the quantitative sections~III--V,
which instead solve the full model globally (their app.\ A.7). This reference
concerns only the analytical (simplified-environment) results.
\end{remark}


% ---- 2. setup: households, FOCs, market clearing, equilibrium, rescaling, DS method ----
% =====================================================================
% 02_setup.tex — households, optimality, market clearing, equilibrium (WP-A)
% =====================================================================
\section{Setup: households, optimality, market clearing, equilibrium}
\label{sec:setup}

This section transcribes the building blocks of the model needed for the proofs:
the Home household problem, the convenience-yield wedge, the supply side and
policy, the household first-order conditions and pricing kernel, the auxiliary
variable definitions, the re-scaled (stationary) recursive equilibrium, and the
Devereux--Sutherland portfolio method. We follow KL's equation numbering
throughout (``their eq.'' for the main paper, ``their app.\ eq.'' for the online
appendix) so that the proof work packages can cross-reference. This is setup,
not yet a proof: where an object is stated rather than derived, that is by
design.

\subsection{The Home household problem}
\label{sec:hh-problem}

The representative Home household has recursive (Epstein--Zin) preferences over
consumption $c_t$, labor $\ell_t$, and the real value of safe dollar bonds
$B_{Ht,s}/P_t$ (their eq.~(1)):
\begin{equation}
\label{eq:ez-value}
\val_t=\left\{(1-\disc)\Bigl[c_t\,\Lphi(\ell_t)\,
\bondu_t\!\bigl(B_{Ht,s}/P_t\bigr)\Bigr]^{1-1/\ies}
+\disc\,\Et\!\Bigl[\val_{t+1}^{1-\rra}\Bigr]^{\frac{1-1/\ies}{1-\rra}}\right\}^{\frac{1}{1-1/\ies}}.
\end{equation}
Consumption is a CES aggregate of Home- and Foreign-produced goods with
home-bias weight $\hbias$ and trade elasticity $\tel$ (their eq.~(2)):
\begin{equation}
\label{eq:ces-c}
c_t=\left[\Bigl(\tfrac{1}{1+\zs}+\hbias\Bigr)^{\!\frac1\tel}c_{Ht}^{\frac{\tel-1}{\tel}}
+\Bigl(\tfrac{\zs}{1+\zs}-\hbias\Bigr)^{\!\frac1\tel}c_{Ft}^{\frac{\tel-1}{\tel}}\right]^{\frac{\tel}{\tel-1}}.
\end{equation}
The disutility of labor follows Shimer (2010) and Trabandt--Uhlig (2011)
(their eq.~(3)):
\begin{equation}
\label{eq:labor-phi}
\Lphi(\ell_t)=\left[1+(1/\ies-1)\,\labdis\,
\frac{(\ell_t)^{1+1/\frisch}}{1+1/\frisch}\right]^{\frac{1/\ies}{1-1/\ies}}.
\end{equation}
The function $\bondu_t(\cdot)$ delivers the nonpecuniary (liquidity/safety) value
of safe dollar bonds, in the spirit of money-in-the-utility (Sidrauski 1967) and
the convenience-yield literature (Krishnamurthy--Vissing-Jorgensen 2012); its
functional form is given in \eqref{eq:Omega-form} below. The household supplies a
continuum of labor varieties $j\in[0,1]$, $\ell_t=\int_0^1\ell_t(j)\,dj$.

Each period the household chooses consumption and a portfolio
$\{B_{Ht,s},B_{Ht,o},B_{Ft},k_t\}$ --- safe dollar bonds (nominal rate $i_t$),
other dollar bonds (rate $\ioth_t$), Foreign nominal bonds (rate $i_t^{*}$), and
capital (price $Q^k_t$, dividend $\Pi_{t+1}$, depreciation $\delta$) --- subject
to the resource constraint (their eq.~(4)):
\begin{align}
\label{eq:budget}
P_{Ht}c_{Ht}+E_t^{-1}P^{*}_{Ft}c_{Ft}
&+B_{Ht,s}+B_{Ht,o}+E_t^{-1}B_{Ft}+Q^k_t k_t \notag\\
\le\;(1+i_{t-1})B_{Ht-1,s}&+(1+\ioth_{t-1})B_{Ht-1,o}+E_t^{-1}(1+i^{*}_{t-1})B_{Ft-1}\notag\\
&+\Bigl[\Pi_t+(1-\dep)Q^k_t\Bigr]k_{t-1}\exp(\varphi_t)
+\int_0^1 W_t(j)\ell_t(j)\,dj\notag\\
&\quad-\int_0^1 AC^W_t(j)\,dj+T_t,
\end{align}
where $P_{Ht},P^{*}_{Ft}$ are Home/Foreign producer prices in their own units of
account, $E_t$ is the nominal exchange rate (Foreign units per dollar), the
disaster term $\exp(\varphi_t)$ scales the capital stock, $T_t$ is a government
transfer, and producer-currency pricing is assumed so the law of one price
holds. The Rotemberg wage-adjustment cost is (their eq.~(5))
\begin{equation}
\label{eq:rotemberg}
AC^W_t(j)=\frac{\wadj}{2}\,W_t\ell_t
\left[\frac{W_t(j)}{W_{t-1}(j)\exp(\varphi_t)}-1\right]^2,
\end{equation}
rebated lump-sum. In the simplified environment of
Definition~\ref{def:simplified}, \eqref{eq:rotemberg} is replaced by flexible or
one-period-predetermined wages.

\paragraph{Foreign analog.} Foreign households (measure $\zs$) solve an
analogous problem with risk aversion $\rraF$ that may differ from $\rra$, the
same $\ies,\hbias,\tel,\frisch,\wadj$, and possibly a different discount factor
$\discF$ and labor-disutility scale $\bar\nu^{*}$. They too receive nonpecuniary
utility $\bondu^{*}_t\bigl(B^{*}_{Ht,s}/(E_t^{-1}P^{*}_t)\bigr)$ from safe dollar
bonds. All starred objects are the Foreign counterparts of the Home objects
above.

\subsection{The convenience yield}
\label{sec:cy}

Because safe dollar bonds and other dollar bonds pay in the same unit of account
and are both risk-free, investor indifference at Home requires that the
\emph{effective} return on safe bonds equal the return on other bonds
(their eq.~(14)):
\begin{equation}
\label{eq:wedge}
\frac{1+i_t}{\,1-c_t\,\bondu_t'(B_{Ht,s}/P_t)/\bondu_t(B_{Ht,s}/P_t)\,}=1+\ioth_t.
\end{equation}
The denominator's second term is the marginal nonpecuniary value of safe bonds.
Equating this nonpecuniary value across Home and Foreign agents on the margin
defines the (common) convenience yield $\cy_t$ (their eq.~(15)):
\begin{equation}
\label{eq:omega-def}
\cy_t\equiv c_t\,\frac{\bondu_t'(B_{Ht,s}/P_t)}{\bondu_t(B_{Ht,s}/P_t)}
=c_t^{*}\,\frac{\bondu_t^{*}{}'\!\bigl(B^{*}_{Ht,s}/(E_t^{-1}P^{*}_t)\bigr)}
{\bondu_t^{*}\!\bigl(B^{*}_{Ht,s}/(E_t^{-1}P^{*}_t)\bigr)}.
\end{equation}
KL adopt the convenient functional form (their displayed equation on p.~1659):
\begin{equation}
\label{eq:Omega-form}
\bondu_t\!\left(\frac{B_{Ht,s}}{P_t}\right)
=\exp\!\left(\omega^d_t\frac{B_{Ht,s}}{P_t\,\bar c_t}
-\frac12\frac1{\edem}\Bigl(\frac{B_{Ht,s}}{P_t\,\bar c_t}\Bigr)^{2}
-\Bigl[\omega^d_t\frac{\bar B_{Ht,s}}{P_t\,\bar c_t}
-\frac12\frac1{\edem}\Bigl(\frac{\bar B_{Ht,s}}{P_t\,\bar c_t}\Bigr)^{2}\Bigr]\right),
\end{equation}
where $\omega^d_t$ is an exogenous driving force (the \emph{safety shock}),
$\edem$ a parameter, and bars denote aggregates the household takes as given. The
first two terms inside the exponential make the nonpecuniary value
time-varying; the bracketed terms ensure $\bondu_t(B_{Ht,s}/P_t)=1$ in
equilibrium, so the convenience yield has no mechanical effect on the
household's position through the level of $\bondu_t$. Given the analogous Foreign
form, the second equality in \eqref{eq:omega-def} together with market clearing
in safe dollar bonds implies (their eq.~(16))
\begin{equation}
\label{eq:safe-bond-clearing}
\frac{B_{Ht,s}}{P_t c_t}=\frac{B^{*}_{Ht,s}}{E_t^{-1}P^{*}_t c^{*}_t}
=\frac{(-B^g_{Ht,s})}{P_t c_t+\zs E_t^{-1}P^{*}_t c^{*}_t},
\end{equation}
and the first equality in \eqref{eq:omega-def} then yields the
\emph{reduced-form} convenience yield (their eq.~(17))
\begin{equation}
\label{eq:omega-reduced}
\cy_t=\omega^d_t-\frac1{\edem}\,
\frac{(-B^g_{Ht,s})}{P_t c_t+\zs E_t^{-1}P^{*}_t c^{*}_t}.
\end{equation}
Intuitively, $\cy_t$ rises with private demand for safe dollar bonds
$\omega^d_t$ (latent demand) and falls with public supply $-B^g_{Ht,s}$; the
relative strength of supply is governed by $\edem$, the elasticity of demand to
the nonpecuniary value. Under the fiscal rule \eqref{eq:fiscal-supply} below the
supply term is proportional to global consumption, so $\cy_t$ inherits the
stochastic properties of the safety shock $\omega^d_t$ and is effectively
exogenous.

\subsection{Supply side and policy}
\label{sec:supply-policy}

\paragraph{Producers (their app.\ A.5.3).} The representative Home producer hires
$\ell_t$ from the domestic labor packer and rents $\kappa_t$ of capital on the
international market, producing with productivity $z_t$ and capital share
$\alpha$ (their eq.~(7)). Its first-order conditions are
\begin{equation}
\label{eq:prod-foc-home}
\rw_t=\frac{P_{Ht}}{P_t}(1-\alpha)z_t^{1-\alpha}\ell_t^{-\alpha}\kappa_t^{\alpha},
\qquad
\divr_t=\frac{P_{Ht}}{P_t}\alpha z_t^{1-\alpha}\ell_t^{1-\alpha}\kappa_t^{\alpha-1}.
\end{equation}
The representative Foreign producer's FOCs are analogous, with relative
productivity $z_{Ft}$ and the real exchange rate $q_t$ converting units:
\begin{equation}
\label{eq:prod-foc-foreign}
\rw_t^{*}=q_t^{-1}\frac{P^{*}_{Ft}}{P^{*}_t}(1-\alpha)(z_t z_{Ft})^{1-\alpha}
(\zs\ell_t^{*})^{-\alpha}\kappa_t^{*\alpha},\qquad
\divr_t=q_t^{-1}\frac{P^{*}_{Ft}}{P^{*}_t}\alpha(z_t z_{Ft})^{1-\alpha}
(\zs\ell_t^{*})^{1-\alpha}\kappa_t^{*\alpha-1}.
\end{equation}
Because capital is rented on a single international market, its dividend
$\divr_t$ (in dollars) is common across countries.

\paragraph{Unions (their app.\ A.5.2).} Each Home union $j$ sets the wage
$W_t(j)$ and supplies labor $\ell_t(j)$ to maximize members' utilitarian
welfare subject to the Rotemberg cost \eqref{eq:rotemberg}. The resulting wage
first-order condition is
\begin{equation}
\label{eq:union-foc}
\begin{aligned}
\rw_t-\frac{c_t\Lphi'(\ell_t)}{\Lphi(\ell_t)}
+\rw_t\frac{\wadj}{\lel}\Bigg[
&\frac{\rw_t}{\rw_{t-1}\exp(\varphi_t)}\frac{P_t}{P_{t-1}}
\left(\frac{\rw_t}{\rw_{t-1}\exp(\varphi_t)}\frac{P_t}{P_{t-1}}-1\right)\\
&-\Et m_{t,t+1}\left(\frac{\rw_{t+1}}{\rw_t\exp(\varphi_{t+1})}\right)^{2}
\frac{P_{t+1}}{P_t}\frac{\ell_{t+1}}{\ell_t}
\left(\frac{\rw_{t+1}}{\rw_t\exp(\varphi_{t+1})}\frac{P_{t+1}}{P_t}-1\right)
\Bigg]=0,
\end{aligned}
\end{equation}
with an analogous condition in Foreign. In the simplified environment this
collapses to a flexible or one-period-predetermined wage equation.

\paragraph{Global capital producer (their eqs.~(8)--(9)).} A representative
global capital producer combines Home and Foreign consumption goods (without home
bias) to produce new capital with adjustment-cost scale $\chi^x$; profit
maximization yields the real price of capital
\begin{equation}
\label{eq:qk}
\qk_t=\left(\frac{\bar k_t}{\bar k_{t-1}\exp(\varphi_t)}\right)^{\chi^x}
\left[\Bigl(\tfrac{1}{1+\zs}+\hbias\Bigr)\Bigl(\tfrac{P_{Ht}}{P_t}\Bigr)^{1-\tel}
+\Bigl(\tfrac{\zs}{1+\zs}-\hbias\Bigr)\Bigl(\tfrac{P^{*}_{Ft}}{P_t}\Bigr)^{1-\tel}\right]^{\frac1{1-\tel}}.
\end{equation}
In the simplified environment $\chi^x\!\to\!\infty$ with $\bar k_t$ fixed.

\paragraph{Monetary policy (their eqs.~(10)--(11)).} A Taylor (1993) rule
governs the Home nominal rate,
\begin{equation}
\label{eq:taylor}
1+i_t=(1+\bar\imath)\left(\frac{P_t}{P_{t-1}}\right)^{\taylor},
\end{equation}
where the ideal CPI is
\begin{equation}
\label{eq:cpi}
P_t=\left[\Bigl(\tfrac{1}{1+\zs}+\hbias\Bigr)P_{Ht}^{1-\tel}
+\Bigl(\tfrac{\zs}{1+\zs}-\hbias\Bigr)\bigl(E_t^{-1}P^{*}_{Ft}\bigr)^{1-\tel}\right]^{\frac1{1-\tel}}.
\end{equation}
An analogous CPI-targeting rule with the same coefficient $\taylor$ determines
$i_t^{*}$ in Foreign.

\paragraph{Fiscal policy (their eqs.~(12)--(13)).} The Home government maintains
a constant ratio of safe dollar debt to global consumption,
\begin{equation}
\label{eq:fiscal-supply}
-B^g_{Ht,s}=\bar b^g\bigl(P_t c_t+\zs E_t^{-1}P^{*}_t c^{*}_t\bigr),
\end{equation}
with $\bar b^g>0$ the empirically relevant case (the government borrows in safe
dollar bonds, i.e.\ issues Treasuries), and rebates with lump-sum transfers
\begin{equation}
\label{eq:fiscal-transfer}
T_t=\int_0^1 AC^W_t(j)\,dj+(1+i_{t-1})B^g_{Ht-1,s}-B^g_{Ht,s}.
\end{equation}
Ricardian equivalence holds for all government activity other than safe-bond
issuance, which alone provides nonpecuniary services. Combining
\eqref{eq:fiscal-supply} with \eqref{eq:omega-reduced} renders $\cy_t$
proportional to the exogenous safety shock $\omega^d_t$ (up to the
constant $\bar b^g/\edem$).

\subsection{Household first-order conditions and pricing kernel}
\label{sec:hh-foc}

Home intratemporal optimality across Home and Foreign goods is (their app.\
A.5.1)
\begin{equation}
\label{eq:intratemporal}
\frac{c_{Ht}}{c_{Ft}}
=\frac{\tfrac{1}{1+\zs}+\hbias}{\tfrac{\zs}{1+\zs}-\hbias}\,\tot_t^{-\tel},
\qquad \tot_t\equiv\frac{E_t P_{Ht}}{P^{*}_{Ft}}\ \ \text{(terms of trade)}.
\end{equation}
The Home \emph{real pricing kernel} is the Epstein--Zin stochastic discount
factor (their app.\ A.5.1), evaluated using $\bondu_t(B_{Ht,s}/P_t)=1$ in
equilibrium:
\begin{equation}
\label{eq:kernel}
\pk_{t,t+1}=\disc\,\frac{c_t}{c_{t+1}}
\left(\frac{c_{t+1}\Lphi(\ell_{t+1})}{c_t\Lphi(\ell_t)}\right)^{1-1/\ies}
\left(\frac{\val_{t+1}}{\ce_t}\right)^{1/\ies-\rra},
\qquad
\ce_t=\Et\!\Bigl[\val_{t+1}^{1-\rra}\Bigr]^{\frac1{1-\rra}},
\end{equation}
where $\ce_t$ is the certainty equivalent of the continuation value. Intertemporal
optimality is characterized by three Euler equations (their app.\ A.5.1) --- for
safe dollar bonds (adjusted by the convenience yield), Foreign bonds (UIP), and
capital, respectively:
\begin{align}
\label{eq:euler-safe}
1&=\Et\,\pk_{t,t+1}\,\frac{1+\rr_{t+1}}{1-\cy_t},\\
\label{eq:euler-uip}
1&=\Et\,\pk_{t,t+1}\,\frac{\rer_t}{\rer_{t+1}}\bigl(1+\fr{r}_{t+1}\bigr),\\
\label{eq:euler-capital}
1&=\Et\,\pk_{t,t+1}\bigl(1+\rk_{t+1}\bigr).
\end{align}
Equation \eqref{eq:euler-safe} is implied by optimality in \emph{either} safe or
other dollar bonds given the definition of $\cy_t$ in \eqref{eq:omega-def}: the
$1/(1-\cy_t)$ factor converts the safe-bond return into the convenience-adjusted
return that the kernel prices. Analogous starred conditions hold for Foreign
households, using the Foreign kernel $\pk^{*}_{t,t+1}$ (their app.\ eqs.\
(47)--(49)).

\paragraph{Safe-bond market clearing and the real resource constraint.} Combining
the cross-country indifference condition \eqref{eq:omega-def} (which gives
$B_{Ht,s}/(P_t c_t)=B^{*}_{Ht,s}/(E_t^{-1}P^{*}_t c^{*}_t)$) with global market
clearing in safe dollar bonds yields the allocation of safe bonds across
countries (their app.\ A.5.1, restating \eqref{eq:safe-bond-clearing}):
\begin{equation}
\label{eq:safe-bond-split}
B_{Ht,s}=\frac{P_t c_t}{P_t c_t+\zs E_t^{-1}P^{*}_t c^{*}_t}(-B^g_{Ht,s}),
\qquad
B^{*}_{Ht,s}=\frac{E_t^{-1}P^{*}_t c^{*}_t}{P_t c_t+\zs E_t^{-1}P^{*}_t c^{*}_t}(-B^g_{Ht,s}).
\end{equation}
Substituting government transfers \eqref{eq:fiscal-transfer} into the budget
constraint \eqref{eq:budget}, dividing by $P_t$, and denoting real positions
$b_{Ht,s}\equiv B_{Ht,s}/P_t$, $b^g_{Ht,s}\equiv B^g_{Ht,s}/P_t$,
$b_{Ht,o}\equiv B_{Ht,o}/P_t$, $b_{Ft}\equiv B_{Ft}/P_t$, $\qk_t\equiv Q^k_t/P_t$,
$\divr_t\equiv\Pi_t/P_t$, $\rw_t\equiv W_t/P_t$, the Home household's real
resource constraint is (their app.\ A.5.1)
\begin{equation}
\label{eq:real-budget-home}
\begin{aligned}
c_t+b_{Ht,s}+b^g_{Ht,s}+b_{Ht,o}+\rer_t^{-1}b_{Ft}+\qk_t k_t
&=\rw_t\ell_t
+(1+\rr_t)\bigl(b_{Ht-1,s}+b^g_{Ht-1,s}\bigr)\\
&\quad+\left(\frac{1+\rr_t}{1-\cy_{t-1}}\right)b_{Ht-1,o}
+\rer_t^{-1}(1+\fr{r}_t)b_{Ft-1}\\
&\quad+\bigl(\divr_t+(1-\dep)\qk_t\bigr)k_{t-1}\exp(\varphi_t).
\end{aligned}
\end{equation}
KL further define each agent's net dollar-bond position
$b_{Ht}\equiv(1-\cy_t)(b_{Ht,s}+b^g_{Ht,s})+b_{Ht,o}$ and
$b^{*}_{Ht}\equiv(1-\cy_t)b^{*}_{Ht,s}+b^{*}_{Ht,o}$, each earning the
``other-bond'' return $1+\ioth_t=(1+i_t)/(1-\cy_t)$; the composition of the
dollar-bond position matters only for the seigniorage Home earns on the safe debt
held by Foreign (their p.~1661 and app.\ A.5.1).

\subsection{Additional variables (their app.\ A.4)}
\label{sec:aux-vars}

The following auxiliary definitions are used throughout the proofs. The real
exchange rate (an increase is a Home appreciation):
\begin{equation}
\label{eq:rer-def}
\rer_t\equiv\frac{E_t P_t}{P^{*}_t}.
\end{equation}
The real interest rates and the real return on capital (in Home consumption
units):
\begin{equation}
\label{eq:returns-def}
1+\rr_{t+1}\equiv(1+i_t)\frac{P_t}{P_{t+1}},\quad
1+\fr{r}_{t+1}\equiv(1+i^{*}_t)\frac{P^{*}_t}{P^{*}_{t+1}},\quad
1+\rk_{t+1}\equiv\frac{\bigl(\Pi_{t+1}+(1-\dep)Q^k_{t+1}\bigr)}{Q^k_t}\frac{P_t}{P_{t+1}}\exp(\varphi_{t+1}).
\end{equation}
Output, the real value of aggregate saving, and real net foreign assets at Home:
\begin{equation}
\label{eq:agg-def}
\out_t\equiv(z_t\ell_t)^{1-\alpha}(\kappa_t)^{\alpha},\quad
\sav_t\equiv\frac{1}{P_t}\bigl(B_{Ht}+E_t^{-1}B_{Ft}+Q^k_t k_t\bigr),\quad
\nfa_t\equiv\sav_t-\frac{Q^k_t}{P_t}\kappa_{t+1}\exp(-\varphi_{t+1}).
\end{equation}
All variables are dated at the end of the period (their app.\ fn.~53); the factor
$\exp(-\varphi_{t+1})$ in $\nfa_t$ undoes capital destruction so that capital
\emph{used} at Home is compared with capital \emph{owned} by Home residents (their
app.\ fn.~54).

\subsection{The re-scaled (stationary) recursive equilibrium}
\label{sec:rescaled}

Because global productivity $z_t$ has a unit root (their eq.~(18)), KL stationarize
by dividing all real variables (and real-deflated nominal variables) by $z_t$,
\emph{labor excepted} (their app.\ A.6). Write $\rs{x}\equiv x/z_t$; lagged stocks
carry a second deflation by the growth term
$\exp(\sigma^z\epsilon^z_t+\varphi_t)$, denoted with a double tilde
$\widetilde{\widetilde{x}}_{t-1}$. After scaling, global-productivity shocks
(inclusive of disasters) only govern the transition across states. The recursive
equilibrium is the system of their app.\ eqs.~(33)--(72); we transcribe the
economically central conditions and refer to the rest by number.

\paragraph{Home block (their app.\ eqs.~(33)--(41)).} The re-scaled value
recursion, certainty equivalent, CES aggregator, intratemporal allocation, and
kernel are their app.\ eqs.~(33)--(37). The three re-scaled Euler equations are
\begin{align}
\label{eq:reuler-safe}
1&=\Et\,\rs{\pk}_{t,t+1}\exp\!\bigl(-\rra[\sigma^z\epsilon^z_{t+1}+\varphi_{t+1}]\bigr)
\frac{1+\rr_{t+1}}{1-\cy_t}, &&\text{(their app.\ eq.~(38))}\\
\label{eq:reuler-uip}
1&=\Et\,\rs{\pk}_{t,t+1}\exp\!\bigl(-\rra[\sigma^z\epsilon^z_{t+1}+\varphi_{t+1}]\bigr)
\frac{\rer_t}{\rer_{t+1}}(1+\fr{r}_{t+1}), &&\text{(their app.\ eq.~(39))}\\
\label{eq:reuler-cap}
1&=\Et\,\rs{\pk}_{t,t+1}\exp\!\bigl(-\rra[\sigma^z\epsilon^z_{t+1}+\varphi_{t+1}]\bigr)
(1+\rk_{t+1}), &&\text{(their app.\ eq.~(40))}
\end{align}
and the re-scaled Home resource constraint is
\begin{equation}
\label{eq:rbudget-home}
\rs{c}_t+\rs{b}_{Ht}+\rer_t^{-1}\rs{b}_{Ft}+\qk_t\rs{k}_t
=\rs{w}_t\ell_t+\theta_t\bigl(\divr_t+(1-\dep)\qk_t\bigr)\widetilde{\widetilde{k}}_{t-1},
\quad\text{(their app.\ eq.~(41))}
\end{equation}
where $\theta_t$ is the Home wealth share. The Foreign block (their app.\
eqs.~(42)--(50)) is symmetric, with risk aversion $\rraF$, Euler equations
\begin{align}
\label{eq:reuler-safe-f}
1&=\Et\,\rs{\pk}^{*}_{t,t+1}\exp\!\bigl(-\rraF[\sigma^z\epsilon^z_{t+1}+\varphi_{t+1}]\bigr)
\frac{\rer_{t+1}}{\rer_t}\frac{1+\rr_{t+1}}{1+\cy_t}, &&\text{(their app.\ eq.~(47))}\\
\label{eq:reuler-uip-f}
1&=\Et\,\rs{\pk}^{*}_{t,t+1}\exp\!\bigl(-\rraF[\sigma^z\epsilon^z_{t+1}+\varphi_{t+1}]\bigr)
(1+\fr{r}_{t+1}), &&\text{(their app.\ eq.~(48))}\\
\label{eq:reuler-cap-f}
1&=\Et\,\rs{\pk}^{*}_{t,t+1}\exp\!\bigl(-\rraF[\sigma^z\epsilon^z_{t+1}+\varphi_{t+1}]\bigr)
\frac{\rer_{t+1}}{\rer_t}(1+\rk_{t+1}), &&\text{(their app.\ eq.~(49))}
\end{align}
and the Foreign resource constraint (their app.\ eq.~(50)).

\noindent\emph{(Note: these match KL's app.\ eqs.\ (47)--(49) verbatim, verified
against their appendix p.~64. Observe the Home/Foreign asymmetry in the
convenience adjustment: the Home safe-bond Euler equation~\eqref{eq:reuler-safe}
carries $1/(1-\cy_t)$, whereas the Foreign safe-bond Euler
equation~\eqref{eq:reuler-safe-f} carries $1/(1+\cy_t)$, because the convenience
benefit accrues to the issuer's residents and is netted out of the price Foreign
agents are willing to pay. The $q_{t+1}/q_t$ placement is exactly as in KL.)}

\paragraph{Wealth-share law of motion (their app.\ eq.~(51)).} The global wealth
share of Home households, inclusive of seigniorage, evolves as
\begin{equation}
\label{eq:wealth-share}
\theta_{t+1}=\frac{1}{(\divr_{t+1}+(1-\dep)\qk_{t+1})\widetilde{\widetilde{k}}_t}
\left[\frac{1+\rr_{t+1}}{1-\cy_t}\rs{b}_{Ht}
+\frac{1}{\rer_{t+1}}(1+\fr{r}_{t+1})\rs{b}_{Ft}
+(\divr_{t+1}+(1-\dep)\qk_{t+1})\widetilde{\widetilde{k}}_t
-\cy_{t+1}\frac{\zs\rer^{-1}_{t+1}\rs{c}^{*}_{t+1}}{\rs{c}_{t+1}+\zs\rer^{-1}_{t+1}\rs{c}^{*}_{t+1}}\rs{b}^g_{Ht+1,s}\right].
\end{equation}
The last term is the seigniorage Home earns on the share of its safe debt held by
Foreign.

\paragraph{Supply, market clearing, returns, prices (their app.\ eqs.~(52)--(72)).}
Supply-side optimality (re-scaled union and producer FOCs) is their app.\
eqs.~(52)--(60); goods-, factor-, and asset-market clearing
\begin{align}
\label{eq:mc-goods}
&\rs{c}_{Ht}+\zs\rs{c}^{*}_{Ht}+\Bigl(\tfrac{\widetilde{\widetilde{k}}_t}{\widetilde{\widetilde{k}}_{t-1}}\Bigr)^{\!\chi^x}\rs{x}_{Ht}=(\ell_t)^{1-\alpha}(\rs\kappa_t)^{\alpha}, &&\text{(their app.\ eq.~(61))}\\
&\rs{c}_{Ft}+\zs\rs{c}^{*}_{Ft}+\Bigl(\tfrac{\widetilde{\widetilde{k}}_t}{\widetilde{\widetilde{k}}_{t-1}}\Bigr)^{\!\chi^x}\rs{x}_{Ft}=(z_{Ft}\zs\ell^{*}_t)^{1-\alpha}(\rs\kappa^{*}_t)^{\alpha}, &&\text{(their app.\ eq.~(62))}\\
&\rs\kappa_t+\rs\kappa^{*}_t=\widetilde{\widetilde{k}}_{t-1}, \quad
\rs{k}_t+\zs\rs{k}^{*}_t=\widetilde{\widetilde{k}}_t,\quad
(1-\dep)\widetilde{\widetilde{k}}_{t-1}+\rs{x}_t=\widetilde{\widetilde{k}}_t, &&\text{(their app.\ eqs.~(63)--(65))}\\
\label{eq:mc-bonds}
&\rs{b}_{Ht}+\zs\rs{b}^{*}_{Ht}=0, &&\text{(their app.\ eq.~(66))}
\end{align}
the return definitions
\begin{equation}
\label{eq:rescaled-returns}
1+\rr_{t+1}=(1+i_t)\tfrac{P_t}{P_{t+1}},\quad
1+\fr{r}_{t+1}=(1+i^{*}_t)\tfrac{P^{*}_t}{P^{*}_{t+1}},\quad
1+\rk_{t+1}=\tfrac{\divr_{t+1}+(1-\dep)\qk_{t+1}}{\qk_t}\exp(\varphi_{t+1}),
\end{equation}
(their app.\ eqs.~(67)--(69)), and the price relations linking the CPI, terms of
trade, and real exchange rate (their app.\ eqs.~(70)--(72)):
\begin{equation}
\label{eq:price-relations}
\frac{P_t}{P_{Ht}}=\Bigl[\bigl(\tfrac{1}{1+\zs}+\hbias\bigr)+\bigl(\tfrac{\zs}{1+\zs}-\hbias\bigr)\tot_t^{\tel-1}\Bigr]^{\frac1{1-\tel}},
\quad
\rer_t=\tot_t\left[\frac{(\tfrac{1}{1+\zs}+\hbias)+(\tfrac{\zs}{1+\zs}-\hbias)\tot_t^{\tel-1}}{(\tfrac{1}{1+\zs}-\tfrac{\hbias}{\zs})\tot_t^{1-\tel}+(\tfrac{\zs}{1+\zs}+\tfrac{\hbias}{\zs})}\right]^{\frac1{1-\tel}}.
\end{equation}

\paragraph{Equilibrium and state vector.} Together with the Taylor and fiscal
rules and the driving forces, their app.\ eqs.~(33)--(72) define the recursive
equilibrium in the sense of the following definition.

\begin{definition}[Equilibrium; KL app.\ A.3 Definition~1]
\label{def:equilibrium}
An equilibrium is a sequence of prices and policies such that: each Home
representative household chooses $\{c_{Ht},c_{Ft},B_{Ht,s},B_{Ht,o},B_{Ft},k_t\}$
to maximize \eqref{eq:ez-value} subject to \eqref{eq:ces-c}--\eqref{eq:rotemberg}
(and analogously in Foreign); each Home union $j$ chooses $\{W_t(j),\ell_t(j)\}$
to maximize members' utilitarian welfare subject to \eqref{eq:rotemberg}
(analogously in Foreign); the Home labor packer chooses $\{\ell_t(j)\}$ to
maximize profits (their eq.~(6), analogously in Foreign); the Home producer
chooses $\{\ell_t,\kappa_t\}$ to maximize profits (their eq.~(7), analogously in
Foreign); the global capital producer chooses $\{x_{Ht},x_{Ft},x_t\}$ to maximize
profits (their eq.~(9)) subject to \eqref{eq:qk}; the Home government sets
$B^g_{Ht,s}$ by \eqref{eq:fiscal-supply} and $\{i_t,\{T_t\}\}$ by
\eqref{eq:taylor} and \eqref{eq:fiscal-transfer} (and the Foreign government the
latter); and goods, factor, and asset markets clear (their app.\ eqs.~(23)--(32),
including $B_{Ft}+\zs B^{*}_{Ft}=0$, their app.\ eq.~(32)).
\end{definition}

After re-scaling, the aggregate state in period $t$ is seven-dimensional:
\begin{equation}
\label{eq:state}
\bigl\{\,p_t,\ \cy_t,\ z_{Ft},\ \theta_t,\ \widetilde{\widetilde{k}}_{t-1},\
\rs{w}_{t-1},\ \rs{w}^{*}_{t-1}\,\bigr\},
\end{equation}
i.e.\ disaster probability, convenience yield, Foreign relative productivity,
Home wealth share, lagged (re-scaled) global capital, and lagged (re-scaled) Home
and Foreign real wages. By Walras' law the Foreign bond market clears given the
rest. In the simplified environment, $p$ and $z_F$ are constant
(Definition~\ref{def:simplified}(iv)) and the wage states are trivial under
flexible or one-period-predetermined wages, so the effective state collapses to
$\{\cy,\theta\}$ around the steady state --- the two variables that carry the
analytical dynamics.

\subsection{The Devereux--Sutherland portfolio method}
\label{sec:ds-method}

Equilibrium \emph{portfolios} --- the steady-state holdings of dollar bonds,
capital, and Foreign bonds that each country chooses --- cannot be read off the
first-order (linearized) equilibrium conditions alone. KL pin them down with the
second-order method of Devereux and Sutherland (2011); we lay out the logic here
and defer the algebra to the portfolio work package (WP-D).

To first order, the three asset returns
$\{1+\rr_{t+1},\,1+\rk_{t+1},\,(\rer_t/\rer_{t+1})(1+\fr{r}_{t+1})\}$ enter the
Euler equations \eqref{eq:euler-safe}--\eqref{eq:euler-capital} only through
their conditional means, which the Euler equations equate. Hence to first order
all assets are \emph{perfect substitutes} and the portfolio weights are
\emph{indeterminate}: any portfolio is consistent with the first-order
equilibrium. The Devereux--Sutherland insight is that portfolios are pinned down
by a \emph{second-order} object. Specifically, the steady-state portfolio is the
unique one for which the first-order \emph{excess-return} (risk-sharing)
conditions --- the differences of the Euler equations
\eqref{eq:euler-safe}--\eqref{eq:euler-capital} across assets --- hold once the
portfolio Euler equations are approximated to \emph{second} order. Equivalently,
the portfolio is chosen so that the conditional covariances between the pricing
kernel and excess returns (which are second-order objects) are consistent with
the first-order dynamics of consumption and the real exchange rate that those
same portfolios induce.

Two structural facts make this tractable and economically sharp in the simplified
environment. First, there are exactly \emph{three} traded assets (safe dollar
bonds, capital, Foreign bonds). Second, by Definition~\ref{def:simplified}(iv)
the economy is subject to exactly \emph{two} aggregate shocks (global
productivity $\epsilon^z$ and safety $\epsilon^\omega$). With one more asset than
shock, the asset span is \emph{locally complete} in the sense of Coeurdacier and
Gourinchas (2016): around the deterministic steady state the available assets
suffice to implement the constrained-efficient (complete-markets) risk-sharing
allocation. Consequently the equilibrium portfolios that the
Devereux--Sutherland method selects are precisely those that implement efficient
risk sharing to first order. This is the formal content behind KL's statement
(their p.~1664) that ``the three available assets implement efficient risk
sharing around the steady state.'' The closed-form characterization of these
portfolios and of the resulting risk premia (their Propositions~4--5) is carried
out in WP-D; here we only fix the method and its scope.

\gapflag{KL state local completeness and efficient risk sharing as a consequence
of the ``one more asset than shock'' count (their p.~1664, citing Coeurdacier and
Gourinchas 2016) but do not, in the main text, prove that the
Devereux--Sutherland-selected portfolio coincides with the efficient one in this
specific three-asset/two-shock setting. We state the equivalence here as KL
assert it; the explicit verification (that the span $\{$bond, capital, FX$\}$
covers the two shocks and that no redundancy or rank-deficiency arises at the
symmetric steady state) belongs to WP-D and is flagged for the critic.}


% ---- 3. impulse-response engine (Appendix B.1) ----
% =====================================================================
% 03_impulse_engine.tex — Appendix B.1: the period-1 log-linearization (WP-B)
% =====================================================================
\section{The impulse-response engine (period-1 log-linearization)}
\label{sec:engine}

This section reconstructs KL's online Appendix~B.1, which reduces the entire
impulse response of the simplified economy to a \emph{static} system in
period~1. The output is a small set of reduced-form expressions --- for the
terms of trade $\dev{\tot}_1$, consumptions $\dev{\rs c}_1,\fr{\dev{\rs c}}_1$,
the expected wealth-share change $\Eone\dev{\wsh}_2$, net foreign assets
$\widehat{\nfa}_1$, the expected capital return $\Eone\rk_2$, the inflation
deviations $\D\dev P_1,\D\fr{\dev P}_1$, and the realized returns
$\rr_1,\fr{r}_1,\rk_1$ --- all written in terms of
$(\dev{\ell}_1,\fr{\dev{\ell}}_1,\dev{\cy}_1,\widetilde{\widetilde{\dev k}}_1)$
and the deep parameters. These are the \emph{engine} that
Propositions~\ref{prop:prices}, \ref{prop:returns}, and~\ref{prop:wealth}
specialize. Throughout we keep the home-bias parameter $\hbias$ general (as KL
do in B.1), so the role of complete home bias $\hbias\to\zs/(1+\zs)$ is visible
when we later impose it. The labels in this section (e.g.\
\eqref{eq:s1-solved}, \eqref{eq:Er2k-solved}) are the cross-references used by
the proof sections.

\subsection{Timing and strategy}
\label{sec:engine-timing}

KL's device (their app.~B.1, first paragraph) is the following. The economy sits
in the deterministic steady state in period~$0$. \emph{All} shocks ---
to the safety/convenience demand $\dev{\cy}_1$ and to productivity $\dev{\epsilon}^z_1$ ---
hit in period~$1$. There are \emph{no} shocks from period~$2$ onward. Two
consequences make the problem static:
\begin{enumerate}[leftmargin=1.6em,itemsep=2pt]
  \item \textbf{Forward dynamics collapse.} With no innovations after period~$1$,
  every period-$\ge2$ variable is a deterministic function of the two
  predetermined states entering period~$2$: the expected Home wealth share
  $\Eone\dev{\wsh}_2$ and the (twice-deflated) lagged global capital stock
  $\widetilde{\widetilde{\dev k}}_1$. The period-$\ge2$ equilibrium conditions
  can then be solved \emph{in closed form} for the conditional means
  $\Eone\dev{\rs c}_2$, $\Eone\fr{\dev{\rs c}}_2$, $\Eone\dev{\tot}_2$,
  $\Eone\dev{\rer}_2$, and the continuation values --- this is
  Subsection~\ref{sec:engine-period2}.
  \item \textbf{Period~$1$ becomes a static system.} Substituting those
  closed-form forward objects into the period-$1$ optimality, market-clearing,
  and policy conditions leaves a finite linear system in the period-$1$
  unknowns. Solving it is Subsection~\ref{sec:engine-period1}.
\end{enumerate}
Because the predetermined capital state enters only as
$\widetilde{\widetilde{\dev k}}_1$ and is, on impact, equal to its steady-state
value (capital is chosen one period in advance), we set
$\widetilde{\widetilde{\dev k}}_1=0$ when evaluating the \emph{impact} response;
we nonetheless carry it symbolically through the engine because the reusable
reduced forms (and Propositions~2--3) reference it.

\subsection{Dynamics from period~2 onward}
\label{sec:engine-period2}

Since there are no shocks from period~$2$ onward, under the parametric
restrictions of Definition~\ref{def:simplified} the period-$\ge2$ equilibrium
conditions (their app.~eqs.~(33)--(72), specialized) reduce to closed-form
relations between the period-$2$ conditional means and the two predetermined
states. KL record them as their app.~eqs.~(73)--(78).

\paragraph{Expected consumptions.} The Home and Foreign Euler/risk-sharing
conditions, evaluated at the deterministic continuation path, give (their
app.~eqs.~(73)--(74))
\begin{align}
\Eone\dev{\rs c}_2 &= \atil\,\Eone\dev{\wsh}_2 + \cshare\,\widetilde{\widetilde{\dev k}}_1, \label{eq:Ec2}\\
\Eone\fr{\dev{\rs c}}_2 &= -\frac{1}{\zs}\,\atil\,\Eone\dev{\wsh}_2 + \cshare\,\widetilde{\widetilde{\dev k}}_1. \label{eq:Ec2-star}
\end{align}
Here $\cshare\,\widetilde{\widetilde{\dev k}}_1$ is the common scale effect of
the predetermined global capital stock (capital share $\cshare$ on the
production side), and the wealth-share term redistributes consumption across
countries: a unit rise in the Home wealth share $\Eone\dev{\wsh}_2$ raises Home
expected consumption by $\atil$ and lowers Foreign expected (per-capita,
measure-$\zs$) consumption by $\atil/\zs$, so that the population-weighted change
nets out against the resource constraint.

\paragraph{Terms of trade and real exchange rate.} The period-$\ge2$ goods-market
clearing and the intratemporal allocation, linearized, pin the expected terms of
trade to the wealth share (their app.~eq.~(75)):
\begin{equation}
\Eone\dev{\tot}_2 = \frac{\hbias\left(\frac{1+\zs}{\zs}\right)^2}
{\dfrac{\cshare}{1-\cshare}+\tel\left(1-\left(\hbias\,\frac{1+\zs}{\zs}\right)^2\right)}\,
\atil\,\Eone\dev{\wsh}_2, \label{eq:Es2}
\end{equation}
and the log-linearized definition of the real exchange rate (their app.~eq.~(76))
\begin{equation}
\Eone\dev{\rer}_2 = \hbias\,\frac{1+\zs}{\zs}\,\Eone\dev{\tot}_2. \label{eq:Eq2}
\end{equation}
Equation~\eqref{eq:Eq2} is the period-$2$ analog of the period-$1$ real-exchange-rate
relation~\eqref{eq:q1-s1} below; it follows from log-linearizing
$q_t\equiv E_tP_t/P^{*}_t$ together with the CES price indices, which makes the
real exchange rate proportional to the terms of trade with constant of
proportionality $\hbias(1+\zs)/\zs$ (their app.~A.4).

\paragraph{Continuation values / certainty equivalents.} Because there is no
uncertainty after period~$1$, the (re-scaled) continuation value and its
certainty equivalent coincide with expected consumption up to the common labor
and capital terms; KL record (their app.~eqs.~(77)--(78))
\begin{align}
\Eone\dev{\rs \val}_2 &= \Eone\dev{\rs c}_2 = \atil\,\Eone\dev{\wsh}_2 + \cshare\,\widetilde{\widetilde{\dev k}}_1, \label{eq:Ev2}\\
\Eone\fr{\dev{\rs \val}}_2 &= \Eone\fr{\dev{\rs c}}_2 = -\frac{1}{\zs}\,\atil\,\Eone\dev{\wsh}_2 + \cshare\,\widetilde{\widetilde{\dev k}}_1, \label{eq:Ev2-star}
\end{align}
i.e.\ the period-$2$ value deviations inherit \eqref{eq:Ec2}--\eqref{eq:Ec2-star}
because, absent further risk, the certainty-equivalent operator is the identity
to first order (their app.~A.5.1; the Jensen correction is second order and
drops at the linearization).

\paragraph{The composite coefficient $\atil$.} The coefficient appearing in
\eqref{eq:Ec2}--\eqref{eq:Ev2-star} is defined by (their app.~eq.\ following~(78))
\begin{equation}
\atil \;\equiv\; \frac{\cshare}{\,1-\left(\dfrac{\zs}{1+\zs}-\hbias\right)
\dfrac{\hbias\left(\frac{1+\zs}{\zs}\right)^2}
{\dfrac{\cshare}{1-\cshare}+\tel\left(1-\left(\hbias\,\frac{1+\zs}{\zs}\right)^2\right)}\,}. \label{eq:atil-def}
\end{equation}
It is the fixed-point coefficient that internalizes the feedback from the wealth
share to the terms of trade \eqref{eq:Es2} and back into the resource
constraint: the denominator subtracts the product of the net Home expenditure
weight $\bigl(\frac{\zs}{1+\zs}-\hbias\bigr)$ and the terms-of-trade response
multiplier from~\eqref{eq:Es2}. The numerator $\cshare$ is the capital-share
scaling inherited from the production function. We follow the notational note in
Section~\ref{sec:notation}: the macro renders $\atil=\tilde\alpha$, and KL refer
to the same object as $\bar\alpha$ in the B.1.1 prose; they are identical.

\medskip
KL remark that ``these are the only conditions we need to solve for the
equilibrium in period~$1$.'' We now turn to that period-$1$ system.

\subsection{The period-1 log-linearized system}
\label{sec:engine-period1}

We transcribe KL's app.~eqs.~(79)--(98) and then perform the substitutions
(their app.~pp.~18--19) that solve the system. We group the conditions into
three blocks: (i) a relative-price block, (ii) an Euler/return block, and (iii)
a wealth/NFA block.

\subsubsection{Block (i): relative prices, allocation, clearing}

Log-linearizing the definition of the real exchange rate $q_1\equiv E_1P_1/P^{*}_1$
(their app.~eq.~(79)) gives the period-$1$ analog of~\eqref{eq:Eq2}:
\begin{equation}
\dev{\rer}_1 = \hbias\,\frac{1+\zs}{\zs}\,\dev{\tot}_1, \label{eq:q1-s1}
\end{equation}
a relation we use repeatedly. Log-linearizing the intratemporal allocation of
consumption, the equilibrium factor prices, the resource constraints, and
goods-market clearing yields the population-weighted resource constraint (their
app.~eq.~(80))
\begin{equation}
\frac{1}{1+\zs}\dev{\rs c}_1 + \frac{\zs}{1+\zs}\fr{\dev{\rs c}}_1
= (1-\cshare)\left[\frac{1}{1+\zs}\dev{\ell}_1 + \frac{\zs}{1+\zs}\fr{\dev{\ell}}_1\right]
+ \cshare\,\widetilde{\widetilde{\dev k}}_1, \label{eq:rc-agg}
\end{equation}
and the cross-country (relative) goods-market condition (their app.~eq.~(81))
\begin{equation}
\hbias\,\frac{1+\zs}{\zs}\left(\dev{\rs c}_1 - \fr{\dev{\rs c}}_1\right)
= \left[\frac{\cshare}{1-\cshare} + \tel\left(1-\left(\hbias\,\frac{1+\zs}{\zs}\right)^2\right)\right]\dev{\tot}_1
+ \dev{\ell}_1 - \fr{\dev{\ell}}_1. \label{eq:rc-rel}
\end{equation}
The bracket multiplying $\dev{\tot}_1$ in \eqref{eq:rc-rel} is exactly the
denominator appearing in the period-$2$ terms-of-trade relation~\eqref{eq:Es2};
it combines the income (capital-share) channel $\cshare/(1-\cshare)$ with the
substitution channel $\tel(1-(\hbias(1+\zs)/\zs)^2)$ of the trade elasticity.

\subsubsection{Block (ii): Euler equations, returns, Fisher/Taylor}

Log-linearizing the Euler equations (their app.~eqs.~(82)--(85)) gives, for the
\emph{expected} (period-$1$ to period-$2$) changes,
\begin{align}
\D\Eone\dev{\rs c}_2 &= \D\Eone\fr{\dev{\rs c}}_2 - \hbias\,\frac{1+\zs}{\zs}\,\D\Eone\dev{\tot}_2, \label{eq:euler-cc}\\
\D\Eone\dev{\rs c}_2 &= \Eone\rk_2, \label{eq:euler-rk}\\
\Eone\rk_2 &= \Eone\rr_2 + \dev{\cy}_1, \label{eq:euler-r}\\
\Eone\fr{r}_2 &= \Eone\rr_2 + \dev{\cy}_1 + \hbias\,\frac{1+\zs}{\zs}\,\D\Eone\dev{\tot}_2, \label{eq:euler-rstar}
\end{align}
where $\D x_2 \equiv x_2-x_1$ denotes the period-$1$-to-$2$ change and
\eqref{eq:euler-rstar} uses~\eqref{eq:Eq2} (KL note ``where we have used~(76)'').
The crucial object is the safety wedge $\dev{\cy}_1$ in \eqref{eq:euler-r}: in
the household's Euler equation the convenience term enters as the factor
$(1+r_{t+1})/(1-\cy_t)$, so to first order the \emph{expected} real return on
the safe bond falls one-for-one with the convenience yield relative to the
capital return --- the wedge $\Eone\rk_2-\Eone\rr_2=\dev{\cy}_1$ in
\eqref{eq:euler-r} is the linearized statement of exactly this.

Log-linearizing the Fisher equations and Taylor rules (their app.~eqs.~(88)--(91))
gives
\begin{align}
\Eone\rr_2 &= \dev{\inom}_1, & \Eone\fr{r}_2 &= \fr{\dev{\inom}}_1, \label{eq:fisher}\\
\dev{\inom}_1 &= \taylor\,\D\dev{P}_1, & \fr{\dev{\inom}}_1 &= \taylor\,\D\fr{\dev{P}}_1, \label{eq:taylor}
\end{align}
where the Taylor rules implement $\D\dev P_2=\D\fr{\dev P}_2=0$ from period~$2$
on (so all forward inflation is zero) and $\taylor$ is the inflation-response
coefficient. Finally, the realized period-$1$ bond returns are pure inflation
surprises (their app.~eqs.~(95)--(96)):
\begin{equation}
\rr_1 = -\D\dev{P}_1, \qquad \fr{r}_1 = -\D\fr{\dev{P}}_1, \label{eq:realized-bond}
\end{equation}
since the period-$0$ nominal rate is predetermined and the real return on a
nominal bond is minus realized inflation to first order.

\subsubsection{Block (iii): wealth share, NFA, capital returns}

Log-linearizing the \emph{expected} evolution of Home's wealth share (using
equilibrium factor prices and the Home resource constraint) gives (their
app.~eq.~(86))
\begin{equation}
\Eone\dev{\wsh}_2 = \frac{1}{\disc\cshare}\Bigl[(1-\disc)\Bigl(\tfrac{\zs}{1+\zs}-\hbias\Bigr)\dev{\tot}_1
+(1-\disc)(1-\cshare)\dev{\ell}_1
+(1-\disc)\cshare\,\widetilde{\widetilde{\dev k}}_1+\cshare\dev{\wsh}_1
-(1-\disc)\dev{\rs c}_1\Bigr]. \label{eq:Ewsh-raw}
\end{equation}
Linearizing the definition of Home net foreign assets (equilibrium factor
prices, Home resource constraint, and cross-country capital allocation) gives
(their app.~eq.~(87))
\begin{multline}
\widehat{\nfa}_1 = \sav\Biggl[\frac{1}{\disc\cshare}\Bigl[(1-\disc)\Bigl(\tfrac{\zs}{1+\zs}-\hbias\Bigr)\dev{\tot}_1
+(1-\disc)(1-\cshare)\dev{\ell}_1\\
+\cshare\,\widetilde{\widetilde{\dev k}}_1+\cshare\dev{\wsh}_1-(1-\disc)\dev{\rs c}_1\Bigr]
-\frac{\zs}{1+\zs}\frac{1}{1-\cshare}\dev{\tot}_2-\widetilde{\widetilde{\dev k}}_1\Biggr], \label{eq:nfa-raw}
\end{multline}
where $\sav$ is steady-state Home saving. Log-linearizing the \emph{realized}
evolution of Home's wealth share (their app.~eq.~(92)) gives
\begin{equation}
\dev{\wsh}_1 = \Bigl(\frac{\qk\,k}{\sav}-1\Bigr)\bigl(\rk_1-\rr_1\bigr)
+\frac{b_F}{\sav}\bigl(\fr{r}_1-\dev{\rer}_1-\rr_1\bigr)
-\disc\,\frac{\zs}{1+\zs}\frac{b^g_{H,s}}{\sav}\,\dev{\cy}_1, \label{eq:wsh-realized}
\end{equation}
where $q^kk$, $b_F$, $b^g_{H,s}$ are the steady-state Home capital, Foreign-bond,
and government-safe-debt positions; this is the revaluation equation that drives
Proposition~\ref{prop:wealth}. Log-linearizing the realized return on capital
(equilibrium profits; their app.~eq.~(93)) gives
\begin{equation}
\rk_1 = (1-\disc)\left[-\hbias\,\dev{\tot}_1 + (1-\cshare)\Bigl(\tfrac{1}{1+\zs}\dev{\ell}_1+\tfrac{\zs}{1+\zs}\fr{\dev{\ell}}_1\Bigr)
-(1-\cshare)\widetilde{\widetilde{\dev k}}_1\right]+\disc\,\dev{\qk}_1, \label{eq:rk1-raw}
\end{equation}
and the \emph{expected} return on capital (their app.~eq.~(94))
\begin{equation}
\dev{\qk}_1 = -\hbias\,\Eone\dev{\tot}_2 - (1-\cshare)\widetilde{\widetilde{\dev k}}_1 - \Eone\rk_2, \label{eq:qk1-raw}
\end{equation}
which is the linearized capital-pricing (no-arbitrage) condition relating today's
capital price to the future terms of trade, capital stock, and expected return.

\subsubsection{Solving the system: substitution chain}
\label{sec:engine-solve}

We now reproduce KL's substitution chain (their app.~pp.~18--19), which solves
the block above for the reduced forms. The strategy is: (1) eliminate
$\fr{\dev{\rs c}}_1$ using the aggregate resource constraint~\eqref{eq:rc-agg};
(2) substitute into the relative condition~\eqref{eq:rc-rel} to get $\dev{\tot}_1$
in terms of consumptions; (3) use the forward Euler block to express $\dev{\rs c}_1$
in terms of $\Eone\dev{\wsh}_2$ and labor; (4) back-substitute to get closed
forms for $\dev{\tot}_1$, $\dev{\rs c}_1$, $\Eone\dev{\wsh}_2$, $\widehat{\nfa}_1$,
$\Eone\rk_2$, and then prices and returns.

\paragraph{Step 1: eliminate $\fr{\dev{\rs c}}_1$.} Solving the aggregate resource
constraint~\eqref{eq:rc-agg} for $\fr{\dev{\rs c}}_1$ (their app.~p.~18, ``(80)
implies''):
\begin{equation}
\fr{\dev{\rs c}}_1 = -\frac{1}{\zs}\dev{\rs c}_1
+\frac{1+\zs}{\zs}\left[(1-\cshare)\Bigl(\tfrac{1}{1+\zs}\dev{\ell}_1+\tfrac{\zs}{1+\zs}\fr{\dev{\ell}}_1\Bigr)+\cshare\,\widetilde{\widetilde{\dev k}}_1\right]. \label{eq:cstar-solved}
\end{equation}

\paragraph{Step 2: terms of trade in terms of $\dev{\rs c}_1$.} Substituting
\eqref{eq:cstar-solved} into the relative condition~\eqref{eq:rc-rel} and solving
for $\dev{\tot}_1$ (their app.~p.~18, ``Substituting in (81) implies''):
\begin{multline}
\dev{\tot}_1 = \frac{1}{\dfrac{\cshare}{1-\cshare}+\tel\left(1-\left(\hbias\,\frac{1+\zs}{\zs}\right)^2\right)}
\Biggl[\bigl(\fr{\dev{\ell}}_1-\dev{\ell}_1\bigr)\\
+\hbias\left(\frac{1+\zs}{\zs}\right)^2\!\left(\dev{\rs c}_1
-\left[(1-\cshare)\Bigl(\tfrac{1}{1+\zs}\dev{\ell}_1+\tfrac{\zs}{1+\zs}\fr{\dev{\ell}}_1\Bigr)+\cshare\,\widetilde{\widetilde{\dev k}}_1\right]\right)\Biggr]. \label{eq:s1-intermediate}
\end{multline}
\gapflag{KL display \eqref{eq:s1-intermediate} directly. The intermediate step is:
substitute \eqref{eq:cstar-solved} into the left side of \eqref{eq:rc-rel}, so
that $\hbias\frac{1+\zs}{\zs}(\dev{\rs c}_1-\fr{\dev{\rs c}}_1)
=\hbias\frac{1+\zs}{\zs}\bigl(\dev{\rs c}_1+\frac{1}{\zs}\dev{\rs c}_1
-\frac{1+\zs}{\zs}[\cdots]\bigr)
=\hbias(\frac{1+\zs}{\zs})^2\bigl(\dev{\rs c}_1-[\cdots]\bigr)$,
using $1+\frac{1}{\zs}=\frac{1+\zs}{\zs}$; then move
$\dev{\ell}_1-\fr{\dev{\ell}}_1$ to the left, divide by the bracket. We have
expanded that algebra; \eqref{eq:s1-intermediate} matches KL's displayed line.}

\paragraph{Step 3: period-1 Home consumption.} Combining \eqref{eq:cstar-solved}
and \eqref{eq:s1-intermediate} with the forward block
\eqref{eq:Ec2}--\eqref{eq:Es2} and the consumption-Euler
relation~\eqref{eq:euler-cc} (their app.~p.~18, ``Combining these with (73),
(74), (75), and (82) implies'') yields the closed form for $\dev{\rs c}_1$:
\begin{multline}
\dev{\rs c}_1 = \atil\,\Eone\dev{\wsh}_2
+(1-\cshare)\Bigl(\tfrac{1}{1+\zs}\dev{\ell}_1+\tfrac{\zs}{1+\zs}\fr{\dev{\ell}}_1\Bigr)
+\cshare\,\widetilde{\widetilde{\dev k}}_1\\
-\frac{\hbias}{\dfrac{\cshare}{1-\cshare}+\tel+(1-\tel)\left(\hbias\,\frac{1+\zs}{\zs}\right)^2}\bigl(\fr{\dev{\ell}}_1-\dev{\ell}_1\bigr). \label{eq:c1-solved}
\end{multline}
\gapflag{KL state \eqref{eq:c1-solved} as ``Combining these \dots\ implies'' with
the intervening algebra suppressed. The mechanism: the Euler
relation~\eqref{eq:euler-cc} written in levels (period~1) sets
$\dev{\rs c}_1=\Eone\dev{\rs c}_2-(\Eone\dev{\rs c}_2-\dev{\rs c}_1)$; using the
forward terms-of-trade response~\eqref{eq:Es2} and the period-2 consumption
\eqref{eq:Ec2}, the wealth-share and capital terms reproduce
$\atil\Eone\dev{\wsh}_2+\cshare\widetilde{\widetilde{\dev k}}_1$, the labor terms
reproduce $(1-\cshare)(\cdots)$, and the relative-labor term collects into the
final coefficient. The denominator
$\frac{\cshare}{1-\cshare}+\tel+(1-\tel)(\hbias\frac{1+\zs}{\zs})^2$ is the same
bracket as in \eqref{eq:rc-rel}/\eqref{eq:Es2} after using
$\tel(1-x^2)=\tel-\tel x^2$ and adding back the $+x^2$ contributed by the
home-bias share, i.e.\ $\tel(1-x^2)+x^2=\tel+(1-\tel)x^2$ with
$x\equiv\hbias\frac{1+\zs}{\zs}$. We verified this identity; it is the recurring
denominator below.}

For brevity define the recurring composite denominator
\begin{equation}
\Lambda \;\equiv\; \frac{\cshare}{1-\cshare}+\tel+(1-\tel)\left(\hbias\,\frac{1+\zs}{\zs}\right)^2, \label{eq:Lambda-def}
\end{equation}
so that the relative-labor coefficient in \eqref{eq:c1-solved} is
$\hbias/\Lambda$. (KL do not name $\Lambda$; we introduce it purely as
shorthand. It is \emph{not} the same as the bracket in \eqref{eq:rc-rel}, which
is $\frac{\cshare}{1-\cshare}+\tel(1-(\hbias\frac{1+\zs}{\zs})^2)
=\Lambda-(\hbias\frac{1+\zs}{\zs})^2$.)

\paragraph{Step 4: the two solved forms for $\dev{\tot}_1$.} Substituting
\eqref{eq:c1-solved} back into \eqref{eq:s1-intermediate} (their app.~p.~18,
``which we can substitute into the previous result to give'') yields the first
solved form,
\begin{equation}
\dev{\tot}_1 = \frac{1}{\Lambda}\bigl(\fr{\dev{\ell}}_1-\dev{\ell}_1\bigr)
+\frac{\hbias\left(\frac{1+\zs}{\zs}\right)^2}
{\dfrac{\cshare}{1-\cshare}+\tel\left(1-\left(\hbias\,\frac{1+\zs}{\zs}\right)^2\right)}\,\atil\,\Eone\dev{\wsh}_2. \label{eq:s1-solved}
\end{equation}
The first term is the \emph{relative-labor} (supply) channel: if Foreign works
relatively more, Home's terms of trade improve. The second term is the
\emph{wealth-share} (demand) channel, with the identical multiplier as in the
period-$2$ relation~\eqref{eq:Es2}. Using the real-exchange-rate
relation~\eqref{eq:q1-s1}, the equivalent statement for the real exchange rate is
\begin{equation}
\dev{\rer}_1 = \hbias\,\frac{1+\zs}{\zs}\,\dev{\tot}_1
=\frac{\hbias\frac{1+\zs}{\zs}}{\Lambda}\bigl(\fr{\dev{\ell}}_1-\dev{\ell}_1\bigr)
+\hbias\,\frac{1+\zs}{\zs}\,\frac{\hbias\left(\frac{1+\zs}{\zs}\right)^2}
{\dfrac{\cshare}{1-\cshare}+\tel\left(1-\left(\hbias\,\frac{1+\zs}{\zs}\right)^2\right)}\,\atil\,\Eone\dev{\wsh}_2. \label{eq:q1-solved}
\end{equation}
The two displayed forms of $\dev{\tot}_1$ that KL refer to are
\eqref{eq:s1-intermediate} (in terms of consumptions) and \eqref{eq:s1-solved}
(in terms of labor and the wealth share).

\paragraph{Step 5: expected wealth share and NFA.} Substituting
\eqref{eq:c1-solved} and \eqref{eq:s1-solved} into the raw wealth-share
linearization~\eqref{eq:Ewsh-raw} (their app.~p.~18, ``Substituting these
into~(86) implies'') gives
\begin{equation}
\Eone\dev{\wsh}_2 = \dev{\wsh}_1
+(1-\disc)\frac{\zs}{1+\zs}(\tel-1)\frac{1-\cshare}{\cshare}\,
\frac{1-\left(\hbias\,\frac{1+\zs}{\zs}\right)^2}{\Lambda}\bigl(\fr{\dev{\ell}}_1-\dev{\ell}_1\bigr). \label{eq:Ewsh-solved}
\end{equation}
\gapflag{KL display \eqref{eq:Ewsh-solved}. The reconstruction: substituting
\eqref{eq:c1-solved} for $\dev{\rs c}_1$ and \eqref{eq:s1-solved} for
$\dev{\tot}_1$ into \eqref{eq:Ewsh-raw}, the $\atil\Eone\dev{\wsh}_2$ pieces
cancel against the $\frac{\cshare}{\disc\cshare}\dev{\wsh}_1$ normalization
\emph{except} for the relative-labor term, and the
$\bigl(\frac{\zs}{1+\zs}-\hbias\bigr)$-weighted terms-of-trade term combines with
the $-(1-\disc)\dev{\rs c}_1$ term. The factor $(\tel-1)$ arises because the
wealth-share revaluation depends on whether expenditure-switching
(elasticity $\tel$) over- or under-compensates the relative price move; at
$\tel=1$ (Cobb--Douglas trade) the channel vanishes. We reproduce KL's
coefficient; the $\frac{1-(\hbias(1+\zs)/\zs)^2}{\Lambda}$ factor is the ratio of
the \eqref{eq:rc-rel}-bracket-residual to $\Lambda$ that survives after the
$\atil$ cancellation. The full term-by-term cancellation is mechanical but
lengthy; we have checked the coefficient against KL's displayed line and flag
that we did not re-expand every one of the $\sim10$ intermediate terms.}
The net-foreign-assets reduced form follows by the identical substitution into
\eqref{eq:nfa-raw} (their app.~pp.~18--19):
\begin{equation}
\widehat{\nfa}_1 = \sav\Biggl[\Eone\dev{\wsh}_2
-\frac{\zs}{1+\zs}\frac{1}{1-\cshare}\dev{\tot}_2\Biggr]
=\sav\Eone\dev{\wsh}_2 - \sav\frac{\zs}{1+\zs}\frac{1}{1-\cshare}\Eone\dev{\tot}_2, \label{eq:nfa-solved}
\end{equation}
where $\Eone\dev{\wsh}_2$ is \eqref{eq:Ewsh-solved} and $\Eone\dev{\tot}_2$ is
\eqref{eq:Es2}; KL write the bracket with the $\frac{\zs}{1+\zs}\frac{1}{1-\cshare}\dev{\tot}_2$
term carried to the next page. The economic content: Home NFA moves with its
wealth share, net of a valuation term in the future terms of trade.

\paragraph{Step 6: expected capital return.} Substituting the solved forms into
the capital-Euler relation~\eqref{eq:euler-rk}--\eqref{eq:euler-r} and using the
forward consumption~\eqref{eq:Ec2} (their app.~p.~19, ``substituting these
into~(83) and making use of~(73) implies'') gives
\begin{equation}
\Eone\rk_2 = -(1-\cshare)\Bigl(\tfrac{1}{1+\zs}\dev{\ell}_1+\tfrac{\zs}{1+\zs}\fr{\dev{\ell}}_1\Bigr)
+\frac{\hbias}{\Lambda}\bigl(\fr{\dev{\ell}}_1-\dev{\ell}_1\bigr). \label{eq:Er2k-solved}
\end{equation}
This is the key supply-side object: the expected capital return falls when global
employment rises (the first term, via diminishing returns / lower future
marginal product) and responds to the relative-labor gap through the same
$\hbias/\Lambda$ coefficient as $\dev{\rs c}_1$.

\paragraph{Step 7: inflation and realized returns.} The Fisher/Taylor
block~\eqref{eq:fisher}--\eqref{eq:taylor} gives the period-$1$ inflation
deviations. From $\Eone\rk_2=\Eone\rr_2+\dev{\cy}_1$ \eqref{eq:euler-r},
$\Eone\rr_2=\dev{\inom}_1$ \eqref{eq:fisher}, and $\dev{\inom}_1=\taylor\D\dev P_1$
\eqref{eq:taylor}, the Home Taylor rule implies (their app.~p.~19)
\begin{equation}
\D\dev{P}_1 = \frac{1}{\taylor}\bigl(\Eone\rk_2-\dev{\cy}_1\bigr), \label{eq:dP1}
\end{equation}
while for Foreign, using \eqref{eq:euler-rstar} with~\eqref{eq:Es2} for
$\D\Eone\dev{\tot}_2$,
\begin{equation}
\D\fr{\dev{P}}_1 = \frac{1}{\taylor}\left(\Eone\rk_2 - \frac{1+\zs}{\zs}\,
\frac{\hbias}{\Lambda}\bigl(\fr{\dev{\ell}}_1-\dev{\ell}_1\bigr)\right). \label{eq:dPstar1}
\end{equation}
\gapflag{For \eqref{eq:dPstar1} we substituted \eqref{eq:Es2} into the Foreign
Euler~\eqref{eq:euler-rstar}: $\fr{\dev{\inom}}_1=\Eone\fr r_2=\Eone\rk_2
-\dev{\cy}_1+\dev{\cy}_1+\hbias\frac{1+\zs}{\zs}\D\Eone\dev{\tot}_2$, and the
term $\hbias\frac{1+\zs}{\zs}\D\Eone\dev{\tot}_2$ evaluates, via the
$\Eone\dev{\tot}_2$ relation and the $\dev{\tot}_1$ solution~\eqref{eq:s1-solved},
to $\frac{1+\zs}{\zs}\frac{\hbias}{\Lambda}(\fr{\dev{\ell}}_1-\dev{\ell}_1)$ on
impact (the wealth-share piece of $\D\Eone\dev{\tot}_2$ is absorbed into
$\Eone\rk_2$). KL display the Foreign inflation line with the second-form
$\frac{\hbias\frac{1+\zs}{\zs}}{\frac{\cshare}{1-\cshare}+\tel+(1-\tel)(\hbias\frac{1+\zs}{\zs})^2}$
coefficient $=\frac{1+\zs}{\zs}\frac{\hbias}{\Lambda}$ (note
$\hbias(\frac{1+\zs}{\zs})/\Lambda=\frac{1+\zs}{\zs}\cdot\hbias/\Lambda$); we
matched the magnitude and flag the one-line forward-substitution we filled.}
The realized bond returns then follow from~\eqref{eq:realized-bond}:
\begin{align}
\rr_1 &= -\D\dev{P}_1 = -\frac{1}{\taylor}\bigl(\Eone\rk_2-\dev{\cy}_1\bigr), \label{eq:r1-solved}\\
\fr{r}_1 &= -\D\fr{\dev{P}}_1 = -\frac{1}{\taylor}\left(\Eone\rk_2 - \frac{1+\zs}{\zs}\,\frac{\hbias}{\Lambda}\bigl(\fr{\dev{\ell}}_1-\dev{\ell}_1\bigr)\right), \label{eq:rstar1-solved}
\end{align}
with $\Eone\rk_2$ given by~\eqref{eq:Er2k-solved}. Writing
\eqref{eq:Er2k-solved} into~\eqref{eq:r1-solved} gives the fully reduced Home
realized bond return
\begin{equation}
\rr_1 = -\frac{1}{\taylor}\left[-(1-\cshare)\Bigl(\tfrac{1}{1+\zs}\dev{\ell}_1+\tfrac{\zs}{1+\zs}\fr{\dev{\ell}}_1\Bigr)
+\frac{\hbias}{\Lambda}\bigl(\fr{\dev{\ell}}_1-\dev{\ell}_1\bigr)-\dev{\cy}_1\right]. \label{eq:r1-full}
\end{equation}

\paragraph{Step 8: realized capital return.} Finally, substituting the realized
capital-return linearization~\eqref{eq:rk1-raw} and the capital-price
relation~\eqref{eq:qk1-raw} (eliminating $\dev{\qk}_1$, using
$\Eone\dev{\tot}_2$ from~\eqref{eq:Es2} and $\Eone\rk_2$
from~\eqref{eq:Er2k-solved}) gives (their app.~p.~19, ``(93) and (94) imply'')
\begin{multline}
\rk_1 = -\frac{\hbias}{\Lambda}\bigl(\fr{\dev{\ell}}_1-\dev{\ell}_1\bigr)
-\frac{\left(\hbias\,\frac{1+\zs}{\zs}\right)^2}
{\dfrac{\cshare}{1-\cshare}+\tel\left(1-\left(\hbias\,\frac{1+\zs}{\zs}\right)^2\right)}\,\atil\,\Eone\dev{\wsh}_2\\
+(1-\cshare)\Bigl(\tfrac{1}{1+\zs}\dev{\ell}_1+\tfrac{\zs}{1+\zs}\fr{\dev{\ell}}_1\Bigr)
-(1-\cshare)\widetilde{\widetilde{\dev k}}_1. \label{eq:rk1-solved}
\end{multline}
\gapflag{KL display \eqref{eq:rk1-solved} as ``(93) and (94) imply.'' The
reconstruction combines the realized profit linearization~\eqref{eq:rk1-raw}
with the capital-pricing condition~\eqref{eq:qk1-raw}: writing
$\dev{\qk}_1=-\hbias\Eone\dev{\tot}_2-(1-\cshare)\widetilde{\widetilde{\dev k}}_1-\Eone\rk_2$
and substituting into the $\disc\dev{\qk}_1$ term of \eqref{eq:rk1-raw}, then
replacing $\dev{\tot}_1$ via~\eqref{eq:s1-solved}, $\Eone\dev{\tot}_2$
via~\eqref{eq:Es2}, and $\Eone\rk_2$ via~\eqref{eq:Er2k-solved}. The
$\rk_1=\divr$-and-revaluation terms collect: the $(1-\disc)$ profit piece plus
the $\disc$ price-revaluation piece combine so that the labor-income terms
appear with coefficient $(1-\cshare)$ (not $(1-\disc)(1-\cshare)$), the
relative-labor term collapses to $-\hbias/\Lambda$, and the wealth-share term
inherits the $\Eone\dev{\tot}_2$ multiplier $\times\atil$. We verified the
\emph{structure} against KL's displayed line (their app.~p.~19): the four terms
(relative-labor, wealth-share, global-employment, lagged-capital) match in form;
we flag that the exact cancellation of the $\disc$ vs.\ $(1-\disc)$ weights in
combining \eqref{eq:rk1-raw} and \eqref{eq:qk1-raw} is the one place we filled
the algebra.}

\subsection{The reusable reduced forms}
\label{sec:engine-summary}

The engine is now complete. The objects cited by
Propositions~\ref{prop:prices}--\ref{prop:wealth} are the boxed reduced forms
below, each a function of
$(\dev{\ell}_1,\fr{\dev{\ell}}_1,\dev{\cy}_1,\widetilde{\widetilde{\dev k}}_1,
\Eone\dev{\wsh}_2)$ and the deep parameters, with $\Eone\dev{\wsh}_2$ itself
given by~\eqref{eq:Ewsh-solved}:
\[
\boxed{\;
\begin{aligned}
&\dev{\tot}_1 \text{ \eqref{eq:s1-solved}},\quad
\dev{\rer}_1 \text{ \eqref{eq:q1-solved}},\quad
\dev{\rs c}_1 \text{ \eqref{eq:c1-solved}},\quad
\Eone\dev{\wsh}_2 \text{ \eqref{eq:Ewsh-solved}},\quad
\widehat{\nfa}_1 \text{ \eqref{eq:nfa-solved}},\\[2pt]
&\Eone\rk_2 \text{ \eqref{eq:Er2k-solved}},\quad
\D\dev P_1 \text{ \eqref{eq:dP1}},\quad
\D\fr{\dev P}_1 \text{ \eqref{eq:dPstar1}},\quad
\rr_1 \text{ \eqref{eq:r1-full}},\quad
\fr{r}_1 \text{ \eqref{eq:rstar1-solved}},\quad
\rk_1 \text{ \eqref{eq:rk1-solved}}.
\end{aligned}\;}
\]
What remains to close the system is a \emph{labor block}: two equations pinning
$\dev{\ell}_1$ and $\fr{\dev{\ell}}_1$. The simplified environment offers two
closures. Under \textbf{flexible wages with infinite Frisch elasticity}
($\frisch\to0$), the union's wage-setting condition forces (their app.~eq.\
before~(97))
\begin{equation}
\dev{\ell}_1 = \fr{\dev{\ell}}_1 = 0. \label{eq:flex-labor}
\end{equation}
Under \textbf{wages set one period in advance}, the predetermined-wage conditions
are (their app.~eqs.\ before~(97))
\begin{equation}
\dev{\rw}_1 = -\D\dev{P}_1 - \tel^z\dev{\epsilon}^z_1, \qquad
\fr{\dev{\rw}}_1 + \dev{\rer}_1 = -\D\fr{\dev{P}}_1 - \tel^z\dev{\epsilon}^z_1, \label{eq:predet-wage}
\end{equation}
where $\tel^z$ scales the productivity surprise $\dev{\epsilon}^z_1$ (KL's
$\sigma^z$). Combining \eqref{eq:predet-wage} with the firms' log-linearized
labor demand gives the $2\times2$ system (their app.~eqs.~(97)--(98))
\begin{align}
&\left[\Bigl(\tfrac{\zs}{1+\zs}-\hbias\Bigr)+\tfrac{\zs}{1+\zs}\tfrac{\cshare}{1-\cshare}\right]\dev{\tot}_1
-\cshare\Bigl(\tfrac{1}{1+\zs}\dev{\ell}_1+\tfrac{\zs}{1+\zs}\fr{\dev{\ell}}_1\Bigr)
-(1-\cshare)\widetilde{\widetilde{\dev k}}_1 = -\D\dev{P}_1, \label{eq:firm-H}\\
&-\left[\Bigl(\tfrac{1}{1+\zs}-\tfrac{\hbias}{\zs}\Bigr)+\tfrac{1}{1+\zs}\tfrac{\cshare}{1-\cshare}\right]\dev{\tot}_1
-\cshare\Bigl(\tfrac{1}{1+\zs}\dev{\ell}_1+\tfrac{\zs}{1+\zs}\fr{\dev{\ell}}_1\Bigr)
-(1-\cshare)\widetilde{\widetilde{\dev k}}_1 = -\D\fr{\dev{P}}_1. \label{eq:firm-F}
\end{align}
Equations~\eqref{eq:firm-H}--\eqref{eq:firm-F} are the combined firm-labor-demand
system; together with the inflation reduced forms \eqref{eq:dP1}--\eqref{eq:dPstar1}
they determine $(\dev{\ell}_1,\fr{\dev{\ell}}_1)$. The proofs of
Propositions~\ref{prop:prices}--\ref{prop:wealth} take exactly this route: impose
the relevant portfolio specialization (which fixes $\dev{\wsh}_1$ and
$\Eone\dev{\wsh}_2$), then solve the labor block, then read off prices, returns,
and revaluation from the boxed reduced forms.


% ---- 4. Proposition 1: prices and production ----
% =====================================================================
% 04_prop1.tex — Proposition 1: prices and production (WP-B)
% =====================================================================
\section{Proposition 1: effects of a safety shock on prices and production}
\label{sec:prop1}

This section proves Proposition~1 of KL: how a flight-to-safety shock
$\dev{\cy}_1>0$ propagates to expected real rates, inflation, the relative price
of capital, and global employment. The proof specializes the
engine of Section~\ref{sec:engine} to identical portfolios and complete home
bias, which forces the wealth share to be inert, so that the only moving parts
are prices and production. We treat the two wage regimes in turn.

\subsection{Statement}

\begin{proposition}[Effects of a safety shock on prices and production]
\label{prop:prices}
Consider the simplified economy of Definition~\ref{def:simplified} with identical
portfolios across countries, no Home government safe debt
($b^g_{H,s}=0$), and complete home bias $\hbias\to\zs/(1+\zs)$, so that
$\dev{\wsh}_1=0$ and $\Eone\dev{\wsh}_2=0$. Consider a flight-to-safety shock
$\dev{\cy}_1>0$ in period~$1$, with $\widetilde{\widetilde{\dev k}}_1=0$ on
impact.
\begin{enumerate}[label=(\alph*),leftmargin=1.8em,itemsep=3pt]
  \item \textbf{Flexible wages} ($\frisch\to0$). The shock is absorbed entirely
  by lower expected real rates and Home deflation, with no effect on the real
  exchange rate or production:
  \begin{equation}
  \Eone\rr_2 = -\dev{\cy}_1, \qquad
  \D\dev{P}_1 = -\frac{1}{\taylor}\dev{\cy}_1, \qquad
  \dev{\rer}_1 = \dev{\ell}_1 = \fr{\dev{\ell}}_1 = 0. \label{eq:prop1-flex}
  \end{equation}
  \item \textbf{Wages set one period in advance.} The shock causes a global
  recession and the responses are \emph{strictly} attenuated relative to the
  flexible-wage benchmark:
  \begin{equation}
  0 > \Eone\rr_2 > -\dev{\cy}_1, \qquad
  0 > \D\dev{P}_1 > -\frac{1}{\taylor}\dev{\cy}_1, \qquad
  \dev{\rer}_1 \propto \dev{\cy}_1 \;(>0), \label{eq:prop1-sticky}
  \end{equation}
  and global employment falls, \emph{disproportionately at Home}:
  \begin{equation}
  \frac{1}{1+\zs}\dev{\ell}_1 + \frac{\zs}{1+\zs}\fr{\dev{\ell}}_1 \;\propto\; -\dev{\cy}_1, \qquad
  \dev{\ell}_1 - \fr{\dev{\ell}}_1 \;\propto\; -\dev{\cy}_1. \label{eq:prop1-labor}
  \end{equation}
\end{enumerate}
\end{proposition}

\noindent (We transcribe the proposition to match KL's main-paper statement; the
proportionality constants in \eqref{eq:prop1-sticky}--\eqref{eq:prop1-labor} are
made explicit in the proof.)

\subsection{Proof}

\paragraph{Specialization of the engine.} Following KL's app.~B.3.1, impose
identical portfolios --- so that the steady-state positions satisfy
$\qk\,k=\sav$, $b_F=0$, $b_H=0$ --- and no Home government safe debt,
$b^g_{H,s}=0$. The realized wealth-share equation~\eqref{eq:wsh-realized} then
has \emph{every} coefficient equal to zero:
\[
\dev{\wsh}_1 = \underbrace{\Bigl(\tfrac{\qk k}{\sav}-1\Bigr)}_{=0}\bigl(\rk_1-\rr_1\bigr)
+\underbrace{\tfrac{b_F}{\sav}}_{=0}\bigl(\fr{r}_1-\dev{\rer}_1-\rr_1\bigr)
-\disc\,\tfrac{\zs}{1+\zs}\underbrace{\tfrac{b^g_{H,s}}{\sav}}_{=0}\dev{\cy}_1
= 0.
\]
Moreover, with complete home bias $\hbias\to\zs/(1+\zs)$, the net Home
expenditure weight in the wealth-share law~\eqref{eq:Ewsh-solved} vanishes,
$\frac{\zs}{1+\zs}-\hbias\to0$, and the relative-labor channel multiplier
collapses; combined with $\dev{\wsh}_1=0$ this gives
\begin{equation}
\Eone\dev{\wsh}_2 = 0. \label{eq:prop1-wsh0}
\end{equation}
\gapflag{KL state ``since $\hbias\to\zs/(1+\zs)$, we have that $\Eone\dev{\wsh}_2=0$''
(their app.~B.3.1). Two facts combine. First, $\dev{\wsh}_1=0$ from the
identical-portfolio specialization above. Second, in \eqref{eq:Ewsh-solved} the
surviving relative-labor term carries the factor
$\bigl(\frac{\zs}{1+\zs}-\hbias\bigr)$ \emph{through} $\Eone\dev{\tot}_2$ and the
$\atil$ feedback: with complete home bias $\hbias=\frac{\zs}{1+\zs}$ one checks
$\hbias\frac{1+\zs}{\zs}=1$, so $1-(\hbias\frac{1+\zs}{\zs})^2=0$, killing the
relative-labor coefficient in \eqref{eq:Ewsh-solved}. Hence $\Eone\dev{\wsh}_2
=\dev{\wsh}_1=0$. We verified $\hbias\frac{1+\zs}{\zs}=1$ at complete home bias
and that this zeroes the $1-(\hbias\frac{1+\zs}{\zs})^2$ factor; KL suppress this
one-line check.}
With \eqref{eq:prop1-wsh0}, the wealth-share terms drop out of every reduced form
in Section~\ref{sec:engine-summary}. In particular, using
$\hbias\frac{1+\zs}{\zs}=1$ (complete home bias) the composite denominator
\eqref{eq:Lambda-def} simplifies to
\begin{equation}
\Lambda = \frac{\cshare}{1-\cshare}+\tel+(1-\tel)\cdot1
=\frac{\cshare}{1-\cshare}+1 = \frac{1}{1-\cshare}, \label{eq:Lambda-cmplt}
\end{equation}
so the relative-labor coefficient becomes
$\hbias/\Lambda=\frac{\zs}{1+\zs}(1-\cshare)$. We record this for the sticky-wage
case below.

\subparagraph{(a) Flexible wages.} With flexible wages and infinite Frisch
elasticity $\frisch\to0$, the union's wage-setting condition gives
\eqref{eq:flex-labor}:
\[
\dev{\ell}_1 = \fr{\dev{\ell}}_1 = 0.
\]
Substitute $\dev{\ell}_1=\fr{\dev{\ell}}_1=0$, $\Eone\dev{\wsh}_2=0$, and
$\widetilde{\widetilde{\dev k}}_1=0$ into the engine. The expected capital
return~\eqref{eq:Er2k-solved} becomes
\[
\Eone\rk_2 = -(1-\cshare)\underbrace{\Bigl(\tfrac{1}{1+\zs}\dev{\ell}_1+\tfrac{\zs}{1+\zs}\fr{\dev{\ell}}_1\Bigr)}_{=0}
+\frac{\hbias}{\Lambda}\underbrace{\bigl(\fr{\dev{\ell}}_1-\dev{\ell}_1\bigr)}_{=0} = 0.
\]
The capital-Euler relation~\eqref{eq:euler-r} then gives the expected real bond
return,
\begin{equation}
\Eone\rr_2 = \Eone\rk_2 - \dev{\cy}_1 = -\dev{\cy}_1, \label{eq:prop1a-Er2}
\end{equation}
which is the first claim in~\eqref{eq:prop1-flex}. The Home Taylor
rule~\eqref{eq:dP1} gives the inflation deviation,
\begin{equation}
\D\dev{P}_1 = \frac{1}{\taylor}\bigl(\Eone\rk_2-\dev{\cy}_1\bigr)
=\frac{1}{\taylor}\bigl(0-\dev{\cy}_1\bigr) = -\frac{1}{\taylor}\dev{\cy}_1, \label{eq:prop1a-dP}
\end{equation}
the second claim. Finally the capital-price relation~\eqref{eq:qk1-raw} with
$\Eone\dev{\tot}_2=0$ (from \eqref{eq:Es2} and \eqref{eq:prop1-wsh0}),
$\widetilde{\widetilde{\dev k}}_1=0$, and $\Eone\rk_2=0$ gives
\[
\dev{\qk}_1 = -\hbias\underbrace{\Eone\dev{\tot}_2}_{=0}-(1-\cshare)\underbrace{\widetilde{\widetilde{\dev k}}_1}_{=0}-\underbrace{\Eone\rk_2}_{=0}=0,
\]
which with \eqref{eq:flex-labor} completes \eqref{eq:prop1-flex}. KL note these
claims ``follow immediately from the above results given
$\dev{\ell}_1=\fr{\dev{\ell}}_1=0$.''

\subparagraph{(b) Wages set one period in advance.} Now wages are predetermined.
Substituting the inflation reduced forms~\eqref{eq:dP1}--\eqref{eq:dPstar1} (with
$\Eone\dev{\wsh}_2=0$) into the firm-labor-demand
system~\eqref{eq:firm-H}--\eqref{eq:firm-F}, and using complete home bias and
$\widetilde{\widetilde{\dev k}}_1=0$, KL reduce the $2\times2$ system to a linear
map from $(\dev{\cy}_1,\widetilde{\widetilde{\dev k}}_1)$ to
$(\dev{\ell}_1,\fr{\dev{\ell}}_1)$. The solution is (their app.~B.3.1)
\begin{align}
\dev{\ell}_1 &= -\frac{1}{\cshare+\frac{1}{\taylor}(1-\cshare)}\,\frac{1}{\taylor}\,\dev{\cy}_1
\;-\;\frac{1}{\cshare+\frac{1}{\taylor}(1-\cshare)}\,(1-\cshare)\,\widetilde{\widetilde{\dev k}}_1, \label{eq:prop1b-l}\\
\fr{\dev{\ell}}_1 &= -\frac{1}{\cshare+\frac{1}{\taylor}(1-\cshare)}\,(1-\cshare)\,\widetilde{\widetilde{\dev k}}_1. \label{eq:prop1b-lstar}
\end{align}
\gapflag{KL display \eqref{eq:prop1b-l}--\eqref{eq:prop1b-lstar} as the solution
of \eqref{eq:firm-H}--\eqref{eq:firm-F} but do not show the elimination. The
reconstruction: at complete home bias $\hbias\frac{1+\zs}{\zs}=1$, so the
terms-of-trade coefficients in \eqref{eq:firm-H}--\eqref{eq:firm-F} simplify
(the $\frac{\zs}{1+\zs}-\hbias$ and $\frac{1}{1+\zs}-\frac{\hbias}{\zs}$ pieces
each $\to0$, leaving only the $\frac{\cshare}{1-\cshare}$-weighted terms-of-trade
terms). Replacing $\D\dev P_1,\D\fr{\dev P}_1$ on the right via
\eqref{eq:dP1}--\eqref{eq:dPstar1} with $\Eone\rk_2$ from~\eqref{eq:Er2k-solved}
(and $\Lambda=1/(1-\cshare)$, so $\frac{\hbias}{\Lambda}=\frac{\zs}{1+\zs}(1-\cshare)$),
the system becomes, after collecting the global-employment combination
$L\equiv\frac{1}{1+\zs}\dev{\ell}_1+\frac{\zs}{1+\zs}\fr{\dev{\ell}}_1$ and the
relative-employment gap $G\equiv\dev{\ell}_1-\fr{\dev{\ell}}_1$:
\[
\cshare L + \tfrac{1}{\taylor}(1-\cshare)L = -\tfrac{1}{\taylor}\dev{\cy}_1
-(1-\cshare)\widetilde{\widetilde{\dev k}}_1,
\qquad
\bigl[\cshare+\tfrac{1}{\taylor}(1-\cshare)\bigr]\,(\text{relative part}) = -\tfrac{1}{\taylor}\dev{\cy}_1,
\]
where the first equation is the population-weighted (Home+Foreign) sum of
\eqref{eq:firm-H}--\eqref{eq:firm-F} and the second is the Home$-$Foreign
difference. The convenience shock $\dev{\cy}_1$ enters only the Home Taylor
rule~\eqref{eq:dP1} (Foreign's \eqref{eq:dPstar1} has no $\dev{\cy}_1$), so it
loads asymmetrically: it appears in the global sum and the difference but
\emph{not} in $\fr{\dev{\ell}}_1$ alone. Solving the two scalar equations for
$\dev{\ell}_1$ and $\fr{\dev{\ell}}_1$ gives
\eqref{eq:prop1b-l}--\eqref{eq:prop1b-lstar}. We checked: (i) $\fr{\dev{\ell}}_1$
has no $\dev{\cy}_1$ term --- the safety shock leaves Foreign employment
unaffected at complete home bias, consistent with KL's remark that
``$\hbias\to\frac{\zs}{1+\zs}$ implies $\fr{\dev{\ell}}_1$ is unaffected by a
safety shock''; (ii) the common denominator
$\cshare+\frac{1}{\taylor}(1-\cshare)$ matches KL's
$\alpha+\frac{1}{\phi}(1-\alpha)$. The forms match KL exactly with
$\frac{1}{\phi}=\frac{1}{\taylor}$.}

Set $\widetilde{\widetilde{\dev k}}_1=0$ on impact. Then
\eqref{eq:prop1b-l}--\eqref{eq:prop1b-lstar} reduce to
\begin{equation}
\dev{\ell}_1 = -\frac{1}{\taylor}\,\frac{1}{\cshare+\frac{1}{\taylor}(1-\cshare)}\,\dev{\cy}_1 < 0,
\qquad \fr{\dev{\ell}}_1 = 0. \label{eq:prop1b-impact}
\end{equation}
Because $\cshare\in(0,1)$ and $\taylor>1$, the denominator
$\cshare+\frac{1}{\taylor}(1-\cshare)>0$, so $\dev{\ell}_1<0$ for $\dev{\cy}_1>0$:
\emph{Home employment falls}, while Foreign employment is unchanged on impact.
We now read off the remaining claims.

\emph{Global employment falls disproportionately at Home.} From
\eqref{eq:prop1b-impact},
\begin{align}
\frac{1}{1+\zs}\dev{\ell}_1+\frac{\zs}{1+\zs}\fr{\dev{\ell}}_1
&= \frac{1}{1+\zs}\,\dev{\ell}_1
= -\frac{1}{1+\zs}\frac{1}{\taylor}\frac{1}{\cshare+\frac{1}{\taylor}(1-\cshare)}\dev{\cy}_1 \;\propto\; -\dev{\cy}_1, \\
\dev{\ell}_1-\fr{\dev{\ell}}_1 &= \dev{\ell}_1
= -\frac{1}{\taylor}\frac{1}{\cshare+\frac{1}{\taylor}(1-\cshare)}\dev{\cy}_1 \;\propto\; -\dev{\cy}_1,
\end{align}
establishing both proportionalities in~\eqref{eq:prop1-labor}: aggregate
employment falls (negative) and the Home$-$Foreign gap is negative (Home falls
by more), exactly the ``disproportionately at Home'' claim.

\emph{Expected capital return and expected real rate.} Substituting
\eqref{eq:prop1b-impact} into~\eqref{eq:Er2k-solved} with
$\Lambda=1/(1-\cshare)$:
\begin{align}
\Eone\rk_2 &= -(1-\cshare)\,\frac{1}{1+\zs}\dev{\ell}_1
+\frac{\zs}{1+\zs}(1-\cshare)\bigl(\fr{\dev{\ell}}_1-\dev{\ell}_1\bigr) \notag\\
&= -(1-\cshare)\frac{1}{1+\zs}\dev{\ell}_1 - \frac{\zs}{1+\zs}(1-\cshare)\dev{\ell}_1
= -(1-\cshare)\dev{\ell}_1 > 0, \label{eq:prop1b-Erk}
\end{align}
using $\fr{\dev{\ell}}_1=0$ and
$\frac{1}{1+\zs}+\frac{\zs}{1+\zs}=1$. Since $\dev{\ell}_1<0$, we have
$\Eone\rk_2>0$. Then the expected real bond return from~\eqref{eq:euler-r} is
\begin{equation}
\Eone\rr_2 = \Eone\rk_2 - \dev{\cy}_1 = -(1-\cshare)\dev{\ell}_1 - \dev{\cy}_1. \label{eq:prop1b-Er2}
\end{equation}
Substituting $\dev{\ell}_1$ from~\eqref{eq:prop1b-impact},
\[
\Eone\rr_2 = (1-\cshare)\frac{1}{\taylor}\frac{1}{\cshare+\frac{1}{\taylor}(1-\cshare)}\dev{\cy}_1 - \dev{\cy}_1
= -\dev{\cy}_1\left[1 - \frac{\frac{1}{\taylor}(1-\cshare)}{\cshare+\frac{1}{\taylor}(1-\cshare)}\right]
= -\dev{\cy}_1\,\frac{\cshare}{\cshare+\frac{1}{\taylor}(1-\cshare)}.
\]
The bracket $\frac{\cshare}{\cshare+\frac{1}{\taylor}(1-\cshare)}\in(0,1)$ for
$\cshare\in(0,1),\taylor>1$, so
\begin{equation}
0 > \Eone\rr_2 = -\frac{\cshare}{\cshare+\frac{1}{\taylor}(1-\cshare)}\dev{\cy}_1 > -\dev{\cy}_1, \label{eq:prop1b-Er2-ineq}
\end{equation}
which is the first strict-inequality chain in~\eqref{eq:prop1-sticky}: the
expected real rate falls, but by \emph{less} than the convenience yield ---
strictly between $0$ and $-\dev{\cy}_1$ --- because part of the safety shock is
now absorbed by the recession (the positive $\Eone\rk_2$) rather than by the real
rate alone.

\emph{Inflation.} From the Home Taylor rule~\eqref{eq:dP1},
\begin{equation}
\D\dev{P}_1 = \frac{1}{\taylor}\bigl(\Eone\rk_2-\dev{\cy}_1\bigr) = \frac{1}{\taylor}\Eone\rr_2
= -\frac{1}{\taylor}\frac{\cshare}{\cshare+\frac{1}{\taylor}(1-\cshare)}\dev{\cy}_1, \label{eq:prop1b-dP}
\end{equation}
so dividing the inequality~\eqref{eq:prop1b-Er2-ineq} by $\taylor>0$,
\[
0 > \D\dev{P}_1 > -\frac{1}{\taylor}\dev{\cy}_1,
\]
the second chain in~\eqref{eq:prop1-sticky}: Home still deflates, but less than in
the flexible benchmark~\eqref{eq:prop1a-dP}.

\emph{Real exchange rate.} At complete home bias $\hbias(1+\zs)/\zs=1$, so the
real-exchange-rate relation~\eqref{eq:q1-s1} reads $\dev{\rer}_1=\dev{\tot}_1$.
From the solved terms of trade~\eqref{eq:s1-solved} with $\Eone\dev{\wsh}_2=0$,
\begin{equation}
\dev{\rer}_1 = \dev{\tot}_1 = \frac{1}{\Lambda}\bigl(\fr{\dev{\ell}}_1-\dev{\ell}_1\bigr)
= -(1-\cshare)\dev{\ell}_1
= (1-\cshare)\frac{1}{\taylor}\frac{1}{\cshare+\frac{1}{\taylor}(1-\cshare)}\dev{\cy}_1 \;>\;0, \label{eq:prop1b-q}
\end{equation}
using $\Lambda=1/(1-\cshare)$ at complete home bias and $\fr{\dev{\ell}}_1=0$,
$\dev{\ell}_1<0$. Since the convention is that a rise in $\rer$ is a \emph{Home
appreciation}, a positive safety shock \emph{appreciates} the Home real exchange
rate by a positive multiple of $\dev{\cy}_1$, establishing the third claim
in~\eqref{eq:prop1-sticky} (in the flexible case $\dev{\ell}_1=0$ so
$\dev{\rer}_1=0$, the third claim in~\eqref{eq:prop1-flex}). \qed
(\,Proposition~\ref{prop:prices}\,)

\begin{remark}[The relative price of capital, an engine by-product]
KL's Proposition~1 makes no claim about the real price of installed capital
$\qk$, but the engine pins it down and it is useful downstream (Propositions~2--3).
From the capital-price relation~\eqref{eq:qk1-raw} with $\Eone\dev{\tot}_2=0$ and
$\widetilde{\widetilde{\dev k}}_1=0$,
\begin{equation}
\dev{\qk}_1 = -\Eone\rk_2 = (1-\cshare)\dev{\ell}_1
= -(1-\cshare)\frac{1}{\taylor}\frac{1}{\cshare+\frac{1}{\taylor}(1-\cshare)}\dev{\cy}_1 \;<\;0. \label{eq:prop1b-qk}
\end{equation}
The price of installed capital \emph{falls} in a flight to safety: the expected
capital return $\Eone\rk_2$ rises (the recession raises the marginal product of
capital), and the no-arbitrage condition $\dev{\qk}_1=-\Eone\rk_2$ revalues
capital down. Note $\dev{\qk}_1=-\dev{\rer}_1$ here: the price of capital and the
real exchange rate move by equal and opposite amounts on impact.
\end{remark}

\subsection{Discussion}

\paragraph{The Euler-equation wedge.} The entire result is driven by the
convenience term in the Home Euler equation. The safe bond enters the
intertemporal condition through the factor $(1+r_{t+1})/(1-\cy_t)$: holders accept
a lower pecuniary return because the bond yields nonpecuniary (safety/liquidity)
services. To first order this is the wedge~\eqref{eq:euler-r},
$\Eone\rk_2-\Eone\rr_2=\dev{\cy}_1$: a flight to safety raises the required
\emph{spread} of capital over safe bonds by exactly $\dev{\cy}_1$.

\paragraph{Flexible wages: one-for-one absorption.} With flexible wages and
prices, output cannot move ($\dev{\ell}_1=\fr{\dev{\ell}}_1=0$, so the marginal
product of capital and hence $\Eone\rk_2$ are unchanged). The required spread is
then met \emph{entirely} by a fall in the expected real rate,
$\Eone\rr_2=-\dev{\cy}_1$. Under the Taylor rule, a lower expected real rate
requires Home deflation, $\D\dev{P}_1=-\frac{1}{\taylor}\dev{\cy}_1$: the central
bank delivers exactly enough deflation that the Fisher relation produces the
needed real-rate decline. There is no recession and no revaluation of capital.

\paragraph{Sticky wages: a deflationary recession, worse at Home.} When nominal
wages are predetermined, the deflationary pressure from the flight to safety
raises the \emph{product wage} $W/P$ (since $P$ falls but $W$ is fixed), which
depresses labor demand and output. The recession raises the expected marginal
product of capital, so $\Eone\rk_2>0$ absorbs part of the required spread; the
expected real rate therefore falls by \emph{less} than $\dev{\cy}_1$
(eq.~\eqref{eq:prop1b-Er2-ineq}), and deflation is correspondingly milder
(eq.~\eqref{eq:prop1b-dP}). The recession is concentrated at Home because the
safety shock enters only the \emph{Home} Taylor rule: the convenience yield is on
\emph{dollar} bonds, so it is Home monetary policy that must accommodate it.
Foreign employment is unchanged on impact ($\fr{\dev{\ell}}_1=0$), so Home's
terms of trade and (via~\eqref{eq:q1-s1}) the real exchange rate \emph{appreciate}
($\dev{\ell}_1<\fr{\dev{\ell}}_1=0\Rightarrow\dev{\tot}_1=\frac{1}{\Lambda}(\fr{\dev{\ell}}_1-\dev{\ell}_1)>0$):
the dollar appreciates in the flight to safety.

\paragraph{Load-bearing assumptions.} (i) \emph{Nominal rigidity} is essential
for the recession: with flexible wages the real allocation is untouched and only
prices move. (ii) The \emph{Taylor rule that does not cut nominal rates
one-for-one with $\cy$} (an inflation coefficient $\taylor$, not a peg on the
convenience-adjusted rate) is what forces deflation to do the work; a central
bank that fully accommodated the safety shock would short-circuit the result.
(iii) The \emph{trade elasticity} $\tel$ governs the size of the terms-of-trade /
real-exchange-rate response through $\Lambda$; here, at complete home bias,
$\Lambda=1/(1-\cshare)$ and the $\tel$-dependence drops from the labor solution,
but it returns away from complete home bias. (iv) \emph{Complete home bias}
$\hbias\to\zs/(1+\zs)$ is what delivers the clean $\dev{\wsh}_1=\Eone\dev{\wsh}_2=0$,
isolating the price/production channel from wealth-share redistribution (which is
the subject of Propositions~\ref{prop:returns}--\ref{prop:wealth}).

\paragraph{Connection to the target facts.} Proposition~\ref{prop:prices} is the
model's account of the empirical regularity that a global flight to safety
coincides with (a) \emph{dollar appreciation} --- here Home's terms of trade and
real exchange rate appreciate because the dollar convenience shock forces Home
deflation and a Home-concentrated contraction --- and (b) a \emph{global
recession that hits the safe-asset issuer's economy hardest}. In the big\_cy
reading, the same mechanism with a \emph{large absolute} convenience yield
$\dev{\cy}_1$ amplifies both the appreciation and the asymmetric contraction,
which is the channel we will quantify downstream.


% ---- 5. Proposition 2: realized and expected excess returns ----
% =====================================================================
% 05_prop2.tex — Proposition 2: realized and expected excess returns (WP-C)
% =====================================================================
\section{Proposition 2: realized and expected excess returns}
\label{sec:prop2}

This section proves Proposition~2 of KL: how a flight-to-safety shock
$\dev{\cy}_1>0$ moves the \emph{realized} and \emph{expected} returns on the
three assets in the simplified economy --- safe dollar bonds, capital, and
Foreign bonds. The entire derivation reuses the impulse-response engine of
Section~\ref{sec:engine}; we do \emph{not} re-derive any reduced form. As in
Proposition~\ref{prop:prices}, we specialize to identical portfolios, no Home
government safe debt ($b^g_{H,s}=0$), and complete home bias
$\hbias\to\zs/(1+\zs)$, so that $\dev{\wsh}_1=0$ and $\Eone\dev{\wsh}_2=0$
(established at \eqref{eq:prop1-wsh0}). The labor responses
$(\dev{\ell}_1,\fr{\dev{\ell}}_1)$ are exactly those of
Proposition~\ref{prop:prices}: zero under flexible wages, and
\eqref{eq:prop1b-impact} under predetermined wages. We treat the two wage
regimes in turn for each asset.

\subsection{Statement}

\begin{proposition}[Realized and expected excess returns]
\label{prop:returns}
Consider the simplified economy of Definition~\ref{def:simplified} with identical
portfolios across countries and no Home government safe debt ($b^g_{H,s}=0$), so
that $\dev{\wsh}_1=0$ and $\Eone\dev{\wsh}_2=0$, and consider a flight-to-safety
shock $\dev{\cy}_1>0$ in period~$1$ with $\widetilde{\widetilde{\dev k}}_1=0$ on
impact. Then, to first order, on impact of the shock:
\begin{enumerate}[label=(\alph*),leftmargin=1.8em,itemsep=3pt]
  \item the realized real return on safe dollar bonds \emph{rises},
  $\rr_1\propto\dev{\cy}_1$ ($>0$), in both wage regimes;
  \item the realized real return on capital is \emph{unaffected} if wages are
  flexible, $\rk_1=0$, but \emph{falls} if wages are set one period in advance,
  $\rk_1\propto-\dev{\cy}_1$ ($<0$);
  \item the realized real return on Foreign bonds (in Home consumption units) is
  \emph{unaffected} if wages are flexible, $\fr{r}_1-\D\dev{\rer}_1=0$, but
  \emph{falls} if wages are set one period in advance,
  $\fr{r}_1-\D\dev{\rer}_1\propto-\dev{\cy}_1$ ($<0$);
  \item the \emph{expected} excess returns on capital and on Foreign bonds over
  dollar bonds are positive and equal to the convenience yield:
  \begin{equation}
  \Eone\bigl[\rk_2-\rr_2\bigr]
  =\Eone\bigl[\fr{r}_2-\D\dev{\rer}_2-\rr_2\bigr]
  =\dev{\cy}_1. \label{eq:prop2-wedge}
  \end{equation}
\end{enumerate}
\end{proposition}

\noindent (We transcribe the proposition to match KL's main-paper statement
(p.~1663). The proportionality constants in (a)--(c) are made explicit in the
proof; here $\D\dev{\rer}_t\equiv\dev{\rer}_t-\dev{\rer}_{t-1}$ is the change in
the real exchange rate, and since the economy is at its deterministic steady
state in period~$0$, $\dev{\rer}_0=0$ so $\D\dev{\rer}_1=\dev{\rer}_1$.)

\subsection{Proof}

\paragraph{Specialization of the engine.} As in the proof of
Proposition~\ref{prop:prices} (their app.~B.3.1), identical portfolios and
$b^g_{H,s}=0$ make the realized wealth-share equation~\eqref{eq:wsh-realized}
vanish, and complete home bias $\hbias\to\zs/(1+\zs)$ gives
$\Eone\dev{\wsh}_2=0$ \eqref{eq:prop1-wsh0}, so the composite denominator
collapses to $\Lambda=1/(1-\cshare)$ \eqref{eq:Lambda-cmplt} and the
relative-labor coefficient is $\hbias/\Lambda=\frac{\zs}{1+\zs}(1-\cshare)$. We
set $\widetilde{\widetilde{\dev k}}_1=0$ on impact throughout. Recall the labor
solutions: flexible wages give $\dev{\ell}_1=\fr{\dev{\ell}}_1=0$
\eqref{eq:flex-labor}, while predetermined wages give \eqref{eq:prop1b-impact}
\[
\dev{\ell}_1 = -\frac{1}{\taylor}\,\frac{1}{\cshare+\frac{1}{\taylor}(1-\cshare)}\,\dev{\cy}_1 < 0,
\qquad \fr{\dev{\ell}}_1 = 0.
\]
For compactness write
\begin{equation}
D \;\equiv\; \cshare+\tfrac{1}{\taylor}(1-\cshare) \;>\;0,
\qquad
\dev{\ell}_1 = -\frac{1}{\taylor D}\,\dev{\cy}_1, \label{eq:prop2-Ddef}
\end{equation}
so that under predetermined wages $\dev{\ell}_1<0$ for $\dev{\cy}_1>0$. We also
reuse, from the proof of Proposition~\ref{prop:prices}, the predetermined-wage
values (eqs.~\eqref{eq:prop1b-Erk}, \eqref{eq:prop1b-q})
\begin{equation}
\Eone\rk_2 = -(1-\cshare)\dev{\ell}_1 \;>\;0,
\qquad
\dev{\rer}_1 = \dev{\tot}_1 = -(1-\cshare)\dev{\ell}_1 \;>\;0, \label{eq:prop2-reuse}
\end{equation}
where the real-exchange-rate value uses $\hbias\frac{1+\zs}{\zs}=1$ at complete
home bias (so $\dev{\rer}_1=\dev{\tot}_1$) and $\Lambda=1/(1-\cshare)$.

\subparagraph{(a) Realized return on safe dollar bonds.} The realized dollar-bond
return is minus realized Home inflation, $\rr_1=-\D\dev P_1$
\eqref{eq:realized-bond}, with $\D\dev P_1$ from the Home Taylor
rule~\eqref{eq:dP1}, so $\rr_1=\frac{1}{\taylor}(\dev{\cy}_1-\Eone\rk_2)$ by
\eqref{eq:r1-solved}.

\emph{Flexible wages.} Here $\dev{\ell}_1=\fr{\dev{\ell}}_1=0$, so by
\eqref{eq:Er2k-solved} $\Eone\rk_2=0$, and
\begin{equation}
\rr_1 = -\D\dev P_1 = \frac{1}{\taylor}\dev{\cy}_1 \;>\;0. \label{eq:prop2a-flex}
\end{equation}
This is KL's ``$\rr_1=\frac{1}{\phi}\dev{\cy}_1>0$.''

\emph{Predetermined wages.} Now $\Eone\rk_2=-(1-\cshare)\dev{\ell}_1>0$ by
\eqref{eq:prop2-reuse}, so by \eqref{eq:r1-solved}
\begin{equation}
\rr_1 = \frac{1}{\taylor}\bigl(\dev{\cy}_1-\Eone\rk_2\bigr)
= \frac{1}{\taylor}\bigl(\dev{\cy}_1+(1-\cshare)\dev{\ell}_1\bigr)
= \frac{1}{\taylor}\dev{\cy}_1\left[1-\frac{1-\cshare}{\taylor D}\right]. \label{eq:prop2a-sticky}
\end{equation}
The bracket is
$1-\frac{1-\cshare}{\taylor D}=\frac{\taylor D-(1-\cshare)}{\taylor D}
=\frac{\taylor\cshare+(1-\cshare)-(1-\cshare)}{\taylor D}
=\frac{\taylor\cshare}{\taylor D}=\frac{\cshare}{D}\in(0,1)$, using
$\taylor D=\taylor\cshare+(1-\cshare)$ from \eqref{eq:prop2-Ddef}. Hence
\begin{equation}
\rr_1 = \frac{1}{\taylor}\frac{\cshare}{D}\dev{\cy}_1
=\frac{\cshare}{\taylor\cshare+(1-\cshare)}\dev{\cy}_1 \;>\;0,
\qquad\text{i.e.}\quad \rr_1\propto\dev{\cy}_1. \label{eq:prop2a-sticky2}
\end{equation}
In both regimes $\rr_1\propto\dev{\cy}_1>0$: the realized real return on dollar
bonds rises, establishing claim~(a). (Consistently, $\rr_1=-\D\dev P_1$ and
$\D\dev P_1<0$ in both regimes by Proposition~\ref{prop:prices}: Home deflates,
so the real return on the predetermined nominal bond rises.)

\subparagraph{(b) Realized return on capital.} We read $\rk_1$ directly from the
solved engine reduced form~\eqref{eq:rk1-solved}. With $\Eone\dev{\wsh}_2=0$ and
$\widetilde{\widetilde{\dev k}}_1=0$ the wealth-share and capital terms drop:
\begin{equation}
\rk_1 = -\frac{\hbias}{\Lambda}\bigl(\fr{\dev{\ell}}_1-\dev{\ell}_1\bigr)
+(1-\cshare)\Bigl(\tfrac{1}{1+\zs}\dev{\ell}_1+\tfrac{\zs}{1+\zs}\fr{\dev{\ell}}_1\Bigr). \label{eq:prop2b-rk}
\end{equation}

\emph{Flexible wages.} With $\dev{\ell}_1=\fr{\dev{\ell}}_1=0$, both terms vanish:
\begin{equation}
\rk_1 = 0. \label{eq:prop2b-flex}
\end{equation}

\emph{Predetermined wages.} Using $\fr{\dev{\ell}}_1=0$,
$\hbias/\Lambda=\frac{\zs}{1+\zs}(1-\cshare)$, and
$\frac{1}{1+\zs}\dev{\ell}_1+\frac{\zs}{1+\zs}\fr{\dev{\ell}}_1
=\frac{1}{1+\zs}\dev{\ell}_1$,
\begin{align}
\rk_1 &= -\frac{\zs}{1+\zs}(1-\cshare)\bigl(0-\dev{\ell}_1\bigr)
+(1-\cshare)\frac{1}{1+\zs}\dev{\ell}_1 \notag\\
&= (1-\cshare)\dev{\ell}_1\left[\frac{\zs}{1+\zs}+\frac{1}{1+\zs}\right]
= (1-\cshare)\dev{\ell}_1 \;<\;0, \label{eq:prop2b-sticky}
\end{align}
since $\frac{\zs}{1+\zs}+\frac{1}{1+\zs}=1$ and $\dev{\ell}_1<0$. Substituting
$\dev{\ell}_1=-\frac{1}{\taylor D}\dev{\cy}_1$,
\begin{equation}
\rk_1 = -(1-\cshare)\frac{1}{\taylor D}\dev{\cy}_1 \;\propto\; -\dev{\cy}_1. \label{eq:prop2b-sticky2}
\end{equation}
This establishes claim~(b): $\rk_1=0$ under flexible wages and
$\rk_1\propto-\dev{\cy}_1<0$ under predetermined wages. (Note
$\rk_1=(1-\cshare)\dev{\ell}_1=-\Eone\rk_2$ by \eqref{eq:prop2-reuse}; the
realized capital return moves \emph{opposite} to the expected capital return on
impact, a fact we use in the discussion.)

\subparagraph{(c) Realized excess return on Foreign bonds.} The object of
interest is the realized return on Foreign bonds expressed in Home consumption
units \emph{relative to} the real-exchange-rate change, $\fr{r}_1-\D\dev{\rer}_1$.
Since $\dev{\rer}_0=0$ (steady state), $\D\dev{\rer}_1=\dev{\rer}_1$. The Foreign
realized bond return is $\fr{r}_1=-\D\fr{\dev P}_1$ \eqref{eq:realized-bond}, with
$\D\fr{\dev P}_1$ from \eqref{eq:dPstar1}:
\begin{equation}
\fr{r}_1 = -\D\fr{\dev P}_1
= -\frac{1}{\taylor}\left(\Eone\rk_2-\frac{1+\zs}{\zs}\frac{\hbias}{\Lambda}\bigl(\fr{\dev{\ell}}_1-\dev{\ell}_1\bigr)\right). \label{eq:prop2c-rstar}
\end{equation}

\emph{Flexible wages.} With $\dev{\ell}_1=\fr{\dev{\ell}}_1=0$ we have
$\Eone\rk_2=0$, so $\fr{r}_1=0$; and $\dev{\rer}_1=\dev{\tot}_1=0$ by
\eqref{eq:s1-solved} (both labor terms and $\Eone\dev{\wsh}_2$ vanish), so
$\D\dev{\rer}_1=0$. Hence
\begin{equation}
\fr{r}_1-\D\dev{\rer}_1 = 0-0 = 0. \label{eq:prop2c-flex}
\end{equation}

\emph{Predetermined wages.} Use $\fr{\dev{\ell}}_1=0$, complete home bias
$\frac{1+\zs}{\zs}\frac{\hbias}{\Lambda}=\frac{\hbias\frac{1+\zs}{\zs}}{\Lambda}
=\frac{1}{\Lambda}=(1-\cshare)$ (since $\hbias\frac{1+\zs}{\zs}=1$ and
$\Lambda=1/(1-\cshare)$), and $\Eone\rk_2=-(1-\cshare)\dev{\ell}_1$. Then
\eqref{eq:prop2c-rstar} becomes
\begin{align}
\fr{r}_1 &= -\frac{1}{\taylor}\Bigl(-(1-\cshare)\dev{\ell}_1-(1-\cshare)\bigl(0-\dev{\ell}_1\bigr)\Bigr)
= -\frac{1}{\taylor}\Bigl(-(1-\cshare)\dev{\ell}_1+(1-\cshare)\dev{\ell}_1\Bigr) \notag\\
&= 0. \label{eq:prop2c-rstar-val}
\end{align}
\gapflag{KL state claim~(c) in the main paper but suppress the algebra. The
cancellation in \eqref{eq:prop2c-rstar-val} is exact: with $\fr{\dev{\ell}}_1=0$
the two pieces inside the parenthesis are $\Eone\rk_2=-(1-\cshare)\dev{\ell}_1$
and $\frac{1+\zs}{\zs}\frac{\hbias}{\Lambda}(\fr{\dev{\ell}}_1-\dev{\ell}_1)
=(1-\cshare)(-\dev{\ell}_1)=-(1-\cshare)\dev{\ell}_1$, which are equal, so
$\fr{r}_1=-\D\fr{\dev P}_1=0$. Economically: at complete home bias and
$\fr{\dev{\ell}}_1=0$, Foreign output and Foreign inflation are unchanged on
impact (the safety shock hits only the Home Taylor rule), so the Foreign nominal
bond's realized real return is zero.} Meanwhile, by \eqref{eq:prop2-reuse},
$\D\dev{\rer}_1=\dev{\rer}_1=-(1-\cshare)\dev{\ell}_1>0$. Therefore
\begin{equation}
\fr{r}_1-\D\dev{\rer}_1 = 0-\bigl(-(1-\cshare)\dev{\ell}_1\bigr)
= (1-\cshare)\dev{\ell}_1 \;<\;0,
\qquad \fr{r}_1-\D\dev{\rer}_1\propto-\dev{\cy}_1. \label{eq:prop2c-sticky}
\end{equation}
This establishes claim~(c): the Home-currency realized return on Foreign bonds is
zero under flexible wages and falls under predetermined wages. The driver is the
real \emph{appreciation} of the Home currency ($\D\dev{\rer}_1>0$): a Foreign
bond loses value in Home consumption units when the dollar appreciates. (Note
that $\fr{r}_1-\D\dev{\rer}_1=(1-\cshare)\dev{\ell}_1=\rk_1$ by
\eqref{eq:prop2b-sticky}: the realized excess returns on capital and on Foreign
bonds coincide on impact under these specializations.)

\subparagraph{(d) Expected excess returns.} The expected wedges are read directly
from the linearized Euler block of the engine, \eqref{eq:euler-r} and
\eqref{eq:euler-rstar} (their app.~eqs.~(84)--(85)). The capital wedge
\eqref{eq:euler-r} is, in difference form,
\begin{equation}
\Eone\rk_2 = \Eone\rr_2 + \dev{\cy}_1
\;\Longrightarrow\;
\Eone\bigl[\rk_2-\rr_2\bigr] = \dev{\cy}_1, \label{eq:prop2d-cap}
\end{equation}
which holds \emph{regardless} of the wage regime (it is the household's
intertemporal condition, not a function of the labor solution). This is the first
equality in \eqref{eq:prop2-wedge}.

For the Foreign-bond wedge, the engine's expected Foreign-return relation
\eqref{eq:euler-rstar} reads
\begin{equation}
\Eone\fr{r}_2 = \Eone\rr_2 + \dev{\cy}_1 + \hbias\,\frac{1+\zs}{\zs}\,\D\Eone\dev{\tot}_2. \label{eq:prop2d-rstar-raw}
\end{equation}
Here $\D\Eone\dev{\tot}_2\equiv\Eone\dev{\tot}_2-\dev{\tot}_1$ is the expected
period-$1$-to-$2$ change in the terms of trade. We must show this collapses to
the clean statement $\Eone[\fr{r}_2-\D\dev{\rer}_2-\rr_2]=\dev{\cy}_1$, i.e.\ that
the terms-of-trade adjustment exactly accounts for the expected
real-exchange-rate change $\Eone\D\dev{\rer}_2$. Two facts close this.

By the period-$2$ real-exchange relation~\eqref{eq:Eq2},
$\Eone\dev{\rer}_2=\hbias\frac{1+\zs}{\zs}\Eone\dev{\tot}_2$, and by the
period-$1$ relation~\eqref{eq:q1-s1},
$\dev{\rer}_1=\hbias\frac{1+\zs}{\zs}\dev{\tot}_1$. Differencing the two,
\begin{equation}
\Eone\D\dev{\rer}_2 \equiv \Eone\dev{\rer}_2-\dev{\rer}_1
= \hbias\,\frac{1+\zs}{\zs}\bigl(\Eone\dev{\tot}_2-\dev{\tot}_1\bigr)
= \hbias\,\frac{1+\zs}{\zs}\,\D\Eone\dev{\tot}_2. \label{eq:prop2d-Dq2}
\end{equation}
Substituting \eqref{eq:prop2d-Dq2} into \eqref{eq:prop2d-rstar-raw} gives,
\emph{identically in all parameters and both wage regimes},
\begin{equation}
\Eone\fr{r}_2 = \Eone\rr_2 + \dev{\cy}_1 + \Eone\D\dev{\rer}_2
\;\Longrightarrow\;
\Eone\bigl[\fr{r}_2-\D\dev{\rer}_2-\rr_2\bigr] = \dev{\cy}_1, \label{eq:prop2d-rstar}
\end{equation}
the second equality in \eqref{eq:prop2-wedge}.
\gapflag{KL write \eqref{eq:euler-rstar}/(85) as
$\Eone\fr r_2=\Eone\rr_2+\dev{\cy}_1+\hbias\frac{1+\zs}{\zs}\D\Eone\dev{\tot}_2$
``where we have used~(76)'' (their \eqref{eq:Eq2}), and state the wedge result
$\Eone[\fr r_2-\D\dev{\rer}_2-\rr_2]=\dev{\cy}_1$ in Proposition~2 without
re-displaying the one-line identification of
$\hbias\frac{1+\zs}{\zs}\D\Eone\dev{\tot}_2$ with $\Eone\D\dev{\rer}_2$. We supply
that identification: it is just the period-$2$-minus-period-$1$ difference of the
two real-exchange-rate relations \eqref{eq:q1-s1} and \eqref{eq:Eq2}, both of
which carry the \emph{same} proportionality constant $\hbias\frac{1+\zs}{\zs}$,
so the difference is
$\hbias\frac{1+\zs}{\zs}\D\Eone\dev{\tot}_2=\Eone\D\dev{\rer}_2$ exactly. The
result therefore holds for general $\hbias$, not only at complete home bias; the
home-bias specialization is \emph{not} needed for part~(d). Under the
Proposition's specialization $\Eone\dev{\wsh}_2=0$ one further has
$\Eone\dev{\tot}_2=0$ \eqref{eq:Es2}, so $\D\Eone\dev{\tot}_2=-\dev{\tot}_1$ and
$\Eone\D\dev{\rer}_2=-\dev{\rer}_1$; but the wedge identity
\eqref{eq:prop2d-rstar} does not rely on this, only on the common constant of
proportionality.}

\medskip
Combining \eqref{eq:prop2d-cap} and \eqref{eq:prop2d-rstar} yields the displayed
equation~\eqref{eq:prop2-wedge}:
\[
\Eone\bigl[\rk_2-\rr_2\bigr]
=\Eone\bigl[\fr{r}_2-\D\dev{\rer}_2-\rr_2\bigr]=\dev{\cy}_1>0,
\]
the expected excess returns on capital and on Foreign bonds over dollar bonds are
both positive and equal to the convenience yield, establishing claim~(d). \qed
(\,Proposition~\ref{prop:returns}\,)

\begin{remark}[Realized vs.\ expected: a sign flip]
Collecting the predetermined-wage results, on impact
\[
\rr_1 = \frac{\cshare}{\taylor\cshare+(1-\cshare)}\dev{\cy}_1 > 0,
\qquad
\rk_1 = \fr{r}_1-\D\dev{\rer}_1 = (1-\cshare)\dev{\ell}_1 < 0,
\]
so the \emph{realized} excess returns on capital and Foreign bonds are
\[
\rk_1-\rr_1 < 0,
\qquad
\fr{r}_1-\D\dev{\rer}_1-\rr_1 < 0,
\]
both strictly negative, while the \emph{expected} excess returns going forward
are $\Eone[\rk_2-\rr_2]=\Eone[\fr r_2-\D\dev{\rer}_2-\rr_2]=\dev{\cy}_1>0$. The
risky assets underperform safe dollar bonds exactly on impact of the flight to
safety, and are expected to outperform thereafter. This realized-vs-expected sign
flip is the engine of the dollar's negative risk premium (Proposition~5).
\end{remark}

\subsection{Discussion}

\paragraph{Realized excess returns are negative on impact.} Proposition~2 says
that on impact of a positive safety shock the realized excess returns on capital
and on Foreign bonds are \emph{negative}: both assets pay badly exactly when the
flight to safety hits. Two distinct forces produce this, with different scope.

\emph{Force 1 (deflation; operative even with flexible wages).} The flight to
safety forces Home deflation, $\D\dev P_1<0$ (Proposition~\ref{prop:prices}),
which raises the realized real return on the predetermined nominal dollar bond,
$\rr_1=-\D\dev P_1>0$. This channel is present \emph{even under flexible wages}:
in that case $\rk_1=0$ and $\fr{r}_1-\D\dev{\rer}_1=0$, yet $\rr_1>0$, so the
realized excess returns $\rk_1-\rr_1$ and $\fr{r}_1-\D\dev{\rer}_1-\rr_1$ are
strictly negative purely because the dollar bond's real return rises. The safe
bond is a hedge against safety shocks through deflation alone.

\emph{Force 2 (recession and real appreciation; requires nominal rigidity).} With
wages set in advance, the safety shock additionally causes a Home-concentrated
recession ($\dev{\ell}_1<0$) and a real dollar appreciation ($\dev{\rer}_1>0$).
The recession raises the expected marginal product of capital, so the price of
installed capital falls today ($\dev{\qk}_1<0$), pulling down the realized
capital return: $\rk_1=(1-\cshare)\dev{\ell}_1<0$. And the real appreciation
directly lowers the Home-currency realized return on Foreign bonds:
$\fr{r}_1-\D\dev{\rer}_1=(1-\cshare)\dev{\ell}_1<0$. So under nominal rigidity
both risky assets fall in \emph{absolute} terms, not merely relative to the
appreciating dollar bond.

\paragraph{Expected excess returns are positive going forward.} For agents to
hold capital and Foreign bonds despite their poor impact returns, these assets
must be \emph{expected} to outperform dollar bonds going forward. This is exactly
the Euler-equation wedge \eqref{eq:prop2-wedge}: the convenience yield
$\dev{\cy}_1$ is the compensation, in expected excess return, that safe dollar
bonds forgo because of the safety/liquidity services they provide. A larger
safety shock raises both the impact underperformance of risky assets and the
forward-looking expected excess return that restores indifference.

\paragraph{Connection to the target facts: the dollar as a hedge.} The
realized-vs-expected sign flip is the model's microfoundation for the empirical
fact that \emph{safe dollar bonds are a hedge}: they deliver high realized real
returns precisely in states (flights to safety) when capital and foreign-currency
bonds deliver low returns. An asset that pays well in bad times commands a
\emph{negative} risk premium and has a \emph{negative beta} with respect to the
global risk factor --- the dollar's ``exorbitant privilege'' as a negative-beta
asset. Proposition~2 supplies the realized covariance structure (dollar bonds up,
risky assets down, in the same shock) that Proposition~\ref{prop:wealth} and
Proposition~5 will turn into wealth redistribution and a signed risk premium. In
the big\_cy reading, a \emph{large absolute} convenience yield $\dev{\cy}_1$
magnifies the wedge in \eqref{eq:prop2-wedge} and hence the magnitude of the
dollar's negative risk premium, without requiring any asymmetry in risk tolerance
across countries --- the convenience yield does the work that KL otherwise load
onto greater US risk-bearing capacity.


% ---- 6. Proposition 3: wealth share / NFA revaluation ----
% =====================================================================
% 06_prop3.tex — Proposition 3: Home wealth share / NFA revaluation (WP-C)
% =====================================================================
\section{Proposition 3: Home wealth share and NFA revaluation}
\label{sec:prop3}

This section proves Proposition~3 of KL: how a flight-to-safety shock
$\dev{\cy}_1>0$ redistributes financial wealth across countries on impact, and
how it revalues Home's net foreign assets. Unlike Propositions~1--2, here the
portfolios are \emph{not} identical: we allow Home to be levered long in capital,
to hold a nonzero Foreign-bond position, and the Home government to have issued
safe debt. The result is the realized wealth-share identity that the engine
already records as~\eqref{eq:wsh-realized} (their app.~eq.~(92)); the work of the
proof is to expose its three components and to combine them with the realized
returns of Proposition~\ref{prop:returns}.

\subsection{Statement}

\begin{proposition}[Home wealth share and NFA revaluation]
\label{prop:wealth}
Consider the simplified economy of Definition~\ref{def:simplified} with portfolios
initially close to the symmetric benchmark and with Home government safe debt
$b^g_{H,s}$ not too far from zero, and consider a flight-to-safety shock
$\dev{\cy}_1>0$ in period~$1$ with $\widetilde{\widetilde{\dev k}}_1=0$ on
impact. Then, to first order, Home's financial wealth share is revalued on impact
by
\begin{equation}
\dev{\wsh}_1 = \Bigl(\frac{\qk\,k}{\sav}-1\Bigr)\bigl(\rk_1-\rr_1\bigr)
+\frac{b_F}{\sav}\bigl(\fr{r}_1-\D\dev{\rer}_1-\rr_1\bigr)
-\disc\,\frac{\zs}{1+\zs}\,\frac{b^g_{H,s}}{\sav}\,\dev{\cy}_1, \label{eq:prop3-wsh}
\end{equation}
where $\qk k$, $b_F$, and $b^g_{H,s}$ are the steady-state Home positions in
capital, Foreign bonds, and government safe debt (as shares of Home saving
$\sav$), and the realized returns $\rk_1,\rr_1,\fr{r}_1,\D\dev{\rer}_1$ are those
of Proposition~\ref{prop:returns}. Home's net foreign assets $\widehat{\nfa}_1$
are revalued in the same direction.
\end{proposition}

\noindent (We transcribe the proposition to match KL's main-paper statement
(p.~1664). As in Proposition~\ref{prop:returns}, $\D\dev{\rer}_1=\dev{\rer}_1$
because $\dev{\rer}_0=0$ at the deterministic steady state.)

\subsection{Proof}

\paragraph{The identity is the engine's realized wealth-share linearization.}
Equation~\eqref{eq:prop3-wsh} is \emph{verbatim} the engine's realized
wealth-share reduced form~\eqref{eq:wsh-realized}, which is KL's app.~eq.~(92).
We reproduce here the derivation of that reduced form (their derivation of (92)),
so that the three coefficients in \eqref{eq:prop3-wsh} are exposed, and then sign
each component using Proposition~\ref{prop:returns}.

\subparagraph{Step 1: linearize the wealth-share law of motion.} Home's financial
wealth share evolves according to the accounting identity that next period's Home
wealth equals the gross return on this period's Home portfolio, divided by the
gross return on world wealth (the setup's wealth-share law of motion; their
app.~eq.~(51), in re-scaled form). Writing $\dev{\wsh}_1$ as the deviation of the
realized period-$1$ wealth share from its steady-state value and linearizing the
portfolio return around the symmetric steady state, the realized wealth share is
the steady-state-weighted sum of the \emph{excess} realized returns Home earns on
each asset it is differentially exposed to, plus the seigniorage Home earns on the
safe debt it has issued:
\begin{equation}
\dev{\wsh}_1 = \sum_{j}\frac{x_j}{\sav}\bigl(R^{j}_1 - R^{\text{ref}}_1\bigr)
\;+\;(\text{seigniorage on Home safe debt}), \label{eq:prop3-decomp-schema}
\end{equation}
where $x_j$ is the steady-state Home position in asset $j$, $R^{j}_1$ its realized
return, and $R^{\text{ref}}_1$ the return on the reference (safe dollar bond)
funding leg. The three differential exposures in the simplified economy are
(i) leverage in capital, (ii) the Foreign-bond position, and (iii) Home-issued
safe government debt; we take them in turn.
\gapflag{KL do not display the linearization of (51) into (92) term by term; they
write (92) directly as ``log-linearizing the realized evolution of Home's wealth
share implies.'' The schema~\eqref{eq:prop3-decomp-schema} is the standard
budget-constraint linearization for a wealth-share / NFA accounting identity (cf.\
the Devereux--Sutherland portfolio method used in the setup, Section~\ref{sec:setup}):
to first order, the change in a portfolio's relative value is the
steady-state-weighted sum of excess realized returns on the legs in which the
portfolio differs from the reference benchmark. The three coefficients
$\frac{\qk k}{\sav}-1$, $\frac{b_F}{\sav}$, and
$-\disc\frac{\zs}{1+\zs}\frac{b^g_{H,s}}{\sav}$ are the steady-state portfolio
weights that multiply each excess return, evaluated at the symmetric benchmark.
We reconstruct the assignment of each weight to its excess-return leg below; the
weights themselves are read from KL's eq.~(92).}

\subparagraph{Step 2: the leverage-in-capital term.} Home holds steady-state
capital $\qk k$ financed partly by issuing safe (dollar) claims. Its net long
position in capital, relative to a portfolio fully funded by the safe leg, is
$\qk k-\sav$ (its capital holdings net of its own wealth $\sav$); as a share of
$\sav$ this is $\frac{\qk k}{\sav}-1$. The realized excess return on this levered
position is the realized capital return over the dollar-bond funding return,
$\rk_1-\rr_1$. Hence the first term of \eqref{eq:prop3-wsh},
\begin{equation}
\Bigl(\frac{\qk k}{\sav}-1\Bigr)\bigl(\rk_1-\rr_1\bigr). \label{eq:prop3-T1}
\end{equation}
Sign on impact: with Home levered long capital, $\frac{\qk k}{\sav}-1>0$; and from
Proposition~\ref{prop:returns}(a)--(b), the realized excess return on capital is
negative, $\rk_1-\rr_1<0$ (zero capital return minus a positive dollar-bond
return under flexible wages, or a negative capital return minus a positive
dollar-bond return under predetermined wages). Therefore
\eqref{eq:prop3-T1}$<0$: Home \emph{loses} wealth share because the risky capital
it is levered into underperforms the dollar bonds it is short.

\subparagraph{Step 3: the Foreign-bond term.} Home holds steady-state Foreign
bonds $b_F$ (financed at the dollar-bond rate). The realized excess return on this
position, in Home consumption units, is the Home-currency Foreign-bond return over
the dollar-bond return, $\fr{r}_1-\D\dev{\rer}_1-\rr_1$ (the $\D\dev{\rer}_1$
converts the Foreign return into Home units). As a share of $\sav$ the position is
$\frac{b_F}{\sav}$, giving the second term of \eqref{eq:prop3-wsh},
\begin{equation}
\frac{b_F}{\sav}\bigl(\fr{r}_1-\D\dev{\rer}_1-\rr_1\bigr). \label{eq:prop3-T2}
\end{equation}
Sign on impact: from Proposition~\ref{prop:returns}(c),
$\fr{r}_1-\D\dev{\rer}_1-\rr_1<0$ (Foreign bonds underperform dollar bonds, both
because of the dollar's deflation and, under nominal rigidity, the real dollar
appreciation). With $b_F>0$ (Home long Foreign bonds), \eqref{eq:prop3-T2}$<0$:
Home loses wealth on its Foreign-bond position as well.

\subparagraph{Step 4: the seigniorage term.} The Home government has issued safe
dollar debt $b^g_{H,s}$ (a \emph{negative} position from Home's private-sector
perspective: $b^g_{H,s}<0$ means Home is a net issuer of safe claims). A share
$\frac{\zs}{1+\zs}$ of this safe debt is held by the Rest of the World. When the
safety shock raises the convenience yield $\dev{\cy}_1$, the issuer of safe debt
earns \emph{seigniorage}: it borrows at a convenience-discounted rate, and the
discount widens by $\dev{\cy}_1$ to first order. The Home private sector captures
the seigniorage on the share held abroad, discounted by $\disc$ (the safe debt
pays off next period). This is the third term of \eqref{eq:prop3-wsh},
\begin{equation}
-\disc\,\frac{\zs}{1+\zs}\,\frac{b^g_{H,s}}{\sav}\,\dev{\cy}_1. \label{eq:prop3-T3}
\end{equation}
Sign on impact: when Home is a net issuer, $b^g_{H,s}<0$, so
$-\disc\frac{\zs}{1+\zs}\frac{b^g_{H,s}}{\sav}>0$, and with $\dev{\cy}_1>0$ this
term is \emph{positive}: Home \emph{gains} wealth share through the seigniorage it
earns on the safe debt it has issued to the Rest of the World.
\gapflag{The coefficient $-\disc\frac{\zs}{1+\zs}\frac{b^g_{H,s}}{\sav}$ is read
directly from KL's eq.~(92); we supply its interpretation. The factor
$\frac{\zs}{1+\zs}$ is the RoW population share (so it is the fraction of Home
safe debt held \emph{abroad}, where the convenience rent is a transfer to Home
rather than an internal wash); the discount $\disc$ converts the next-period
payoff to present value; and $\dev{\cy}_1$ is the first-order widening of the
convenience discount. The convenience term enters the household Euler equation as
$(1+r_{t+1})/(1-\cy_t)$ (Section~\ref{sec:engine}, discussion after
\eqref{eq:euler-r}), so a one-unit rise in $\dev{\cy}_1$ lowers the real cost of
the safe debt by one unit to first order; the issuer's private-sector wealth
gains $\disc\frac{\zs}{1+\zs}\frac{|b^g_{H,s}|}{\sav}\dev{\cy}_1$. We reconstruct
the sign and economic content; the magnitude is KL's.}

\subparagraph{Step 5: assemble.} Summing \eqref{eq:prop3-T1}, \eqref{eq:prop3-T2},
and \eqref{eq:prop3-T3} reproduces \eqref{eq:prop3-wsh}, i.e.\ the engine's
\eqref{eq:wsh-realized}:
\[
\dev{\wsh}_1 = \underbrace{\Bigl(\tfrac{\qk k}{\sav}-1\Bigr)\bigl(\rk_1-\rr_1\bigr)}_{<0}
+\underbrace{\tfrac{b_F}{\sav}\bigl(\fr{r}_1-\D\dev{\rer}_1-\rr_1\bigr)}_{<0}
\underbrace{-\disc\tfrac{\zs}{1+\zs}\tfrac{b^g_{H,s}}{\sav}\dev{\cy}_1}_{>0\text{ when }b^g_{H,s}<0},
\]
the realized wealth-share revaluation, with the first two components negative and
the third positive (for a net issuer) on impact of a positive safety shock.

\paragraph{The ``close to the symmetric benchmark'' qualifier.} The linear
decomposition \eqref{eq:prop3-wsh} is valid because the coefficients
$\frac{\qk k}{\sav}-1$, $\frac{b_F}{\sav}$, and $\frac{b^g_{H,s}}{\sav}$ are
\emph{steady-state} portfolio weights, evaluated at the symmetric benchmark. KL's
qualifier ``portfolios initially close to the symmetric benchmark and $b^g_{H,s}$
not too far from zero'' is exactly what licenses this first-order treatment: it
keeps the deviations $\dev{\wsh}_1$ small enough that the realized returns
$\rk_1,\rr_1,\fr{r}_1,\D\dev{\rer}_1$ --- which were computed in
Proposition~\ref{prop:returns} \emph{at} the symmetric benchmark
($\Eone\dev{\wsh}_2=0$) --- remain valid inputs to \eqref{eq:prop3-wsh} to first
order. Were the steady-state portfolio far from symmetric, the realized returns
would themselves depend on the wealth-share dynamics (through the
$\Eone\dev{\wsh}_2$ terms in the engine), and the clean three-way decomposition
would break.
\gapflag{KL phrase the qualifier as ``portfolios initially close to the symmetric
benchmark.'' We make explicit \emph{why} it is needed: it decouples the
realized-return inputs (computed at symmetry in Proposition~\ref{prop:returns})
from the wealth-share output, so that the portfolio weights enter only as
constant coefficients. This is the standard first-order portfolio logic; KL leave
it implicit.}

\paragraph{NFA revaluation.} Home's net foreign assets are revalued in the same
direction as the wealth share. The engine's NFA reduced form is
\eqref{eq:nfa-solved}, $\widehat{\nfa}_1=\sav[\Eone\dev{\wsh}_2-\frac{\zs}{1+\zs}\frac{1}{1-\cshare}\Eone\dev{\tot}_2]$,
which is the linearization of the NFA definition (their app.~eq.~(87),
\eqref{eq:nfa-raw}). The realized revaluation channel that drives $\dev{\wsh}_1$
in \eqref{eq:prop3-wsh} feeds NFA through the same realized excess returns: a
fall in Home's wealth share on impact is, to first order, a fall in the market
value of Home's net foreign asset position. Concretely, the realized
return-differential terms~\eqref{eq:prop3-T1}--\eqref{eq:prop3-T2} are a
\emph{valuation} loss on Home's external balance sheet (Home's external assets,
levered long risky and long Foreign bonds, lose value relative to its external
liabilities, the safe dollar claims it has issued), exactly the revaluation that
moves $\widehat{\nfa}_1$ down; the seigniorage term~\eqref{eq:prop3-T3} is a
partial offset. Hence ``Home's NFA are revalued in the same way.'' \qed
(\,Proposition~\ref{prop:wealth}\,)

\subsection{Discussion}

\paragraph{Two transfers out, one rent in.} Proposition~3 decomposes the wealth
redistribution in a flight to safety into three pieces. The first two ---
\eqref{eq:prop3-T1} and \eqref{eq:prop3-T2} --- are a direct consequence of
Proposition~\ref{prop:returns}: Home's wealth falls in its \emph{leveraged
capital} position and in its \emph{Foreign-bond} position, because both of these
risky assets underperform the safe dollar bonds that fund them exactly when the
safety shock hits. The third --- \eqref{eq:prop3-T3} --- runs the other way:
Home's wealth \emph{rises} in the safe debt issued by its government, because Home
earns the convenience rent (seigniorage) on the share of that safe debt held by
the Rest of the World. The net revaluation is the sum: on impact, the two
valuation losses dominate for a country that is levered long risky assets and
short safe debt, so Home loses wealth and transfers it to the Rest of the World.

\paragraph{Connection to the target facts: the US as the world's insurer.} This
is the model's account of the United States as the ``world's insurer'' and of the
valuation channel of US external adjustment. The US runs an external balance sheet
that is \emph{long risky, short safe}: it holds foreign equity and risky claims
(here, levered capital and Foreign bonds) financed by issuing safe dollar claims
(Treasuries and dollar bank liabilities) to the rest of the world. In a global
flight to safety, this balance sheet delivers a large \emph{valuation loss} ---
the US transfers wealth to the rest of the world precisely in bad global states,
which is the insurance it provides. Proposition~3 is the structural counterpart of
the empirical NFA-revaluation fact: US net foreign assets fall in value in risk-off
episodes, by an amount governed by the leverage weight $\frac{\qk k}{\sav}-1$ and
the Foreign-bond weight $\frac{b_F}{\sav}$, partially offset by the seigniorage
the US earns on its safe-debt issuance, $-\disc\frac{\zs}{1+\zs}\frac{b^g_{H,s}}{\sav}\dev{\cy}_1$.

\paragraph{The big\_cy reading.} A \emph{large absolute} convenience yield
$\dev{\cy}_1$ does two things to \eqref{eq:prop3-wsh}. Through
Proposition~\ref{prop:returns} it magnifies the realized excess-return terms
\eqref{eq:prop3-T1}--\eqref{eq:prop3-T2} (larger deflation and real appreciation
$\Rightarrow$ larger underperformance of risky assets), deepening the valuation
transfer the US makes as the world's insurer. But it also magnifies the
seigniorage term \eqref{eq:prop3-T3}: the larger the absolute convenience yield,
the more the US earns on its safe-debt issuance, partially insulating its wealth.
The balance between these is precisely the quantitative question big\_cy poses ---
whether a DiTella-magnitude convenience yield lets the model match the observed
US insurance role and home-currency bias \emph{without} the counterfactual US
risk-tolerance asymmetry that Kekre and Lenel otherwise require.


% ---- 7. pricing kernels, portfolios, risk premia engine (Appendix B.2) ----
% =====================================================================
% 07_portfolio_engine.tex — Appendix B.2: pricing kernels, portfolios,
% risk premia (second-order portfolio engine) (WP-D)
% =====================================================================
\section{The portfolio and risk-premium engine (second-order)}
\label{sec:portfolio-engine}

The first-order engine of Section~\ref{sec:engine} determines prices, returns,
and the realized wealth-share revaluation \emph{given} the steady-state
portfolio $(\qk k,b_F,b^g_{H,s})$. It does not pin that portfolio down: to first
order, all assets are perfect substitutes and the portfolio is indeterminate.
KL close the model with the Devereux--Sutherland (2011) device: the
\emph{equilibrium} portfolio is the one that makes the \emph{second-order}
optimal-portfolio (excess-return) conditions hold, and to leading order it is
determined by a first-order risk-sharing condition. This section reconstructs
KL's online Appendix~B.2 (their app.~pp.~73--75). The outputs are: (i) the
second-order portfolio-Euler conditions; (ii) the log pricing-kernel deviation
$\dev{\pk}_{0,1}$ and its Home and Foreign ``risk-adjusted'' forms
$\dev{\pk}_{0,1}-\rra\tel^z\dev{\epsilon}^z_1$ and
$\fr{\dev{\pk}}_{0,1}-\rraF\tel^z\dev{\epsilon}^z_1$; (iii) the
local-completeness (efficient-risk-sharing) condition; and (iv) the
international-portfolio equation, KL's eq.~(99), with its composite coefficient
$\Gam$. These feed directly into Propositions~\ref{prop:portfolios}
and~\ref{prop:riskpremium}.

Throughout we follow KL's app.~B.2 in writing $\tel^z$ for the loading of the
productivity surprise (their $\sigma^z$; macro \texttt{\textbackslash tel}$^z$ in
the engine), $\dev{\epsilon}^z_1$ for the period-$1$ productivity innovation, and
$\lwedge$ for the steady-state \emph{labor wedge}. As in the engine we keep
$\hbias$ general until completeness is imposed.

\subsection{The Devereux--Sutherland portfolio-Euler conditions}
\label{sec:portfolio-euler}

A second-order approximation of the period-$0$ optimal portfolio-choice
conditions (their app.~p.~73) gives, for the excess return of capital over the
safe dollar bond and for the excess return of the Foreign bond over the safe
dollar bond,
\begin{align}
\Ezero\bigl[\rk_1-\rr_1\bigr] + \text{Jensen terms}
&= \dev{\cy}_0 - \Ezero\Bigl[\bigl(\dev{\pk}_{0,1}-\rra\tel^z\dev{\epsilon}^z_1\bigr)\bigl(\rk_1-\rr_1\bigr)\Bigr], \label{eq:ds-capital}\\
\Ezero\bigl[\fr{r}_1-\D\dev{\rer}_1-\rr_1\bigr] + \text{Jensen terms}
&= \dev{\cy}_0 - \Ezero\Bigl[\bigl(\dev{\pk}_{0,1}-\rra\tel^z\dev{\epsilon}^z_1\bigr)\bigl(\fr{r}_1-\D\dev{\rer}_1-\rr_1\bigr)\Bigr], \label{eq:ds-fbond}
\end{align}
and analogously for the Foreign household (with $\fr{\dev{\pk}}_{0,1}$ and
$\rraF$ in place of $\dev{\pk}_{0,1}$ and $\rra$). Here $\dev{\pk}_{0,1}$ is the
log deviation of the Home household's real pricing kernel (period~$0$ to
period~$1$), defined in Subsection~\ref{sec:log-kernel} below, and
$\dev{\cy}_0$ is the period-$0$ (steady-state) convenience yield on the safe
dollar bond, which appears because the dollar bond yields liquidity services that
capital and Foreign bonds do not.

\paragraph{Reading the conditions.} The left-hand side of each equation is an
expected excess return plus a \emph{Jensen} (variance-of-returns) correction. KL
note (their app.~p.~73) that ``\emph{Jensen terms} reflect the component of
excess returns which do not reflect safety shocks nor the covariance with agents'
pricing kernels (and instead reflect the variance of returns).'' The
economically active part of each condition is therefore the right-hand side: the
convenience premium $\dev{\cy}_0$ that the safe bond commands, net of the
covariance of the (risk-adjusted) pricing kernel with the excess return. The
equilibrium portfolio is the one that makes these covariances take the values
required for \eqref{eq:ds-capital}--\eqref{eq:ds-fbond} to hold. Because the
Jensen terms are pure second moments of returns that are common across the
household problems and do not load on the safety shock or on the kernel
covariance, KL's strategy is to read the equilibrium holdings off the
\emph{first-order} risk-sharing condition that the covariance terms encode; the
Jensen terms then determine only the (common, second-order) level of expected
excess returns, not the cross-sectional portfolio. We make this precise via the
local-completeness condition of Subsection~\ref{sec:local-completeness}.

\subsection{The log pricing kernel}
\label{sec:log-kernel}

In period~$1$, the log deviation of the representative Home household's pricing
kernel is (their app.~p.~73)
\begin{equation}
\dev{\pk}_{0,1} = -\dev{\rs c}_1 + (1-\rra)\,\dev{\rs \val}_1, \label{eq:m01}
\end{equation}
which is the log-linearization of the Epstein--Zin/recursive stochastic discount
factor: with unit IES (the simplified environment), the kernel is
$m_{0,1}\propto (c_1)^{-1}(v_1/ce_0)^{1-\rra}$, whose log deviation is
$-\dev{\rs c}_1$ plus $(1-\rra)$ times the (re-scaled) value deviation
$\dev{\rs\val}_1$ (the certainty-equivalent normalizer $ce_0$ is non-stochastic
as of period~$0$ and drops from the period-$0$-to-$1$ innovation).
\gapflag{KL display \eqref{eq:m01} directly as ``the log deviation in the
representative Home household's pricing kernel is given by
$\dev{\pk}_{0,1}=-\dev{\rs c}_1+(1-\rra)\dev{\rs\val}_1$.'' We supply the
one-line provenance: under unit IES the recursive kernel is the product of the
consumption-growth term (exponent $-1$ on $c$, since $1/\text{IES}=1$) and the
continuation-value/certainty-equivalent term carrying the risk-aversion exponent
$1-\rra$. The macro renders $\dev{\rs c}_1=\hat{\tilde c}_1$ and
$\dev{\rs\val}_1=\hat{\tilde v}_1$, the period-$1$ re-scaled (stationarized)
consumption and value deviations of Section~\ref{sec:setup}.}

\paragraph{The value linearization.} The (re-scaled) recursive value linearizes
to (their app.~p.~73)
\begin{equation}
\dev{\rs \val}_1 = (1-\disc)\dev{\rs c}_1 - (1-\disc)(1-\lwedge)(1-\cshare)\dev{\ell}_1 + \disc\,\dev{\rs{ce}}_1, \label{eq:v1}
\end{equation}
where $\lwedge$ is the labor wedge in the deterministic steady state and
$\dev{\rs{ce}}_1$ is the log deviation of the certainty equivalent of the
period-$2$ continuation value. The term
$-(1-\disc)(1-\lwedge)(1-\cshare)\dev{\ell}_1$ is the flow disutility of labor,
valued at the steady-state \emph{after-wedge} marginal rate; the factor
$(1-\lwedge)$ scales the labor term down by the wedge because, with a positive
labor wedge, the private marginal disutility of work is below the social marginal
product, so a unit of employment contributes
$(1-\lwedge)(1-\cshare)$ to flow utility (in consumption units) rather than the
full $(1-\cshare)$.
\gapflag{KL display \eqref{eq:v1} as
$\dev{\rs\val}_1=(1-\disc)\dev{\rs c}_1-(1-\disc)(1-\lwedge)(1-\cshare)\dev{\ell}_1+\disc\dev{\rs{ce}}_1$.
The labor coefficient $(1-\lwedge)(1-\cshare)$ is KL's; we supply the reading.
Note this is the only place the steady-state labor wedge $\lwedge$ enters the
kernel, and it is precisely the channel that makes \emph{employment} (not only
productivity and wealth) a priced state variable in Props~4--5; it vanishes if
$\lwedge=0$.}

\paragraph{The certainty equivalent.} Because there is no uncertainty after
period~$1$, the period-$1$ certainty equivalent equals the period-$2$ expected
value (their app.~p.~74),
\begin{equation}
\dev{\rs{ce}}_1 = \Eone\dev{\rs \val}_2 = \atil\,\Eone\dev{\wsh}_2 + \cshare\,\widetilde{\widetilde{\dev k}}_1, \label{eq:ce1}
\end{equation}
where the second equality is the engine's period-$2$ value
relation~\eqref{eq:Ev2} (their app.~eq.~(77)). To first order the
certainty-equivalent operator is the identity here, since the only remaining
randomness (the period-$1$ innovations) is already realized at the point where
$\dev{\rs{ce}}_1$ is evaluated.

\paragraph{Assembling $\dev{\rs c}_1$ from the engine.} The period-$1$ Home
consumption is the engine's reduced form~\eqref{eq:c1-solved}, which we restate
here in the app.~B.2 notation (writing the recurring denominator from the engine
explicitly rather than as $\Lambda$, to match KL's displayed lines):
\begin{equation}
\dev{\rs c}_1 = \atil\,\Eone\dev{\wsh}_2
+(1-\cshare)\Bigl(\tfrac{1}{1+\zs}\dev{\ell}_1+\tfrac{\zs}{1+\zs}\fr{\dev{\ell}}_1\Bigr)
+\cshare\,\widetilde{\widetilde{\dev k}}_1
-\frac{\hbias}{\frac{\cshare}{1-\cshare}+\tel+(1-\tel)\left(\hbias\frac{1+\zs}{\zs}\right)^2}\bigl(\fr{\dev{\ell}}_1-\dev{\ell}_1\bigr). \label{eq:c1-b2}
\end{equation}
Equation~\eqref{eq:c1-b2} is identically \eqref{eq:c1-solved} with the
denominator $\Lambda$ \eqref{eq:Lambda-def} written out; KL display it in this
expanded form on app.~p.~73.

\paragraph{The combined value deviation.} Substituting \eqref{eq:c1-b2} and
\eqref{eq:ce1} into the value linearization~\eqref{eq:v1} and collecting terms,
KL obtain (their app.~p.~74)
\begin{multline}
\dev{\rs \val}_1 = \atil\,\Eone\dev{\wsh}_2 - \cshare\,\tel^z\dev{\epsilon}^z_1
+\lwedge(1-\disc)(1-\cshare)\dev{\ell}_1\\
+(1-\disc)\left((1-\cshare)\frac{\zs}{1+\zs}
-\frac{\hbias}{\frac{\cshare}{1-\cshare}+\tel+(1-\tel)\left(\hbias\frac{1+\zs}{\zs}\right)^2}\right)\bigl(\fr{\dev{\ell}}_1-\dev{\ell}_1\bigr). \label{eq:v1-combined}
\end{multline}
\gapflag{\eqref{eq:v1-combined} is KL's displayed combined value line (app.~p.~74).
The reconstruction: write \eqref{eq:v1} as
$\dev{\rs\val}_1=(1-\disc)\dev{\rs c}_1-(1-\disc)(1-\lwedge)(1-\cshare)\dev{\ell}_1+\disc\dev{\rs{ce}}_1$.
The $\disc\dev{\rs{ce}}_1=\disc[\atil\Eone\dev{\wsh}_2+\cshare\widetilde{\widetilde{\dev k}}_1]$
piece and the $(1-\disc)\atil\Eone\dev{\wsh}_2$ piece inside $(1-\disc)\dev{\rs c}_1$
combine to the full $\atil\Eone\dev{\wsh}_2$ (coefficients $\disc+(1-\disc)=1$);
likewise the two $\cshare\widetilde{\widetilde{\dev k}}_1$ pieces combine to
$\cshare\widetilde{\widetilde{\dev k}}_1$. KL then \emph{re-express} the symmetric
employment piece using the impact wage-setting relation
\eqref{eq:predet-wage}: under wages set one period in advance the common labor
movement satisfies $(1-\cshare)(\frac{1}{1+\zs}\dev{\ell}_1+\frac{\zs}{1+\zs}\fr{\dev{\ell}}_1)+\cshare\widetilde{\widetilde{\dev k}}_1$
contributes a $-\cshare\tel^z\dev{\epsilon}^z_1$ term once the productivity
surprise is substituted (the same $\tel^z\dev{\epsilon}^z_1$ that loads the wage
equation), so the aggregate-employment + capital block collapses into the
$-\cshare\tel^z\dev{\epsilon}^z_1$ displayed term. The labor coefficient
collects as follows: $(1-\disc)\dev{\rs c}_1$ contributes
$+(1-\disc)(1-\cshare)\frac{1}{1+\zs}$ to $\dev{\ell}_1$ and the relative-labor
$-\hbias/[\cdots]$ term; the explicit disutility term contributes
$-(1-\disc)(1-\lwedge)(1-\cshare)\dev{\ell}_1$. Adding the
$+(1-\disc)(1-\cshare)\dev{\ell}_1$ implicit in the symmetric piece, the
$\dev{\ell}_1$ coefficient nets to $+\lwedge(1-\disc)(1-\cshare)$ (the
$-(1-\lwedge)$ and $+1$ combine to $+\lwedge$), which is the displayed
$\lwedge(1-\disc)(1-\cshare)\dev{\ell}_1$ wedge term; the residual relative-labor
piece collects into the displayed
$(1-\disc)\bigl((1-\cshare)\frac{\zs}{1+\zs}-\hbias/[\cdots]\bigr)(\fr{\dev{\ell}}_1-\dev{\ell}_1)$
coefficient. This collection is the single tersest algebra step in B.2; we have
reproduced KL's displayed coefficients and flag that we filled the term-by-term
collection rather than re-deriving each cancellation from the primitive Euler
block.}

\paragraph{The Home risk-adjusted kernel.} Substituting \eqref{eq:c1-b2} and
\eqref{eq:v1-combined} into the kernel definition~\eqref{eq:m01}, and grouping the
productivity loading on the kernel as $\rra\tel^z\dev{\epsilon}^z_1$, KL obtain
the central B.2 expression (their app.~p.~74)
\begin{multline}
\dev{\pk}_{0,1} - \rra\,\tel^z\dev{\epsilon}^z_1
= -\rra\Bigl[\atil\,\Eone\dev{\wsh}_2 + (1-\cshare)\tel^z\dev{\epsilon}^z_1\Bigr]
- (1-\cshare)\Bigl(\tfrac{1}{1+\zs}\dev{\ell}_1+\tfrac{\zs}{1+\zs}\fr{\dev{\ell}}_1\Bigr)\\
+\frac{\hbias}{\frac{\cshare}{1-\cshare}+\tel+(1-\tel)\left(\hbias\frac{1+\zs}{\zs}\right)^2}\bigl(\fr{\dev{\ell}}_1-\dev{\ell}_1\bigr)\\
+(1-\rra)(1-\disc)\left[\lwedge(1-\cshare)\dev{\ell}_1
+\left((1-\cshare)\frac{\zs}{1+\zs}-\frac{\hbias}{\frac{\cshare}{1-\cshare}+\tel+(1-\tel)\left(\hbias\frac{1+\zs}{\zs}\right)^2}\right)\bigl(\fr{\dev{\ell}}_1-\dev{\ell}_1\bigr)\right]. \label{eq:m01-home}
\end{multline}
\gapflag{\eqref{eq:m01-home} is KL's displayed Home risk-adjusted-kernel line
(app.~p.~74). The reconstruction is purely mechanical given \eqref{eq:m01},
\eqref{eq:c1-b2}, \eqref{eq:v1-combined}: $\dev{\pk}_{0,1}=-\dev{\rs c}_1+(1-\rra)\dev{\rs\val}_1$.
The first bracketed group $-\rra[\atil\Eone\dev{\wsh}_2+(1-\cshare)\tel^z\dev{\epsilon}^z_1]$
arises because the $\atil\Eone\dev{\wsh}_2$ term appears in $\dev{\rs c}_1$ with
coefficient $-1$ and in $(1-\rra)\dev{\rs\val}_1$ with coefficient $(1-\rra)$,
netting to $-\rra\atil\Eone\dev{\wsh}_2$; we then \emph{move} the
$\rra\tel^z\dev{\epsilon}^z_1$ loading to the left-hand side (KL's convention,
isolating the part of the kernel orthogonal to the directly-priced productivity
risk), which converts the $-\cshare\tel^z\dev{\epsilon}^z_1$ pieces of
$\dev{\rs c}_1$ and $\dev{\rs\val}_1$ into the displayed
$-\rra(1-\cshare)\tel^z\dev{\epsilon}^z_1$ (using $-(-\cshare)-(1-\rra)\cshare+\rra=\rra(1-\cshare)$
on the $\tel^z\dev{\epsilon}^z_1$ coefficient). The symmetric-employment and
relative-labor terms from $-\dev{\rs c}_1$ appear with the signs shown, and the
$(1-\rra)(1-\disc)[\cdots]$ block is exactly $(1-\rra)$ times the
wedge-and-relative-labor part of \eqref{eq:v1-combined}. We reproduce KL's line
term-for-term; the only judgement call is the left-hand-side grouping of
$\rra\tel^z\dev{\epsilon}^z_1$, which is KL's own (it is the quantity that enters
the Devereux--Sutherland conditions \eqref{eq:ds-capital}--\eqref{eq:ds-fbond}).}

\paragraph{The Foreign risk-adjusted kernel.} Analogous steps in Foreign
(replacing $\rra\to\rraF$, the wealth-share loading $\atil\to-\frac{1}{\zs}\atil$
from the period-$2$ Foreign value~\eqref{eq:Ev2-star}, and the symmetric-labor
weight as appropriate) yield (their app.~p.~74)
\begin{multline}
\fr{\dev{\pk}}_{0,1} - \rraF\,\tel^z\dev{\epsilon}^z_1
= -\rraF\Bigl[-\tfrac{1}{\zs}\atil\,\Eone\dev{\wsh}_2 + (1-\cshare)\tel^z\dev{\epsilon}^z_1\Bigr]
- (1-\cshare)\Bigl(\tfrac{1}{1+\zs}\dev{\ell}_1+\tfrac{\zs}{1+\zs}\fr{\dev{\ell}}_1\Bigr)\\
-\frac{1}{\zs}\frac{\hbias}{\frac{\cshare}{1-\cshare}+\tel+(1-\tel)\left(\hbias\frac{1+\zs}{\zs}\right)^2}\bigl(\fr{\dev{\ell}}_1-\dev{\ell}_1\bigr)\\
+(1-\rraF)(1-\disc)\left[\lwedge(1-\cshare)\fr{\dev{\ell}}_1
-\left((1-\cshare)\frac{1}{1+\zs}-\frac{1}{\zs}\frac{\hbias}{\frac{\cshare}{1-\cshare}+\tel+(1-\tel)\left(\hbias\frac{1+\zs}{\zs}\right)^2}\right)\bigl(\fr{\dev{\ell}}_1-\dev{\ell}_1\bigr)\right]. \label{eq:m01-foreign}
\end{multline}
\gapflag{\eqref{eq:m01-foreign} is KL's displayed Foreign risk-adjusted-kernel
line (app.~p.~74). The reconstruction mirrors the Home case under the Foreign
relabeling: the Foreign household's period-$2$ value carries
$-\frac{1}{\zs}\atil\Eone\dev{\wsh}_2$ (\eqref{eq:Ev2-star}), so its
wealth-share loading flips sign and scales by $1/\zs$; the relative-labor
coefficient picks up the $1/\zs$ that distinguishes the Foreign cross-country
weight; the wedge term carries $\fr{\dev{\ell}}_1$ rather than $\dev{\ell}_1$.
We reproduce KL's displayed signs; the structural mirror to \eqref{eq:m01-home}
is exact, and we flag that we transcribed the Foreign weights ($1/\zs$ factors,
$\frac{1}{1+\zs}$ vs.\ $\frac{\zs}{1+\zs}$) from KL's line rather than
re-deriving the Foreign value linearization from scratch.}

\subsection{Local completeness and the risk-sharing condition}
\label{sec:local-completeness}

KL observe (their app.~p.~74) that the simplified environment is \emph{locally
complete} in the sense of Coeurdacier and Gourinchas (2016): with three traded
assets (safe dollar bond, capital, Foreign bond) and two shocks (safety
$\dev{\cy}$, productivity $\dev{\epsilon}^z$), the asset span is rich enough that
the equilibrium portfolios implement constrained-efficient risk sharing to first
order. The implication is the \emph{risk-sharing condition}
\begin{equation}
\boxed{\;\dev{\pk}_{0,1} - \rra\,\tel^z\dev{\epsilon}^z_1
= \fr{\dev{\pk}}_{0,1} - \rraF\,\tel^z\dev{\epsilon}^z_1 + \D\dev{\rer}_1.\;} \label{eq:risk-sharing}
\end{equation}
This is the open-economy Backus--Smith / law-of-one-price relation for stochastic
discount factors: the Home and Foreign risk-adjusted kernels must agree up to the
real-exchange-rate change $\D\dev{\rer}_1$, which converts Foreign marginal
utility into Home units. Equivalently, the two countries' marginal valuations of
state-contingent wealth coincide once expressed in a common numeraire --- the
hallmark of (locally) complete markets.

\paragraph{Why this pins the portfolio.} The kernels \eqref{eq:m01-home} and
\eqref{eq:m01-foreign} are functions of the endogenous outcomes
$(\Eone\dev{\wsh}_2,\dev{\ell}_1,\fr{\dev{\ell}}_1)$, which in turn depend on the
portfolio through the realized wealth-share identity \eqref{eq:wsh-realized}
(their eq.~(92)) and the labor block. Imposing \eqref{eq:risk-sharing} therefore
delivers one scalar equation linking the portfolio weights to the shocks; this is
the equation that the international-portfolio relation below makes explicit. The
Jensen terms of \eqref{eq:ds-capital}--\eqref{eq:ds-fbond} do not enter
\eqref{eq:risk-sharing} because they are common across the Home and Foreign
conditions and cancel in the difference --- which is precisely why KL can read the
equilibrium portfolio off the first-order risk-sharing condition rather than the
full second-order excess-return level.
\gapflag{KL state \eqref{eq:risk-sharing} as the local-completeness implication
``$\dev{\pk}_{0,1}-\rra\tel^z\dev{\epsilon}^z_1=\fr{\dev{\pk}}_{0,1}-\rraF\tel^z\dev{\epsilon}^z_1+\D\dev{\rer}_1$''
citing Coeurdacier and Gourinchas (2016) for the local-completeness notion. We
supply the interpretation (Backus--Smith with $\D\dev{\rer}_1$ as the numeraire
conversion) and the explanation of why imposing it pins the portfolio while the
Jensen terms drop. We do not re-derive the Coeurdacier--Gourinchas completeness
result; we take it as KL do, namely as the statement that three assets span the
two shocks plus the redistributive (wealth-share) margin.}

\subsection{The international-portfolio equation (their eq.~99)}
\label{sec:eq99}

Imposing the risk-sharing condition~\eqref{eq:risk-sharing} --- i.e.\ subtracting
\eqref{eq:m01-foreign} from \eqref{eq:m01-home} and setting the difference equal
to $\D\dev{\rer}_1$ --- and substituting the realized wealth-share
identity~\eqref{eq:wsh-realized} for $\Eone\dev{\wsh}_2$ via $\dev{\wsh}_1$
(their app.~eq.~(92)), KL collect terms to obtain the
\emph{international-portfolio equation} (their app.~eq.~(99)):
\begin{multline}
\Bigl(\frac{\qk k}{\sav}-1\Bigr)\bigl(\rk_1-\rr_1\bigr)
+\frac{b_F}{\sav}\bigl(\fr{r}_1-\dev{\rer}_1-\rr_1\bigr)
-\disc\,\frac{\zs}{1+\zs}\,\frac{b^g_{H,s}}{\sav}\,\dev{\cy}_1
=\\
\frac{1}{\Gam}(\rraF-\rra)(1-\cshare)\tel^z\dev{\epsilon}^z_1
+\frac{1}{\Gam}(\rraF-\rra)\lwedge(1-\disc)(1-\cshare)\Bigl(\tfrac{1}{1+\zs}\dev{\ell}_1+\tfrac{\zs}{1+\zs}\fr{\dev{\ell}}_1\Bigr)\\
+(1-\disc)\frac{\zs}{1+\zs}(\tel-1)\frac{1-\cshare}{\cshare}
\frac{1-\left(\hbias\frac{1+\zs}{\zs}\right)^2}{\frac{\cshare}{1-\cshare}+\tel+(1-\tel)\left(\hbias\frac{1+\zs}{\zs}\right)^2}\bigl(\fr{\dev{\ell}}_1-\dev{\ell}_1\bigr)\\
-\frac{1}{\Gam}\Bigl(\tfrac{1}{\zs}(\rraF-1)+(\rra-1)\Bigr)(1-\disc)(1-\lwedge)(1-\cshare)\frac{\zs}{1+\zs}\bigl(\fr{\dev{\ell}}_1-\dev{\ell}_1\bigr)\\
+\frac{1}{\Gam}\Bigl(\tfrac{1}{\zs}(\rraF-1)+(\rra-1)\Bigr)(1-\disc)
\frac{\hbias}{\frac{\cshare}{1-\cshare}+\tel+(1-\tel)\left(\hbias\frac{1+\zs}{\zs}\right)^2}\bigl(\fr{\dev{\ell}}_1-\dev{\ell}_1\bigr), \label{eq:eq99}
\end{multline}
where the composite coefficient is
\begin{equation}
\Gam \;\equiv\; \left[\rra + \frac{1}{\zs}\rraF
+\frac{\hbias^2\left(\frac{1+\zs}{\zs}\right)^3}{\frac{\cshare}{1-\cshare}+\tel\left(1-\left(\hbias\frac{1+\zs}{\zs}\right)^2\right)}\right]\atil. \label{eq:Gamma-def}
\end{equation}
\gapflag{Equations~\eqref{eq:eq99}--\eqref{eq:Gamma-def} are KL's displayed
eq.~(99) and the definition of $\Gam$ (their app.~p.~75). This is the single
densest derivation in B.2 and KL display only the final collected form
(``Substituting in using the above results and those of the previous section and
collecting terms, we obtain''). We reconstruct the structure as follows. (i) The
\emph{left-hand side} is exactly the realized wealth-share
identity~\eqref{eq:wsh-realized}/(92): the three portfolio-weighted excess-return
terms (leverage, Foreign bond, seigniorage) equal $\dev{\wsh}_1=\Eone\dev{\wsh}_2$
on impact (recall \eqref{eq:Ewsh-solved} sets
$\Eone\dev{\wsh}_2=\dev{\wsh}_1+(\text{relative-labor})$). KL move all of
\eqref{eq:wsh-realized} to the LHS so the RHS is the risk-sharing
\emph{requirement} on $\Eone\dev{\wsh}_2$. (ii) The \emph{right-hand side} is the
difference of the two risk-adjusted kernels~\eqref{eq:m01-home}$-$\eqref{eq:m01-foreign}
minus $\D\dev{\rer}_1$, solved for $\Eone\dev{\wsh}_2$. The
$\Eone\dev{\wsh}_2$ coefficient on that difference is
$-\rra\atil-(-\rraF(-\frac{1}{\zs}\atil))=-(\rra+\frac{1}{\zs}\rraF)\atil$
\emph{plus} the $\D\dev{\rer}_1$ contribution: $\D\dev{\rer}_1=\dev{\rer}_1$
carries an $\Eone\dev{\wsh}_2$ loading through \eqref{eq:q1-solved}, namely
$\hbias\frac{1+\zs}{\zs}\cdot\frac{\hbias(\frac{1+\zs}{\zs})^2}{\frac{\cshare}{1-\cshare}+\tel(1-(\hbias\frac{1+\zs}{\zs})^2)}\atil
=\frac{\hbias^2(\frac{1+\zs}{\zs})^3}{\frac{\cshare}{1-\cshare}+\tel(1-(\hbias\frac{1+\zs}{\zs})^2)}\atil$,
which is the third term inside the bracket of $\Gam$. Collecting,
$\Gam=[\rra+\frac{1}{\zs}\rraF+\frac{\hbias^2(\frac{1+\zs}{\zs})^3}{\frac{\cshare}{1-\cshare}+\tel(1-(\hbias\frac{1+\zs}{\zs})^2)}]\atil$
is exactly the coefficient multiplying $\Eone\dev{\wsh}_2$ in the
kernel-difference; dividing through by it gives the $1/\Gam$ on the RHS terms.
(iii) The RHS \emph{numerators} are the remaining (non-$\Eone\dev{\wsh}_2$) terms
of \eqref{eq:m01-home}$-$\eqref{eq:m01-foreign}: the
$(\rraF-\rra)(1-\cshare)\tel^z\dev{\epsilon}^z_1$ productivity term (from the
$-\rra(1-\cshare)$ vs.\ $-\rraF(1-\cshare)$ loadings on $\tel^z\dev{\epsilon}^z_1$);
the $(\rraF-\rra)\lwedge(1-\disc)(1-\cshare)(\cdots)$ aggregate-employment term
(from the wedge terms, which carry $(1-\rra)$ vs.\ $(1-\rraF)$ and net to
$(\rraF-\rra)\lwedge$ on the symmetric-employment piece); and the
relative-labor $(\fr{\dev{\ell}}_1-\dev{\ell}_1)$ terms, which split into the
$(\tel-1)$ ``terms-of-trade revaluation'' piece (matching the $\Eone\dev{\wsh}_2$
relative-labor coefficient of \eqref{eq:Ewsh-solved}, carried over because the
LHS used \eqref{eq:wsh-realized} while the RHS prices $\Eone\dev{\wsh}_2$) and the
$(\frac{1}{\zs}(\rraF-1)+(\rra-1))$ pieces from the $(1-\rra)$/$(1-\rraF)$
relative-labor coefficients. We have matched KL's displayed eq.~(99)
term-by-term: the five RHS lines correspond, in order, to the productivity term,
the aggregate-employment term, the $(\tel-1)$ terms-of-trade-revaluation term,
and the two relative-labor wedge/home-bias terms. We flag this as a \emph{genuine
reconstruction}: KL display only the collected result, and while every
coefficient below matches their eq.~(99), we filled the intermediate
kernel-difference algebra and the identification of the $\Gam$ bracket with the
$\Eone\dev{\wsh}_2$ loading on $\dev{\rer}_1$. The Props~4--5 proofs use only the
\emph{structure} of \eqref{eq:eq99} (which RHS terms vary with $\dev{\epsilon}^z$
vs.\ $\dev{\cy}$, and the signs of their coefficients), not the exact algebra of
every term.}

\paragraph{The three drivers of Home wealth.} The left-hand side of
\eqref{eq:eq99} is the realized revaluation of Home's wealth share
(Proposition~\ref{prop:wealth}); the right-hand side is what efficient risk
sharing \emph{requires} that revaluation to be. KL read off three economic
drivers (their app.~p.~75); international risk sharing calls for Home wealth (the
LHS) to rise with:
\begin{itemize}[leftmargin=1.6em,itemsep=3pt]
  \item \textbf{productivity, provided $\rraF>\rra$}: a positive TFP shock
  $\dev{\epsilon}^z_1>0$ raises aggregate production and hence world consumption;
  the more risk-tolerant country (lower risk aversion) should bear more of this
  aggregate risk, so if Home is more risk-tolerant ($\rra<\rraF$) Home wealth
  should rise in good productivity states. This is the
  $\frac{1}{\Gam}(\rraF-\rra)(1-\cshare)\tel^z\dev{\epsilon}^z_1$ term, positive
  when $\rraF>\rra$.
  \item \textbf{aggregate employment, provided $\lwedge(\rraF-\rra)>0$}: with a
  positive labor wedge, an increase in employment raises welfare (the marginal
  product exceeds the private marginal disutility), and again the more
  risk-tolerant country should be exposed to this welfare gain. This is the
  $\frac{1}{\Gam}(\rraF-\rra)\lwedge(1-\disc)(1-\cshare)(\cdots)$ term.
  \item \textbf{Foreign employment less Home employment}, i.e.\ the relative-labor
  gap $\fr{\dev{\ell}}_1-\dev{\ell}_1$, under either of two stated sign
  conditions (their app.~p.~75): (a)
  $(\tel-1)\bigl(1-(\hbias\frac{1+\zs}{\zs})^2\bigr)>0$, ``since this implies that
  Foreign labor income rises relative to Home labor income''; or (b)
  $-\bigl(\frac{1}{\zs}(\rraF-1)+(\rra-1)\bigr)\bigl(\frac{\hbias/[\cdots]}{1}-(1-\lwedge)(1-\cshare)\frac{\zs}{1+\zs}\bigr)>0$,
  ``since this implies that Home requires more wealth when its real exchange rate
  appreciates, despite the larger disutility of labor in Foreign.''
\end{itemize}
KL close B.2 with the statement that the equilibrium risk premium on Foreign
bonds relative to dollar bonds is ``the covariance of (the negative of) the log
deviation in any agent's pricing kernel with the excess log return,'' i.e.\
\begin{equation}
\text{risk premium on Foreign vs.\ dollar bonds}
\;=\; -\Cov_0\bigl(\dev{\pk}_{0,1}-\rra\tel^z\dev{\epsilon}^z_1,\;\fr{r}_1-\D\dev{\rer}_1-\rr_1\bigr), \label{eq:rp-cov}
\end{equation}
which is the object Proposition~\ref{prop:riskpremium} signs. (By local
completeness~\eqref{eq:risk-sharing} the same covariance obtains using the
Foreign household's kernel, so the risk premium is well-defined independent of
whose kernel one prices with.)


% ---- 8. Proposition 4: equilibrium portfolios ----
% =====================================================================
% 08_prop4.tex — Proposition 4: equilibrium portfolios (WP-D)
% =====================================================================
\section{Proposition 4: equilibrium portfolios}
\label{sec:prop4}

This section proves Proposition~4 of KL: the comparative statics of the
equilibrium international portfolio --- Home's holdings of capital (a claim to the
world equity return) and its issuance of safe dollar bonds --- with respect to
the supply of safe dollar debt $b^g_{H,s}$ and to the cross-country risk-tolerance
ratio $\rraF/\rra$. The proof combines the international-portfolio
equation~\eqref{eq:eq99} of Section~\ref{sec:portfolio-engine} with the
\emph{arbitrary-portfolio} labor solutions (generalizing the
Proposition~\ref{prop:prices} labor block away from symmetric portfolios), and
then applies the method of undetermined coefficients at the symmetric-portfolio
point.

\subsection{Statement}

\begin{proposition}[Equilibrium portfolios]
\label{prop:portfolios}
Consider the simplified economy of Definition~\ref{def:simplified} with wages set
one period in advance and a positive steady-state labor wedge $\lwedge>0$ in each
country, and impose complete home bias $\hbias\to\zs/(1+\zs)$. Then, at least
around the symmetric-country-portfolio case:
\begin{enumerate}[label=(\alph*),leftmargin=1.8em,itemsep=3pt]
  \item Home's equilibrium portfolio share in capital (the world equity claim) is
  \emph{unaffected} by the supply of safe dollar debt, $\dfrac{dk}{db^g_{H,s}}=0$,
  while Home's portfolio share in dollar bonds \emph{falls} (Home issues more
  dollar bonds) as $-b^g_{H,s}$ rises, $\dfrac{db_H}{db^g_{H,s}}>0$;
  \item holding $\rra+\frac{1}{\zs}\rraF$ fixed, Home's portfolio share in capital
  \emph{rises} as Home becomes more risk-tolerant relative to Foreign,
  $\dfrac{dk}{d[\rraF/\rra]}\Big|_{\rra+\frac{1}{\zs}\rraF}>0$, while Home's
  portfolio share in dollar bonds \emph{falls} (Home issues more dollar bonds),
  $\dfrac{db_H}{d[\rraF/\rra]}\Big|_{\rra+\frac{1}{\zs}\rraF}<0$.
\end{enumerate}
\end{proposition}

\gapflag{We transcribe Proposition~4 to match KL's main-paper statement
(p.~1664) as relayed in the work-package brief; the project does not hold a local
copy of the main-paper PDF, so the verbatim wording is taken from the brief. The
mathematical content of (a)--(b) --- the four signed derivatives
$dk/db^g_{H,s}=0$, $db_H/db^g_{H,s}>0$, $dk/d[\rraF/\rra]>0$,
$db_H/d[\rraF/\rra]<0$ at the symmetric portfolio, the latter two holding
$\rra+\frac{1}{\zs}\rraF$ fixed --- matches the comparative statics displayed in
the appendix proof (their app.~p.~78), which we hold directly. Here $k$ denotes
Home's portfolio share in capital (the levered world-equity position) and $b_H$
Home's portfolio share in safe dollar bonds; ``Home issues more dollar bonds''
corresponds to $b_H$ falling / $-b^g_{H,s}$ rising.}

\subsection{Proof}

\paragraph{Strategy.} The proof has three steps. (i) Solve the labor block for
\emph{arbitrary} portfolios, generalizing the Proposition~\ref{prop:prices}
solution; this expresses $(\dev{\ell}_1,\fr{\dev{\ell}}_1)$ in terms of the shocks
$(\dev{\cy}_1,\dev{\epsilon}^z_1)$ and the forward state $\Eone\dev{\wsh}_2$. (ii)
Substitute into the international-portfolio equation~\eqref{eq:eq99} and into the
realized wealth-share identity~\eqref{eq:wsh-realized} to obtain a fixed-point in
$\Eone\dev{\wsh}_2$, yielding the eq.~(99)-with-$\Gamp$ identity. (iii) At the
symmetric-portfolio point (where a normalization constant $\Delta=1$), match
coefficients on $\dev{\epsilon}^z_1$ and $\dev{\cy}_1$ by the method of
undetermined coefficients to read off the portfolio derivatives.

\subparagraph{Step 1: arbitrary-portfolio labor solutions.} Generalizing the
Proposition~\ref{prop:prices} labor block (their app.~B.3.1) to arbitrary
portfolios, KL show (their app.~p.~77) that the symmetric and relative labor
movements solve
\begin{align}
\frac{1}{1+\zs}\dev{\ell}_1 + \frac{\zs}{1+\zs}\fr{\dev{\ell}}_1
&= -\frac{1}{\cshare+\frac{1}{\taylor}(1-\cshare)}\frac{1}{\taylor}\frac{1}{1+\zs}\dev{\cy}_1
+\frac{1}{\cshare+\frac{1}{\taylor}(1-\cshare)}(1-\cshare)\tel^z\dev{\epsilon}^z_1, \label{eq:lab-sym}\\
\fr{\dev{\ell}}_1 - \dev{\ell}_1
&= \frac{1}{\cshare+\frac{1}{\taylor}(1-\cshare)}\frac{1}{\taylor}\dev{\cy}_1
-\frac{1}{\cshare+\frac{1}{\taylor}(1-\cshare)}\frac{1+\zs}{\zs}\atil\,\Eone\dev{\wsh}_2. \label{eq:lab-rel}
\end{align}
These generalize the Proposition~1 labor solution: \eqref{eq:lab-sym} is the
aggregate-employment response (a positive safety shock contracts global
employment; a positive productivity surprise expands it), and \eqref{eq:lab-rel}
is the relative-employment response (a positive safety shock raises Foreign-less-
Home employment, since the shock hits the Home Taylor rule; and a positive Home
wealth share raises Home relative employment through the demand channel, with the
$\frac{1+\zs}{\zs}\atil$ wealth-share multiplier). The denominator
$\cshare+\frac{1}{\taylor}(1-\cshare)$ is the same $D$ of the
Proposition~\ref{prop:returns} proof.
\gapflag{KL display \eqref{eq:lab-sym}--\eqref{eq:lab-rel} as ``For arbitrary
portfolios, it is straightforward to show that [\dots], generalizing the results
given in the proof of Proposition~1 to arbitrary portfolios.'' We transcribe
their two displayed lines verbatim (app.~p.~77). The only difference from the
Proposition~1 labor block is the appearance of the $\Eone\dev{\wsh}_2$ term in
\eqref{eq:lab-rel}: under symmetric portfolios $\Eone\dev{\wsh}_2=0$ and
\eqref{eq:lab-rel} collapses to the Proposition~1 relative-labor response; for
arbitrary portfolios $\Eone\dev{\wsh}_2\neq0$ and the wealth share feeds back into
relative labor. We did not re-derive the labor block from the firm/wage system
\eqref{eq:firm-H}--\eqref{eq:firm-F}; we take KL's displayed arbitrary-portfolio
solution.}

\subparagraph{Step 2: substitute into the international-portfolio equation.}
Substituting \eqref{eq:lab-sym}--\eqref{eq:lab-rel} into the
international-portfolio equation~\eqref{eq:eq99}, and using the realized
wealth-share identity~\eqref{eq:wsh-realized} (their (92)) for $\Eone\dev{\wsh}_2$
on the left-hand side, KL obtain (their app.~pp.~76--77) the eq.~(99)-with-$\Gamp$
identity
\begin{multline}
\Bigl(\frac{\qk k}{\sav}-1\Bigr)\bigl(\rk_1-\rr_1\bigr)
+\frac{b_F}{\sav}\bigl(\fr{r}_1-\dev{\rer}_1-\rr_1\bigr)
-\disc\,\frac{\zs}{1+\zs}\,\frac{b^g_{H,s}}{\sav}\,\dev{\cy}_1
=\\
\frac{1}{\Gamp}(\rraF-\rra)(1-\cshare)\left[1+\lwedge(1-\disc)\frac{1}{\cshare+\frac{1}{\taylor}(1-\cshare)}(1-\cshare)\right]\tel^z\dev{\epsilon}^z_1\\
-\frac{1}{\Gamp}(\rraF-\rra)\lwedge(1-\disc)(1-\cshare)\frac{1}{\cshare+\frac{1}{\taylor}(1-\cshare)}\frac{1}{\taylor}\frac{1}{1+\zs}\dev{\cy}_1\\
+\frac{1}{\Gamp}\Bigl(\tfrac{1}{\zs}(\rraF-1)+(\rra-1)\Bigr)(1-\disc)\lwedge(1-\cshare)\frac{\zs}{1+\zs}\frac{1}{\cshare+\frac{1}{\taylor}(1-\cshare)}\frac{1}{\taylor}\dev{\cy}_1, \label{eq:eq99-prime}
\end{multline}
where the composite coefficient is now
\begin{equation}
\Gamp \;\equiv\; \Gam + \Bigl(\tfrac{1}{\zs}(\rraF-1)+(\rra-1)\Bigr)(1-\disc)\lwedge(1-\cshare)\frac{1}{\cshare+\frac{1}{\taylor}(1-\cshare)}\cshare \;>\;0. \label{eq:Gamma-prime-def}
\end{equation}
\gapflag{\eqref{eq:eq99-prime}--\eqref{eq:Gamma-prime-def} are KL's displayed
eq.~(99)-with-$\Gamp$ and the definition of $\Gamp$ (their app.~p.~77). The
reconstruction: substituting the relative-labor solution~\eqref{eq:lab-rel} into
the $(\fr{\dev{\ell}}_1-\dev{\ell}_1)$ terms of \eqref{eq:eq99} introduces a new
$\Eone\dev{\wsh}_2$ loading on the RHS (through the
$-\frac{1}{\cshare+\frac{1}{\taylor}(1-\cshare)}\frac{1+\zs}{\zs}\atil\Eone\dev{\wsh}_2$
piece of \eqref{eq:lab-rel}); moving that loading to the LHS augments the
$\Eone\dev{\wsh}_2$ coefficient $\Gam$ by exactly
$(\frac{1}{\zs}(\rraF-1)+(\rra-1))(1-\disc)\lwedge(1-\cshare)\frac{1}{\cshare+\frac{1}{\taylor}(1-\cshare)}\cshare$,
which is the added term in $\Gamp$. (The $\frac{1+\zs}{\zs}\atil$ from
\eqref{eq:lab-rel} cancels against the $\frac{\zs}{1+\zs}$ and one power of
$\atil$ in the relative-labor coefficient, leaving the $\cshare$ shown; the
positivity $\Gamp>0$ holds because, with $\rra,\rraF\ge1$ and $\lwedge\in(0,1)$,
the added term is non-negative and $\Gam>0$.) Substituting the symmetric-labor
solution~\eqref{eq:lab-sym} into the aggregate-employment term of \eqref{eq:eq99}
produces the $\dev{\epsilon}^z_1$ and $\dev{\cy}_1$ loadings shown on the RHS: the
productivity term collects the direct $(\rraF-\rra)(1-\cshare)$ loading plus the
wedge contribution $(\rraF-\rra)\lwedge(1-\disc)(1-\cshare)\times[\text{symmetric-labor coefficient on }\tel^z\dev{\epsilon}^z_1]$,
giving the displayed $[1+\lwedge(1-\disc)\frac{1-\cshare}{\cshare+\frac{1}{\taylor}(1-\cshare)}]$
bracket; the safety-shock terms collect the wedge and relative-labor
contributions through the $\dev{\cy}_1$ loadings of
\eqref{eq:lab-sym}--\eqref{eq:lab-rel}. We have matched KL's displayed eq.~(99)$'$
(app.~p.~77) line-by-line and flag the substitution as a genuine reconstruction:
KL display only ``Substituting these into the expression for international
portfolios [\dots] we obtain'' followed by the collected result. We reproduce the
collected result; the comparative statics below depend only on \emph{which} RHS
terms carry $\dev{\epsilon}^z_1$ versus $\dev{\cy}_1$ and on the sign of their
coefficients, which are robust to any residual algebra in the exact magnitudes.}

\subparagraph{Step 3: the symmetric-portfolio normalization and undetermined
coefficients.} The left-hand side of \eqref{eq:eq99-prime} is, by the realized
wealth-share identity~\eqref{eq:wsh-realized}, equal to $\Eone\dev{\wsh}_2$.
Solving \eqref{eq:eq99-prime} for $\Eone\dev{\wsh}_2$ and substituting the
realized returns of Proposition~\ref{prop:returns} (which, at the symmetric
portfolio, are the closed forms of Section~\ref{sec:engine}), KL express
$\Eone\dev{\wsh}_2$ as a linear combination of $\dev{\cy}_1$ and
$\dev{\epsilon}^z_1$ (their app.~p.~78):
\begin{multline}
\Eone\dev{\wsh}_2 =
\frac{1}{\Delta}\left[\frac{1}{\taylor}+\Bigl(1-\frac{1}{\taylor}\Bigr)\frac{\frac{1}{\taylor}(1-\cshare)}{\cshare+\frac{1}{\taylor}(1-\cshare)}\right]\frac{b_H}{\sav}\,\dev{\cy}_1
-\frac{1}{\Delta}\disc\,\frac{\zs}{1+\zs}\frac{b^g_{H,s}}{\sav}\,\dev{\cy}_1\\
+\frac{1}{\Delta}\Bigl(\frac{\qk k}{\sav}-1\Bigr)\left[(1-\cshare)+\Bigl(1-\frac{1}{\taylor}\Bigr)(1-\cshare)^2\frac{1}{\cshare+\frac{1}{\taylor}(1-\cshare)}\right]\tel^z\dev{\epsilon}^z_1, \label{eq:wsh-final}
\end{multline}
where
\begin{equation}
\Delta \;\equiv\; 1 + \left[\Bigl(\frac{\qk k}{\sav}-1\Bigr)+\frac{b_F}{\sav}\frac{\zs}{1+\zs}\right]
\frac{\frac{1}{\taylor}(1-\cshare)}{\cshare+\frac{1}{\taylor}(1-\cshare)}, \label{eq:Delta-def}
\end{equation}
which evaluates to $\Delta=1$ at the symmetric portfolio (where the bracketed
portfolio weights vanish). Here $b_H$ is Home's safe-dollar-bond position.
\gapflag{\eqref{eq:wsh-final}--\eqref{eq:Delta-def} are KL's displayed lines
(their app.~p.~78). The reconstruction: substituting the realized excess returns
$\rk_1-\rr_1$ and $\fr{r}_1-\dev{\rer}_1-\rr_1$ (Proposition~\ref{prop:returns},
both $\propto\dev{\cy}_1$ at the symmetric portfolio plus productivity loadings)
into the LHS of \eqref{eq:eq99-prime} makes the LHS itself a linear function of
$\Eone\dev{\wsh}_2$ (through the portfolio weights times the
$\Eone\dev{\wsh}_2$-dependence of the realized returns); collecting that
dependence onto the LHS produces the normalization factor $\Delta$ in
\eqref{eq:Delta-def}. At the symmetric portfolio $\frac{\qk k}{\sav}-1=0$ and
$\frac{b_F}{\sav}=0$, so $\Delta=1$ and the realized returns reduce to their
symmetric-portfolio forms. We transcribe KL's $\Delta$ and the
$\Eone\dev{\wsh}_2$ expression; the $\frac{b_H}{\sav}$ coefficient on the
$\dev{\cy}_1$ term comes from the leverage and Foreign-bond excess returns
(both $\propto\dev{\cy}_1$) collected through their portfolio weights, plus the
seigniorage term; the $(\frac{\qk k}{\sav}-1)$ coefficient on the
$\tel^z\dev{\epsilon}^z_1$ term comes from the productivity loading of
$\rk_1-\rr_1$. We flag that we reproduced KL's coefficients rather than
re-deriving each from the realized-return reduced forms.}

\medskip
We now read off the comparative statics by the method of undetermined
coefficients, evaluating at the symmetric portfolio ($\Delta=1$). Write the
equilibrium portfolio as the pair $(k,b_H)$ that makes \eqref{eq:wsh-final} hold
for \emph{both} independent shocks $\dev{\cy}_1$ and $\dev{\epsilon}^z_1$
simultaneously; equivalently, the coefficients on $\dev{\cy}_1$ and on
$\tel^z\dev{\epsilon}^z_1$ in \eqref{eq:wsh-final} must each equal the
corresponding coefficient required by the risk-sharing RHS of
\eqref{eq:eq99-prime}.

\emph{Comparative statics in $b^g_{H,s}$.} The supply of safe debt $b^g_{H,s}$
enters \eqref{eq:wsh-final} only through the seigniorage term
$-\frac{1}{\Delta}\disc\frac{\zs}{1+\zs}\frac{b^g_{H,s}}{\sav}\dev{\cy}_1$ and,
via the household budget, through $b_H$ (more government safe debt is held by
households). It does \emph{not} enter the $\tel^z\dev{\epsilon}^z_1$ coefficient,
which depends only on the capital weight $\frac{\qk k}{\sav}-1$. Matching the
$\dev{\epsilon}^z_1$ coefficient therefore leaves the equilibrium capital share
$k$ unchanged as $b^g_{H,s}$ varies:
\begin{equation}
\frac{dk}{db^g_{H,s}} = 0. \label{eq:cs-k-bg}
\end{equation}
Matching the $\dev{\cy}_1$ coefficient, an increase in $-b^g_{H,s}$ (more safe
debt issued) raises the seigniorage Home earns on impact of a positive safety
shock; to keep the risk-sharing condition satisfied at unchanged $k$, the
household must hold more dollar bonds, i.e.\ $b_H$ adjusts so that
\begin{equation}
\frac{db_H}{db^g_{H,s}} > 0. \label{eq:cs-bH-bg}
\end{equation}
\gapflag{Equations~\eqref{eq:cs-k-bg}--\eqref{eq:cs-bH-bg} are KL's displayed
comparative statics (app.~p.~78). The undetermined-coefficients argument:
\eqref{eq:wsh-final} must hold as an identity in the two independent shocks, so
the $\dev{\cy}_1$-coefficient equation and the $\dev{\epsilon}^z_1$-coefficient
equation are two scalar equations in the two unknown portfolio shares $(k,b_H)$.
Because $b^g_{H,s}$ appears \emph{only} in the $\dev{\cy}_1$ equation (the
seigniorage term), the $\dev{\epsilon}^z_1$ equation pins $k$ independently of
$b^g_{H,s}$, giving \eqref{eq:cs-k-bg}; differentiating the $\dev{\cy}_1$ equation
then gives the sign of $db_H/db^g_{H,s}$. We supply the
$2\times2$-block-triangular structure that makes $dk/db^g_{H,s}=0$ exact (the
$\dev{\epsilon}^z_1$ equation does not involve $b^g_{H,s}$); the sign of
$db_H/db^g_{H,s}$ follows from the positive seigniorage loading. KL state the two
results without displaying the linear-algebra block structure; we make it
explicit and flag this as a reconstruction of the undetermined-coefficients
step.}

\emph{Comparative statics in $\rraF/\rra$, holding $\rra+\frac{1}{\zs}\rraF$
fixed.} The combination $\rra+\frac{1}{\zs}\rraF$ enters $\Gam$ (and hence
$\Gamp$) in the denominator of every RHS term of \eqref{eq:eq99-prime}; holding it
fixed therefore freezes $\Gamp$. As KL note (their app.~p.~77), ``in comparative
statics with respect to $\rraF/\rra$ holding fixed $\rra+\frac{1}{\zs}\rraF$, on
the right-hand side it is clear that only the first two terms vary; the third,
capturing the effect of the relative labor response on international portfolios,
is constant.'' The first two RHS terms both carry the factor $(\rraF-\rra)$: the
productivity term $\frac{1}{\Gamp}(\rraF-\rra)(1-\cshare)[\cdots]\tel^z\dev{\epsilon}^z_1$
and the aggregate-employment safety term
$-\frac{1}{\Gamp}(\rraF-\rra)\lwedge(1-\disc)(1-\cshare)[\cdots]\dev{\cy}_1$. An
increase in $\rraF/\rra$ at fixed $\rra+\frac{1}{\zs}\rraF$ raises $\rraF-\rra$.
Matching the $\dev{\epsilon}^z_1$ coefficient, the productivity loading on the RHS
rises, so the equilibrium capital weight must rise to deliver a larger Home
wealth gain in good productivity states:
\begin{equation}
\frac{dk}{d[\rraF/\rra]}\Big|_{\rra+\frac{1}{\zs}\rraF} > 0. \label{eq:cs-k-gamma}
\end{equation}
Matching the $\dev{\cy}_1$ coefficient (with $\lwedge>0$), the safety-shock
loading on the RHS becomes more negative, so to satisfy risk sharing Home must
issue more dollar bonds (its dollar-bond share falls):
\begin{equation}
\frac{db_H}{d[\rraF/\rra]}\Big|_{\rra+\frac{1}{\zs}\rraF} < 0
\qquad(\text{assuming } \lwedge>0). \label{eq:cs-bH-gamma}
\end{equation}
\gapflag{Equations~\eqref{eq:cs-k-gamma}--\eqref{eq:cs-bH-gamma} are KL's
displayed comparative statics (app.~p.~78), with the second holding
$\rra+\frac{1}{\zs}\rraF$ fixed and assuming $\lwedge>0$. The
undetermined-coefficients argument mirrors the $b^g_{H,s}$ case: the
$\dev{\epsilon}^z_1$ equation pins $k$ and, because its RHS coefficient rises with
$\rraF-\rra$ (the productivity term), $dk/d[\rraF/\rra]>0$; the $\dev{\cy}_1$
equation then pins $b_H$, and because its RHS coefficient (the
aggregate-employment safety term, which carries $-(\rraF-\rra)\lwedge$) becomes
more negative with $\rraF-\rra$ when $\lwedge>0$, $db_H/d[\rraF/\rra]<0$. We
supply the matching argument and flag (i) the reliance on KL's own statement that
only the first two RHS terms vary at fixed $\rra+\frac{1}{\zs}\rraF$, and (ii) the
$\lwedge>0$ requirement for the sign of $db_H/d[\rraF/\rra]$, which KL state
explicitly (``assumes $\lwedge>0$''). The signs depend on the positivity of
$\Gamp$ \eqref{eq:Gamma-prime-def} and on the steady-state portfolio being levered
long capital and short dollar bonds, as at the symmetric benchmark.}

\medskip
This establishes parts (a) and (b) of Proposition~\ref{prop:portfolios}. \qed
(\,Proposition~\ref{prop:portfolios}\,)

\subsection{Discussion}

\paragraph{Two risk factors, two portfolio margins.} Proposition~4 organizes the
equilibrium portfolio around the two shocks the asset span must insure: a
\emph{global productivity} factor $\dev{\epsilon}^z_1$ and a \emph{safe-asset
demand} (safety) factor $\dev{\cy}_1$. The two portfolio margins --- the capital
(world-equity) share $k$ and the dollar-bond share $b_H$ --- map to the two
factors through the block structure exposed in the proof: the capital share is
pinned by the productivity-coefficient equation, the dollar-bond share by the
safety-coefficient equation.

\paragraph{The role of safe-debt supply.} As Home's government borrows more in
safe dollar bonds ($-b^g_{H,s}$ rises), Home receives more seigniorage on impact
of a positive safety shock: it is the issuer of the asset whose convenience value
spikes exactly when safety is demanded. This makes Home a \emph{natural insurer}
of the safety shock, and --- crucially --- it can provide that insurance
\emph{without loading up on productivity risk}. Concretely, Home borrows more in
dollar bonds (lower $b_H$) to hold more Foreign bonds, leaving its capital
(world-equity) exposure $k$ unchanged: this is exactly the content of
$dk/db^g_{H,s}=0$ together with $db_H/db^g_{H,s}>0$. The supply of safe dollar
debt is a margin for insuring the safety factor that is orthogonal to the
equity-risk margin.

\paragraph{The role of relative risk tolerance.} As Home becomes more
risk-tolerant than Foreign ($\rraF/\rra$ rises, holding
$\rra+\frac{1}{\zs}\rraF$ fixed so that the aggregate price of risk $\Gam$ is
held constant), Home optimally insures \emph{both} the negative-productivity and
the positive-safety shocks: it holds more capital (so it bears more of the global
equity risk, $dk/d[\rraF/\rra]>0$) and it borrows more in dollar bonds (so it
provides more safety insurance, $db_H/d[\rraF/\rra]<0$). KL emphasize that this is
the sense in which \emph{greater risk tolerance is necessary} to explain why the
United States takes a disproportionate exposure to equity returns: in their model,
the US equity-tilted external balance sheet requires the US to be more
risk-tolerant than the rest of the world.

\paragraph{The big\_cy reading.} Proposition~4 is precisely the KL tension that
big\_cy aims to revisit: KL need greater US risk tolerance ($\rra<\rraF$) to
rationalize the US equity tilt, yet high US risk tolerance is the very feature the
project regards as counterfactual. The two margins isolated here ---
$dk/d[\rraF/\rra]>0$ (equity tilt needs risk tolerance) and $db_H/db^g_{H,s}>0$
(safe-debt supply moves the dollar-bond margin without touching equity) ---
identify exactly where a large \emph{absolute} convenience yield might substitute
for the risk-tolerance asymmetry. We record KL's interpretation faithfully here
and defer the big\_cy extension to the dedicated extension memo.


% ---- 9. Proposition 5: sign of the risk premium / reserve-currency paradox ----
% =====================================================================
% 09_prop5.tex — Proposition 5: sign of the risk premium / reserve-currency
% paradox (WP-D)
% =====================================================================
\section{Proposition 5: the sign of the risk premium and the reserve-currency paradox}
\label{sec:prop5}

This section proves Proposition~5 of KL: the sign of the equilibrium risk premium
on Foreign bonds relative to safe dollar bonds. The risk premium is the covariance
of (the negative of) any agent's log pricing-kernel deviation with the excess log
return (eq.~\eqref{eq:rp-cov} of Section~\ref{sec:portfolio-engine}). The
proposition decomposes this covariance into a productivity component, signed by
$\rra-\rraF$, and a safety component, which is unambiguously positive --- the
result that safety shocks resolve the ``reserve-currency paradox'' of Maggiori
(2017).

\subsection{Statement}

\begin{proposition}[Sign of the risk premium; reserve-currency paradox]
\label{prop:riskpremium}
Consider the simplified economy of Proposition~\ref{prop:portfolios} (wages set
one period in advance, $\lwedge>0$, complete home bias), with the safety shock
$\dev{\cy}$ and the productivity shock $\dev{\epsilon}^z$ independent. Then, at
least around the symmetric-country-portfolio case, the equilibrium risk premium on
Foreign bonds relative to dollar bonds,
\begin{equation}
\Cov_t\bigl(-\dev{\pk}_{t,t+1},\;\fr{r}_{t+1}-\D\dev{\rer}_{t+1}-\rr_{t+1}\bigr), \label{eq:prop5-cov}
\end{equation}
satisfies:
\begin{enumerate}[label=(\alph*),leftmargin=1.8em,itemsep=3pt]
  \item if $\sigma^{\cy}=0$ (no safety shocks), the covariance~\eqref{eq:prop5-cov}
  is proportional to $\rra-\rraF$ (it is negative if Home is more risk-tolerant,
  $\rra<\rraF$);
  \item the covariance~\eqref{eq:prop5-cov} is \emph{rising} in $\sigma^{\cy}$
  (more volatile safety shocks raise the risk premium).
\end{enumerate}
The same statements hold using the Foreign household's pricing kernel
$\fr{\dev{\pk}}$ in place of $\dev{\pk}$.
\end{proposition}

\gapflag{We transcribe Proposition~5 to match KL's main-paper statement
(p.~1665) as relayed in the work-package brief; as for Proposition~4, the project
does not hold a local copy of the main-paper PDF, so the verbatim wording is taken
from the brief. The mathematical content --- the $(\sigma^z)^2$ coefficient
signed by $\rra-\rraF$ and the positive $(\sigma^{\cy})^2$ coefficient --- matches
the decomposition displayed at the end of the appendix proof (their app.~p.~78),
which we hold directly. Here $\sigma^{\cy}$ denotes the standard deviation of the
safety shock (KL's $\sigma^\omega$) and $\sigma^z$ that of the productivity shock;
$\D\dev{\rer}_{t+1}$ is the real-exchange-rate change converting the Foreign-bond
return to Home units. We write the statement at a generic date $t$ following KL;
in the period-$1$ engine $t=0$, so the relevant object is the period-$0$
conditional covariance, evaluated below at $\Cov_0$.}

\subsection{Proof}

\paragraph{Strategy.} By local completeness~\eqref{eq:risk-sharing}, the risk
premium equals
$-\Ezero[(\dev{\pk}_{0,1}-\rra\tel^z\dev{\epsilon}^z_1)(\fr{r}_1-\D\dev{\rer}_1-\rr_1)]$
(the directly-priced productivity loading $\rra\tel^z\dev{\epsilon}^z_1$ is the
part of the kernel that is spanned and does not contribute a premium beyond the
covariance; KL form the premium with the risk-adjusted kernel). We expand this
covariance using the B.2 results, exploit the independence of the two shocks to
split it into a $(\sigma^z)^2$ piece and a $(\sigma^{\cy})^2$ piece, and sign each.

\subparagraph{Step 1: the covariance object.} The risk premium is (their
app.~p.~78)
\begin{equation}
\text{RP} \;=\; -\Ezero\Bigl[\bigl(\dev{\pk}_{0,1}-\rra\tel^z\dev{\epsilon}^z_1\bigr)\bigl(\fr{r}_1-\D\dev{\rer}_1-\rr_1\bigr)\Bigr]. \label{eq:prop5-RP}
\end{equation}
Both factors are mean-zero to first order (period-$0$ conditional), so
$\text{RP}=-\Cov_0(\dev{\pk}_{0,1}-\rra\tel^z\dev{\epsilon}^z_1,\;\fr{r}_1-\D\dev{\rer}_1-\rr_1)$,
which has the \emph{same sign structure} as~\eqref{eq:prop5-cov}: the two differ
only by $-\Cov_0(\rra\tel^z\dev{\epsilon}^z_1,\fr{r}_1-\D\dev{\rer}_1-\rr_1)$, a
pure $(\sigma^z)^2$ term that shifts the magnitude of the $(\sigma^z)^2$ piece but
not its $\rra-\rraF$ sign and leaves the $(\sigma^{\cy})^2$ piece untouched (this
is why KL state Proposition~5 with the plain kernel $-\dev{\pk}$ while proving it
with the risk-adjusted kernel $\dev{\pk}_{0,1}-\rra\tel^z\dev{\epsilon}^z_1$). The
risk-adjusted kernel is \eqref{eq:m01-home} and the excess return
$\fr{r}_1-\D\dev{\rer}_1-\rr_1$ is the
Proposition~\ref{prop:returns}(c)-minus-(a) realized excess return.

\subparagraph{Step 2: load each factor on the two shocks.} At the symmetric
portfolio, both factors in \eqref{eq:prop5-RP} are linear in the two independent
innovations $(\dev{\epsilon}^z_1,\dev{\cy}_1)$. From the
Proposition~\ref{prop:returns} realized returns (specialized to the symmetric
portfolio, complete home bias, predetermined wages):
\begin{align}
\fr{r}_1-\D\dev{\rer}_1-\rr_1
&= \underbrace{A_z\,\tel^z\dev{\epsilon}^z_1}_{\text{productivity}}
+\underbrace{A_\omega\,\dev{\cy}_1}_{\text{safety}}, \label{eq:exret-decomp}
\end{align}
with $A_\omega<0$: from Proposition~\ref{prop:returns}(a),(c),
$\fr{r}_1-\D\dev{\rer}_1-\rr_1=(1-\cshare)\dev{\ell}_1-\rr_1$, and with
$\dev{\ell}_1\propto-\dev{\cy}_1$ and $\rr_1\propto+\dev{\cy}_1$ the safety
loading is negative (Foreign bonds underperform dollar bonds on impact of a
positive safety shock). The productivity loading $A_z$ is the response of the same
excess return to a productivity surprise.

Likewise the risk-adjusted kernel~\eqref{eq:m01-home} loads on the two shocks as
\begin{equation}
\dev{\pk}_{0,1}-\rra\tel^z\dev{\epsilon}^z_1
= B_z\,\tel^z\dev{\epsilon}^z_1 + B_\omega\,\dev{\cy}_1, \label{eq:kernel-decomp}
\end{equation}
where $B_z$ collects the kernel's productivity loading (through
$\Eone\dev{\wsh}_2$, aggregate employment, and the explicit
$-\rra(1-\cshare)\tel^z\dev{\epsilon}^z_1$ term) and $B_\omega$ its safety loading
(through the safety-driven contraction in employment and the wealth share).

\subparagraph{Step 3: split the covariance by independence.} Because
$\dev{\epsilon}^z_1\perp\dev{\cy}_1$, the cross terms vanish in expectation
($\Ezero[\dev{\epsilon}^z_1\dev{\cy}_1]=0$), so \eqref{eq:prop5-RP} becomes
\begin{equation}
\text{RP} = -\Bigl[B_z A_z\,(\tel^z)^2\Ezero(\dev{\epsilon}^z_1)^2
+ B_\omega A_\omega\,\Ezero(\dev{\cy}_1)^2\Bigr]
= -B_z A_z\,(\sigma^z)^2 - B_\omega A_\omega\,(\sigma^{\cy})^2, \label{eq:RP-split}
\end{equation}
absorbing $(\tel^z)^2\Ezero(\dev{\epsilon}^z_1)^2$ into $(\sigma^z)^2$ and
$\Ezero(\dev{\cy}_1)^2$ into $(\sigma^{\cy})^2$. Thus the risk premium is a linear
combination of $(\sigma^z)^2$ and $(\sigma^{\cy})^2$, as claimed.

\subparagraph{Step 4: sign the productivity coefficient.} The coefficient on
$(\sigma^z)^2$ is $-B_z A_z$. From the B.2 results, the productivity loading of
the risk-adjusted kernel carries the cross-country risk-aversion difference: a
positive productivity shock raises Home consumption and the Home wealth share, and
the kernel's response is governed by $\rra$ versus $\rraF$ through the
risk-sharing condition~\eqref{eq:risk-sharing}. KL show (their app.~p.~78) that
the resulting coefficient on $(\sigma^z)^2$ ``takes the sign of $\rra-\rraF$.''
\begin{equation}
-B_z A_z \;\propto\; \rra-\rraF. \label{eq:coef-z}
\end{equation}
\gapflag{Equation~\eqref{eq:coef-z} is KL's stated sign result (app.~p.~78: ``the
coefficient on the former takes the sign of $\rra-\rraF$''). The reconstruction:
the excess return $A_z$ (the productivity loading of
$\fr{r}_1-\D\dev{\rer}_1-\rr_1$) is common across agents, while the kernel loading
$B_z$ differs across Home and Foreign exactly through the risk-aversion difference
--- the risk-sharing condition~\eqref{eq:risk-sharing} forces
$\dev{\pk}_{0,1}-\rra\tel^z\dev{\epsilon}^z_1=\fr{\dev{\pk}}_{0,1}-\rraF\tel^z\dev{\epsilon}^z_1+\D\dev{\rer}_1$,
so the productivity loadings of the two risk-adjusted kernels differ, and the
covariance of the \emph{Home} risk-adjusted kernel with the common excess return
inherits the sign of $\rra-\rraF$. (Intuitively: when $\rra<\rraF$, Home is more
risk-tolerant, so in good productivity states Home's marginal utility falls by
less than Foreign's relative to the spanned benchmark; the Home kernel covaries
\emph{positively} with the dollar-appreciating excess return, making $-B_z A_z<0$,
i.e.\ a \emph{negative} productivity contribution to the Foreign-bond premium.)
We supply the mechanism; KL display only the sign. The sign $\rra-\rraF$ (rather
than $\rraF-\rra$) is KL's, and is consistent with the Proposition~4 reading that
a more risk-tolerant Home holds the dollar (the appreciating asset) as the
negative-beta leg. We flag that we did not compute the magnitude of $-B_z A_z$,
only reproduce KL's sign.}

\subparagraph{Step 5: sign the safety coefficient.} The coefficient on
$(\sigma^{\cy})^2$ is $-B_\omega A_\omega$. We have $A_\omega<0$ (Step~2:
Foreign bonds underperform dollar bonds on impact of a positive safety shock). The
kernel's safety loading $B_\omega$ is also negative: a positive safety shock
contracts Home employment and lowers the Home wealth share / consumption, raising
the (risk-adjusted) marginal-utility kernel $\dev{\pk}_{0,1}$ --- but in the
covariance with the \emph{negative} of the kernel the relevant sign works out so
that $-B_\omega A_\omega>0$. KL state (their app.~p.~78) that ``the coefficient on
the latter is positive,'' independent of the risk-aversion difference:
\begin{equation}
-B_\omega A_\omega \;>\;0. \label{eq:coef-omega}
\end{equation}
\gapflag{Equation~\eqref{eq:coef-omega} is KL's stated sign result (app.~p.~78:
``the coefficient on the latter is positive''). The reconstruction: the safety
shock moves both the kernel and the excess return in the \emph{same} direction ---
a positive $\dev{\cy}_1$ depresses the Foreign-bond excess return ($A_\omega<0$)
\emph{and} (because it contracts global employment and Home wealth) makes
state-prices high, so $-\dev{\pk}_{0,1}$ covaries positively with the
underperformance. The product $-B_\omega A_\omega$ is therefore positive
\emph{regardless of $\rra$ versus $\rraF$}, because the safety shock's propagation
(Propositions~\ref{prop:prices}--\ref{prop:returns}) does not rely on a
risk-aversion difference --- it operates through the convenience-yield wedge and
nominal rigidity alone. This is the crux of KL's point that safety shocks raise
the risk premium even absent differences in risk tolerance. We supply the
sign-chasing; KL display only ``positive.'' We flag that we reproduce KL's sign
via the comovement argument rather than computing $B_\omega$ and $A_\omega$
explicitly; the comovement (both negative on impact, product positive in the
$-\Cov$ premium) is exactly the Proposition~\ref{prop:returns} structure.}

\subparagraph{Step 6: conclude.} Combining \eqref{eq:RP-split},
\eqref{eq:coef-z}, and \eqref{eq:coef-omega}: if $\sigma^{\cy}=0$ then
$\text{RP}=-B_z A_z(\sigma^z)^2\propto\rra-\rraF$, establishing part~(a); and
$\partial\text{RP}/\partial(\sigma^{\cy})^2=-B_\omega A_\omega>0$, so RP is rising
in $\sigma^{\cy}$, establishing part~(b). By the local-completeness
condition~\eqref{eq:risk-sharing}, the Foreign risk-adjusted kernel
$\fr{\dev{\pk}}_{0,1}-\rraF\tel^z\dev{\epsilon}^z_1$ differs from the Home one only
by $\D\dev{\rer}_1$, which is itself a linear combination of the same two shocks;
the covariance with the excess return is therefore identical, so the statements
hold for the Foreign household's kernel as well. \qed
(\,Proposition~\ref{prop:riskpremium}\,)

\subsection{Discussion}

\paragraph{Part 1: the reserve-currency paradox.} Absent safety shocks
($\sigma^{\cy}=0$), part~(a) says Foreign bonds earn a \emph{negative} risk
premium relative to dollar bonds whenever Home is more risk-tolerant
($\rra<\rraF$). The mechanism: with only productivity risk, the dollar appreciates
in \emph{good} times. A positive productivity shock raises world output; the more
risk-tolerant Home, levered long in capital (Proposition~\ref{prop:portfolios}),
sees its wealth and consumption rise; Home demand shifts toward Home goods and the
dollar appreciates exactly in this good state. The dollar (safe-bond leg) thus
pays well in good times and badly in bad times --- it is a \emph{high}-beta,
positive-premium asset, so Foreign bonds carry a \emph{negative} premium. This is
the ``reserve-currency paradox'' of Maggiori (2017): the would-be reserve currency
counterfactually appreciates in good states. KL's contribution here is to show the
paradox is \emph{robust} to endogenous production and nominal rigidities --- it
survives in their richer environment, not just in an endowment model.

\paragraph{Part 2: safety shocks resolve the paradox.} Part~(b) says that adding
safety shocks raises the Foreign-bond risk premium, and sufficiently volatile
safety shocks flip its sign to \emph{positive}. The mechanism is the mirror image
of part~1: a positive safety shock makes the dollar appreciate in \emph{bad} times
(high safe-asset demand coincides with a global employment contraction,
Propositions~\ref{prop:prices}--\ref{prop:returns}). Now the dollar pays well
exactly when marginal utility is high --- it is a \emph{negative}-beta, hedge
asset --- so Foreign bonds must earn a \emph{positive} premium to be held. KL
emphasize the key feature: safety shocks raise the risk premium \emph{even absent
any difference in risk tolerance}, because their propagation in
Propositions~\ref{prop:prices}--\ref{prop:returns} does not rely on
$\rra\neq\rraF$. The safety channel and the risk-tolerance channel are distinct;
the former alone can deliver the dollar's negative beta.

\paragraph{Connection to the three target facts.} Proposition~5 ties together the
paper's empirical program. (i) The \emph{dollar's negative beta}: with sufficiently
volatile safety shocks, the dollar appreciates in bad global states, giving safe
dollar bonds a negative beta and a negative risk premium --- the structural
counterpart of the dollar's hedge property. (ii) The \emph{US as the world's
insurer}: the positive Foreign-bond premium is the compensation Home demands for
providing dollar safety, the asset-pricing face of the wealth-transfer mechanism
of Proposition~\ref{prop:wealth}. (iii) The \emph{resolution of the
reserve-currency paradox}: safety shocks overturn the Maggiori (2017) sign,
reconciling the model with the observed positive carry/term premia on
foreign-currency bonds without requiring an implausible risk-tolerance asymmetry.

\paragraph{The big\_cy reading.} Proposition~5 is the asset-pricing hinge of the
big\_cy thesis. KL obtain the dollar's negative premium through two channels ---
risk-tolerance asymmetry (part~1, sign $\rra-\rraF$) and safety-shock volatility
(part~2, the positive $(\sigma^{\cy})^2$ coefficient). The second channel
\emph{does not require} $\rra\neq\rraF$. A \emph{large absolute} convenience yield
raises the effective volatility and propagation of the safety factor, magnifying
the positive $(\sigma^{\cy})^2$ contribution to the Foreign-bond premium; the
project's question is whether this channel, scaled to DiTella-magnitude
convenience yields, can deliver the dollar's negative beta on its own ---
substituting for the risk-tolerance asymmetry that Proposition~4 otherwise
requires. We record KL's result faithfully here and defer the quantitative
assessment to the big\_cy extension.


% ---- 10. efficient risk sharing / LOOP + specialness of safe dollar bonds ----
% =====================================================================
% 10_risksharing_specialness.tex — efficient risk sharing / LOOP,
% the specialness of safe dollar bonds (main paper II.D / II.C),
% Foreign-only demand (App. B.4), distinct capital + sticky prices
% (App. B.5) (WP-E)
% =====================================================================
\section{Efficient risk sharing, the law of one price, and the specialness of safe dollar bonds}
\label{sec:risksharing-specialness}

The previous sections built the analytical machine: the period-$1$ engine
(Section~\ref{sec:engine}), Propositions~\ref{prop:prices}--\ref{prop:wealth} on
prices/returns/wealth, the second-order portfolio engine
(Section~\ref{sec:portfolio-engine}), and Propositions~\ref{prop:portfolios}--%
\ref{prop:riskpremium} on portfolios and the risk premium. This section collects
the three remaining analytical statements of KL's Section~II and the two
robustness exercises of their online Appendix~B. First, we make explicit the
\emph{efficient-risk-sharing / law-of-one-price} content of the model: the
result, stated cleanly in II.C's ``\emph{Efficient Risk Sharing and the Law of
One Price}'' paragraph (their p.~1667), that the dollar appreciation in bad times
is a \emph{necessary consequence} of efficient risk sharing plus the law of one
price, manifested as a tight link between the real exchange rate and relative
output. Second, we transcribe and explain KL's three ``specialness of safe dollar
bonds'' arguments (their II.D, p.~1666). Third, we reconstruct the
Foreign-only-demand robustness (their App.~B.4) and its multiplier result
$\mu_t=\cy_t$. Fourth, we summarize compactly the distinct-capital-stocks-and-%
sticky-prices environment (their App.~B.5), stating the analogs of
Propositions~1--5 and flagging that KL exclude their proofs.

\subsection{The efficient-risk-sharing / law-of-one-price relation}
\label{sec:loop}

\paragraph{From the kernel relation to a relation in quantities.} The
local-completeness condition~\eqref{eq:risk-sharing} of
Section~\ref{sec:local-completeness} is the Backus--Smith / law-of-one-price
relation in \emph{kernel} form,
\begin{equation}
\dev{\pk}_{0,1} - \rra\,\tel^z\dev{\epsilon}^z_1
= \fr{\dev{\pk}}_{0,1} - \rraF\,\tel^z\dev{\epsilon}^z_1 + \D\dev{\rer}_1.
\tag{\ref{eq:risk-sharing} revisited}
\end{equation}
KL observe (their p.~1667) that this relation has a sharp counterpart in
\emph{quantities} once one specializes to the unitary-risk-aversion benchmark
$\rra=\rraF=1$, which --- recalling the maintained unit IES of the simplified
environment --- makes preferences time-separable. In that benchmark the kernel is
the pure consumption-growth term, $\dev{\pk}_{0,1}=-\dev{\rs c}_1$ and
$\fr{\dev{\pk}}_{0,1}=-\fr{\dev{\rs c}}_1$ (the continuation-value term in
\eqref{eq:m01} carries the factor $1-\rra=0$), and the productivity-loading terms
$\rra\tel^z\dev{\epsilon}^z_1$, $\rraF\tel^z\dev{\epsilon}^z_1$ are equal. The
risk-sharing relation~\eqref{eq:risk-sharing} then reduces to the
\emph{Backus--Smith relation in levels} (their displayed equation, p.~1667,
``assuming no shocks prior to period $t$''):
\begin{equation}
\fr{\dev{\rs c}}_1 - \dev{\rs c}_1 = \dev{\rer}_1, \label{eq:bs-levels}
\end{equation}
i.e.\ Home consumption rises relative to Foreign exactly when the Home real
exchange rate depreciates ($\dev{\rer}_1<0$ raises $\dev{\rs c}_1$ relative to
$\fr{\dev{\rs c}}_1$). This is the standard complete-markets prediction that
relative consumption tracks the real exchange rate.
\gapflag{KL display the benchmark relation~\eqref{eq:bs-levels} as
``$\hat c^{*}_t-\hat c_t=\hat q_t$'' for the $\rra=\rraF=1$ case (their p.~1667).
We supply the one-line provenance: it is the local-completeness
condition~\eqref{eq:risk-sharing} evaluated at $\rra=\rraF=1$, where the
continuation-value term of the kernel~\eqref{eq:m01} drops (its coefficient
$1-\rra$ vanishes) and the symmetric productivity loadings cancel, leaving
$-\dev{\rs c}_1=-\fr{\dev{\rs c}}_1+\dev{\rer}_1$, i.e.\
$\fr{\dev{\rs c}}_1-\dev{\rs c}_1=\dev{\rer}_1$. KL's sign convention writes
$\hat q$ as a Home appreciation (Section~\ref{sec:aux-vars},
eq.~\eqref{eq:rer-def}); their displayed line $\hat c^{*}_t-\hat c_t=\hat q_t$ is
exactly~\eqref{eq:bs-levels} in our notation.}

\paragraph{Linking the real exchange rate to relative output.} KL then combine
\eqref{eq:bs-levels} with the engine's relative-goods-market and
intratemporal-allocation relations to obtain a relation between the real exchange
rate and \emph{relative output}. The relevant engine inputs are the
real-exchange-rate--terms-of-trade relation~\eqref{eq:q1-s1},
$\dev{\rer}_1=\hbias\frac{1+\zs}{\zs}\dev{\tot}_1$, and the cross-country
goods-market / intratemporal condition. KL's displayed combination (their
p.~1667, with an \emph{arbitrary} degree of home bias $\hbias$ retained ``to
emphasize the generality of the result'') reads
\begin{equation}
\hbias\,\frac{1+\zs}{\zs}\bigl(\fr{\dev{\rs c}}_1-\dev{\rs c}_1\bigr)
+\Bigl(1+\hbias\,\tfrac{1+\zs}{\zs}\Bigr)
\Bigl(1-\hbias\,\tfrac{1+\zs}{\zs}\Bigr)
\frac{1}{\hbias\,\frac{1+\zs}{\zs}}\,\tel\,\dev{\rer}_1
= \fr{\out}_1 - \out_1, \label{eq:loop-rel-raw}
\end{equation}
where $\out_1,\fr{\out}_1$ are Home and Foreign log output deviations. Using the
Backus--Smith relation~\eqref{eq:bs-levels} to replace
$\fr{\dev{\rs c}}_1-\dev{\rs c}_1=\dev{\rer}_1$ in the first term, the two
left-hand terms collapse into a single multiple of $\dev{\rer}_1$, and solving
for $\dev{\rer}_1$ yields KL's displayed law-of-one-price relation (their
p.~1667):
\begin{equation}
\boxed{\;
\dev{\rer}_1 = \frac{1}
{\hbias\,\frac{1+\zs}{\zs}
+\Bigl(1+\hbias\,\frac{1+\zs}{\zs}\Bigr)\Bigl(1-\hbias\,\frac{1+\zs}{\zs}\Bigr)
\dfrac{1}{\hbias\,\frac{1+\zs}{\zs}}\,\tel}\,
\bigl(\fr{\out}_1 - \out_1\bigr).\;} \label{eq:loop-solved}
\end{equation}
\gapflag{Equations~\eqref{eq:loop-rel-raw}--\eqref{eq:loop-solved} are KL's two
displayed lines on p.~1667 (the relation in
$\varsigma\frac{1+\zeta^{*}}{\zeta^{*}}(\hat c^{*}_t-\hat c_t)$ and $\sigma\hat q_t$
equal to $\hat y^{*}_t-\hat y_t$, and the solved expression for $\hat q_t$). We
transcribe them in the document's macros (KL's
$\varsigma\frac{1+\zeta^{*}}{\zeta^{*}}$ is our $\hbias\frac{1+\zs}{\zs}$, the
recurring relative-price loading of~\eqref{eq:q1-s1}; their $\sigma$ is our
$\tel$). The first line is the log-linearized statement of ``intratemporal
optimality between Home- and Foreign-produced goods in each country, together with
goods market clearing''; it is the same content as the engine's cross-country
goods-market relation~\eqref{eq:rc-rel}, re-expressed in output rather than
terms-of-trade units and with the relative-consumption term written using
$\hbias\frac{1+\zs}{\zs}(\fr{\dev{\rs c}}_1-\dev{\rs c}_1)$. We did not re-derive
the output-form of~\eqref{eq:rc-rel} from primitives; we transcribe KL's
displayed relation and verify that substituting~\eqref{eq:bs-levels} collapses
it to~\eqref{eq:loop-solved}.}

\paragraph{Reading the result.} The denominator of~\eqref{eq:loop-solved} is
strictly positive whenever there is consumption home bias, $\hbias>0$ (so that
$0<\hbias\frac{1+\zs}{\zs}<1$ and the middle factor
$(1-\hbias\frac{1+\zs}{\zs})>0$). Hence~\eqref{eq:loop-solved} delivers KL's
central qualitative statement (their p.~1667):
\begin{quote}
\emph{Given consumption home bias $\hbias>0$, any shock that appreciates the Home
real exchange rate ($\dev{\rer}_1>0$) must be reflected in a relative fall in
Home output ($\out_1<\fr{\out}_1$).}
\end{quote}
KL identify the two distinct channels behind this link, both visible
in~\eqref{eq:loop-rel-raw}:
\begin{enumerate}[leftmargin=1.8em,itemsep=3pt]
  \item \textbf{Efficient risk sharing.} The Backus--Smith
  relation~\eqref{eq:bs-levels} forces relative Home \emph{consumption} to fall
  when the Home real exchange rate appreciates. This is the
  $\hbias\frac{1+\zs}{\zs}(\fr{\dev{\rs c}}_1-\dev{\rs c}_1)$ term: efficient
  risk sharing calls for Home consumption to fall relative to Foreign exactly in
  states of Home appreciation.
  \item \textbf{The law of one price (expenditure switching).} Under
  producer-currency pricing (the law of one price), a Home appreciation raises the
  relative price of Home goods --- the associated improvement in Home's terms of
  trade triggers a global expenditure switch \emph{away} from Home-produced goods,
  depressing Home output. This is the $\tel\,\dev{\rer}_1$ term, with the trade
  elasticity $\tel$ governing its strength.
\end{enumerate}
Both channels push Home output down when the dollar appreciates. KL note (their
p.~1667) that relaxing the law of one price --- e.g.\ under sticky local-currency
pricing --- can sever the \emph{second} channel by decoupling the consumer-facing
relative price from the real exchange rate, but the \emph{first} (efficient
risk-sharing) channel survives as long as risk sharing is efficient. This is why
the result is ``a natural consequence of efficient risk sharing and the law of
one price'' rather than of either alone.

\paragraph{Relation to the literature.} KL place this result against three
alternatives (their p.~1667). \emph{Sauzet (2023)} shows that a fall in relative
US output in crises can resolve the reserve-currency paradox following exactly
this logic; KL's contribution is to \emph{identify the shock} --- the flight to
safe dollar assets --- that produces this output pattern endogenously.
\emph{Maggiori (2017)} resolves the paradox instead through an expenditure shift
toward US goods in crises, which operates through the goods-market term
of~\eqref{eq:loop-rel-raw}, allowing the dollar to appreciate even with relative
output unchanged. \emph{Dahlquist et al.\ (2023)} generate a dollar appreciation
from a fall in US output as in Sauzet but augment with deep habits, so the output
decline also shifts relative demand toward US goods. KL's safety-shock mechanism
delivers the output channel without an exogenous taste shift.

\paragraph{The unitary-risk-aversion benchmark and constant wealth share.} The
benchmark $\rra=\rraF=1$ used to derive~\eqref{eq:bs-levels} is the same benchmark
referenced in Propositions~\ref{prop:wealth}--\ref{prop:portfolios}. KL note
(their footnote~58, App.~B.5) that
\begin{quote}
\emph{consistent with eq.~(99) in the baseline simplified environment, unitary
risk aversions $\rra=\rraF=1$ imply that equilibrium financial portfolios hedge
only labor-income risk,}
\end{quote}
and, correspondingly, that efficient risk sharing keeps Home's \emph{overall}
wealth share (inclusive of labor income) constant in response to all shocks. The
analytical content is visible in the international-portfolio
equation~\eqref{eq:eq99}: every right-hand-side term carries either
$(\rraF-\rra)$ or $\bigl(\frac{1}{\zs}(\rraF-1)+(\rra-1)\bigr)$, both of which
vanish at $\rra=\rraF=1$, so the risk-sharing requirement on $\Eone\dev{\wsh}_2$
collapses to the labor-income-hedging terms alone. This is the benchmark in which
``Home's overall wealth share is constant,'' against which Proposition~4's
comparative statics in $\rraF/\rra$ are taken.
\gapflag{The connection between the $\rra=\rraF=1$ benchmark and the constant
overall wealth share is KL's (their footnote~58 and the discussion at p.~1668:
``in the benchmark with identical risk tolerance across countries, efficient risk
sharing calls for the Home wealth share to increase upon a flight to safety
\emph{provided risk aversion exceeds unity}''). We reconstruct the analytical
reason --- that the $(\rraF-\rra)$ and $\bigl(\frac{1}{\zs}(\rraF-1)+(\rra-1)\bigr)$
loadings of~\eqref{eq:eq99} both vanish at $\rra=\rraF=1$, leaving only the
labor-income-hedging terms, so the overall wealth share (financial plus human) is
constant. Note the distinction KL draw: at $\rra=\rraF=1$ the \emph{overall}
wealth share is constant, whereas for $\rra=\rraF>1$ the overall wealth share
\emph{rises} on a flight to safety (the real-exchange-rate-hedging motive of
relative-supply shocks); the $\rra=\rraF=1$ knife-edge is the one in which the
financial portfolio hedges labor income exactly. We flag that we supply this
reading from the structure of~\eqref{eq:eq99} rather than re-deriving the overall
wealth share from a separate human-wealth accounting.}

\subsection{The specialness of safe dollar bonds}
\label{sec:specialness}

KL devote their Section~II.D (p.~1666) to three dimensions in which the demand
for safe dollar bonds is \emph{special} relative to the demand for other assets.
These are stated as arguments / extensions rather than as formal propositions; we
transcribe each with the available analytical content and match KL's level of
detail (expository, with the formal verification deferred to the cited
extensions).

\subsubsection{Zero net supply}
\label{sec:zero-net-supply}

Safe dollar bonds are in \emph{zero net supply}: the outstanding stock is the debt
issued by the Home government (their fiscal rule~\eqref{eq:fiscal-supply}), held
in positive amounts by private agents, so that private and public positions sum to
zero. KL's argument (their p.~1666) is that this is precisely what makes a flight
to safety \emph{contractionary}:
\begin{quote}
\emph{The demand for safe dollar bonds triggers a Keynesian recession because
these assets are in zero net supply.}
\end{quote}
The contrast is with a counterfactual in which the heightened demand were instead
for \emph{capital} (equity). Formally, suppose capital also entered the utility
function and its nonpecuniary value rose on the margin. Then --- because capital
is in \emph{positive} net supply and is a produced, priced asset --- the increase
in its nonpecuniary value would raise the price of capital $\qk_t$, raise
household wealth (existing capital holders are richer), and thereby \emph{raise}
consumption demand and aggregate output. A demand shift toward capital is
\emph{expansionary}. KL note further that relaxing the assumption of a fixed
global capital stock would reinforce this: the demand for capital would also
stimulate output through increased \emph{investment}.

The economic content is the bond-vs-equity distinction. A safety shock that raises
demand for an asset in \emph{zero} net supply cannot be accommodated by a quantity
expansion (there is no net stock to bid up in aggregate); the only margin is the
price of \emph{other} assets and, through the convenience-yield
wedge~\eqref{eq:euler-r} and nominal rigidity, the real interest rate and
employment. The flight to the safe bond therefore manifests as the inability of
the real rate to clear the goods market at full employment --- a Keynesian
recession --- rather than as a wealth-and-investment expansion. KL frame this as
underscoring the importance of distinguishing convenience yields in \emph{bond}
versus \emph{equity} markets, as in Koijen and Yogo (2020), for understanding
their macroeconomic effects.
\gapflag{The zero-net-supply argument is stated by KL as prose (their p.~1666);
they do not write out the capital-in-utility counterfactual formally. We supply
the mechanism --- that a positive-net-supply, priced asset (capital) absorbs a
demand shift through its price $\qk_t$ and a wealth/investment expansion, whereas
a zero-net-supply asset (the safe bond) can only transmit the shift through the
convenience-yield wedge and nominal rigidity, producing the Keynesian recession of
Proposition~\ref{prop:prices}. This is consistent with, but not a re-derivation
of, KL's argument; the formal capital-in-utility extension is not carried out by
KL and we do not carry it out here. We flag that the ``stimulating'' sign of a
capital-demand shift is asserted, following KL's appeal to Koijen and Yogo (2020).}

\subsubsection{Country size}
\label{sec:country-size}

The second dimension is the \emph{large size} of the US economy. KL argue (their
p.~1666) that the demand for safe dollar bonds is special relative to the demand
for the bonds of \emph{small} issuers. The thought experiment is to augment the
two-country model with an \emph{infinitesimal} third country --- for concreteness,
Switzerland --- that is a negligible part of the global economy. An increase in
the portfolio demand for Swiss-franc bonds would:
\begin{itemize}[leftmargin=1.8em,itemsep=2pt]
  \item generate an \emph{appreciation} of the franc;
  \item cause a \emph{decline} in Swiss production (the same
  efficient-risk-sharing-plus-LOOP logic of Subsection~\ref{sec:loop}: an
  appreciation of the small country's real exchange rate is reflected in a
  relative fall in its output); but
  \item have \emph{negligible spillovers} on global demand and production, because
  the country's share of the world economy is infinitesimal.
\end{itemize}
By contrast, the US is large enough that a flight to safe dollar bonds moves
\emph{global} aggregates --- world employment and output --- not merely US relative
output. KL connect this to the message of Hassan (2013) that the United States'
relative size in the global economy may be an important contributor to the
dollar's negative beta: it is precisely because the dollar issuer is large that
the dollar's appreciation in bad times has aggregate, rather than purely relative,
consequences and hence commands a sizeable (negative-beta) risk premium.
\gapflag{The country-size argument is stated by KL as a thought experiment (their
p.~1666); the infinitesimal-third-country model is described but not formally
written out. We state it exactly as KL do, noting that the ``appreciation $+$
relative-output decline $+$ negligible global spillover'' conclusion for the small
country follows from applying the law-of-one-price / efficient-risk-sharing
relation~\eqref{eq:loop-solved} to a country of vanishing global weight. We do not
formalize the three-country extension; KL do not, and we match their expository
level. The link to Hassan (2013) is KL's.}

\subsubsection{Currency of invoicing}
\label{sec:invoicing}

The third dimension is the widespread use of the dollar in \emph{pricing}. KL
argue (their p.~1666) that dollar invoicing renders the demand for safe dollar
bonds special relative to the currencies of other large countries. The concrete
example: suppose nominal wages \emph{even in Foreign} are denominated and sticky in
dollars (rather than in Foreign currency). Then the global decline in employment
following a positive safety shock is \emph{exacerbated} relative to the baseline
model. The intuition is that dollar deflation following a positive safety shock
(the safety shock is deflationary in the dollar bloc,
Proposition~\ref{prop:prices}) now implies that product wages \emph{in both
countries} are too high --- the sticky dollar wage does not adjust down with the
Foreign price level --- generating a more severe and more uniform global recession.

KL note that similar results obtain under dollar pricing of \emph{exports} (the
dominant-currency-pricing paradigm of Gopinath (2015), Gopinath et al.\ (2020),
and Mukhin (2022)): dollar invoicing of trade among countries not involving the
United States is analogous, in the two-country setup, to an internal Foreign price
(such as the nominal wage) being denominated in dollars. The general message is
that the dollar's role in goods markets (invoicing) and its role in financial
markets (safe-asset demand) interact: dollar invoicing amplifies the employment
consequences of dollar safe-asset demand, suggesting ``potentially rich
interactions between the dollar's role in financial and goods markets.''
\gapflag{The currency-of-invoicing argument is stated by KL as prose (their
p.~1666); the dollar-wage and dollar-export-pricing extensions are described but
not formally derived. We transcribe KL's argument and its mapping (dollar
invoicing of trade $\approx$ an internal Foreign price denominated in dollars, in
the two-country reduction). We do not formalize either extension; KL state them as
suggestive, citing Gopinath (2015), Gopinath et al.\ (2020), and Mukhin (2022) for
the dominant-currency-pricing channel. We match their expository level.}

\subsection{Robustness: Foreign-only demand for safe dollar bonds (App.~B.4)}
\label{sec:foreign-only}

KL's online Appendix~B.4 shows that the analysis is robust to the realistic case
in which \emph{only Foreign} households have a special (nonpecuniary) demand for
safe dollar bonds, with a non-negativity constraint preventing private agents from
\emph{creating} safe dollar assets (only the Home government can). We reconstruct
the FOC algebra, which is the load-bearing part of B.4.

\paragraph{The augmented household FOCs.} KL augment the model with a
non-negativity constraint on each household's position in the safe dollar bond,
reflecting that households cannot manufacture safe assets. This is irrelevant when
the demand shock is \emph{global} (as in the baseline), since then all agents hold
a positive share of the government's outstanding safe debt; it binds when the
demand is \emph{Foreign-only}, because Home households then have no nonpecuniary
motive to hold the safe bond and would, absent the constraint, want to short it.
The Home representative agent's FOCs for safe dollar bonds, other dollar bonds,
Foreign bonds, and capital become (their app.~p.~79)
\begin{align}
1 - \mu_t &= \Et\,\pk_{t,t+1}(1+\inom_t)\frac{P_t}{P_{t+1}}, \label{eq:b4-home-safe}\\
1 &= \Et\,\pk_{t,t+1}(1+\ioth_t)\frac{P_t}{P_{t+1}}, \label{eq:b4-home-other}\\
1 &= \Et\,\pk_{t,t+1}\frac{\rer_t}{\rer_{t+1}}(1+\fr{r}_{t+1}), \label{eq:b4-home-fbond}\\
1 &= \Et\,\pk_{t,t+1}(1+\rk_{t+1}), \label{eq:b4-home-cap}
\end{align}
where $\mu_t\ge0$ is the (scaled) multiplier on the non-negativity constraint on
the Home agent's safe-dollar-bond position, and the last three FOCs are exactly as
in the baseline model. The Foreign representative agent --- who \emph{does} receive
the nonpecuniary safe-bond utility, with the same functional form $\bondu^{*}_t$ as
in the baseline --- has FOCs (their app.~p.~79)
\begin{align}
1 - \fr c_t\,\frac{\bondu^{*}_t{}'\!\bigl(\fr B_{Ht,s}/(E_t^{-1}\fr P_t)\bigr)}
{\bondu^{*}_t\!\bigl(\fr B_{Ht,s}/(E_t^{-1}\fr P_t)\bigr)}
&= \Et\,\fr{\pk}_{t,t+1}\frac{\rer_{t+1}}{\rer_t}(1+\inom_t)\frac{P_t}{P_{t+1}}, \label{eq:b4-foreign-safe}\\
1 &= \Et\,\fr{\pk}_{t,t+1}\frac{\rer_{t+1}}{\rer_t}(1+\ioth_t)\frac{P_t}{P_{t+1}}, \label{eq:b4-foreign-other}\\
1 &= \Et\,\fr{\pk}_{t,t+1}(1+\fr{r}_{t+1}), \label{eq:b4-foreign-fbond}\\
1 &= \Et\,\fr{\pk}_{t,t+1}\frac{\rer_{t+1}}{\rer_t}(1+\rk_{t+1}), \label{eq:b4-foreign-cap}
\end{align}
as in the baseline model. Given the same functional form for $\bondu^{*}_t$ as in
the baseline (eq.~\eqref{eq:Omega-form}, starred), the Foreign marginal
nonpecuniary value evaluates to
\begin{equation}
\fr c_t\,\frac{\bondu^{*}_t{}'\!\bigl(\fr B_{Ht,s}/(E_t^{-1}\fr P_t)\bigr)}
{\bondu^{*}_t\!\bigl(\fr B_{Ht,s}/(E_t^{-1}\fr P_t)\bigr)}
= \omega^{d*}_t - \frac{\fr B_{Ht,s}}{E_t^{-1}\fr P_t\,\fr c_t}, \label{eq:b4-foreign-wedge}
\end{equation}
where $\omega^{d*}_t$ is the Foreign latent demand shock for safe dollar bonds. KL
assume for expositional simplicity that the non-negativity constraint for
\emph{Foreign} agents does not bind, which is the case when $\omega^{d*}_t$ is
sufficiently high.

\paragraph{The multiplier equals the convenience yield.} The key step. Each
agent's pair of FOCs for safe and other dollar bonds imply, by dividing the
safe-bond FOC by the other-bond FOC. For the Home agent,
\eqref{eq:b4-home-safe}/\eqref{eq:b4-home-other} gives
\begin{equation}
\frac{1+\inom_t}{1-\mu_t} = 1+\ioth_t, \label{eq:b4-home-ratio}
\end{equation}
and for the Foreign agent,
\eqref{eq:b4-foreign-safe}/\eqref{eq:b4-foreign-other} gives, using
\eqref{eq:b4-foreign-wedge},
\begin{equation}
1+\ioth_t = \frac{1+\inom_t}{1-\Bigl(\omega^{d*}_t-\dfrac{\fr B_{Ht,s}}{E_t^{-1}\fr P_t\,\fr c_t}\Bigr)}. \label{eq:b4-foreign-ratio}
\end{equation}
Both expressions equal the common gross other-dollar-bond return $1+\ioth_t$, so
their denominators must coincide:
\begin{equation}
\frac{1+\inom_t}{1-\mu_t}
= \frac{1+\inom_t}{1-\Bigl(\omega^{d*}_t-\dfrac{\fr B_{Ht,s}}{E_t^{-1}\fr P_t\,\fr c_t}\Bigr)}
\;\Longrightarrow\;
\boxed{\;\mu_t = \omega^{d*}_t - \frac{\fr B_{Ht,s}}{E_t^{-1}\fr P_t\,\fr c_t} \equiv \cy_t.\;} \label{eq:b4-mu-omega}
\end{equation}
That is, the Lagrange multiplier on the non-negativity constraint for safe dollar
bonds faced by Home agents must equal the convenience yield $\cy_t$ perceived by
Foreign agents on those bonds. Intuitively, the shadow value of being unable to
short the safe bond (for Home) is exactly the convenience premium that Foreign
agents attach to it.
\gapflag{Equations~\eqref{eq:b4-home-ratio}--\eqref{eq:b4-mu-omega} reconstruct
KL's App.~B.4 result
$\mu_t=\omega^{d*}_t-\fr B_{Ht,s}/(E_t^{-1}\fr P_t\fr c_t)\equiv\cy_t$ (their
app.~p.~79--80). KL display the FOCs~\eqref{eq:b4-home-safe}--%
\eqref{eq:b4-foreign-cap}, the Foreign wedge~\eqref{eq:b4-foreign-wedge}, the
intermediate identity
$\frac{1+i_t}{1-\mu_t}=1+\ioth_t=\frac{1+i_t}{1-(\omega^{d*}_t-\fr B_{Ht,s}/(E_t^{-1}\fr P_t\fr c_t))}$,
and conclude
$\mu_t=\omega^{d*}_t-\fr B_{Ht,s}/(E_t^{-1}\fr P_t\fr c_t)\equiv\cy_t$. We supply the
one-line division step (each agent's safe-bond FOC over its other-bond FOC
produces the gross convenience-adjusted return, and equating the two expressions
for $1+\ioth_t$ forces the denominators to coincide). The definition
$\cy_t\equiv\omega^{d*}_t-\fr B_{Ht,s}/(E_t^{-1}\fr P_t\fr c_t)$ is the Foreign-only
analog of the baseline reduced-form convenience yield~\eqref{eq:omega-reduced},
now reflecting only the Foreign latent demand $\omega^{d*}_t$ (and the supply of safe
bonds by the Home government). We reproduce KL's algebra faithfully.}

\paragraph{Consequence: the equilibrium system reduces to the baseline with
$b_{Ht,s}=0$.} With the multiplier pinned by~\eqref{eq:b4-mu-omega}, the Home
agent is at the corner --- it holds \emph{zero} safe dollar bonds, $b_{Ht,s}=0$ ---
and the Home safe-bond FOC~\eqref{eq:b4-home-safe} reads
$1-\cy_t=\Et\pk_{t,t+1}(1+\inom_t)P_t/P_{t+1}$, which is exactly the baseline
convenience-adjusted Euler equation~\eqref{eq:euler-safe} with the Home agent
holding none of the safe asset. KL conclude (their app.~p.~80):
\begin{quote}
\emph{The equilibrium system is thus exactly as in the baseline model, except with
$B_{Ht,s}=0$ (zero holdings of safe dollar bonds by Home agents). It follows that
the propagation of Foreign demand shocks for safe dollar bonds is exactly like
global demand shocks for these assets, up to any differences in seigniorage
because in this case only Foreign agents hold them.}
\end{quote}
In particular, when the Home government issues zero net safe debt ($b^g_{H,s}=0$)
there is no seigniorage, and the propagation of Foreign demand shocks is
\emph{identical} to that of global demand shocks (their footnote~55). A Foreign
demand shift propagates exactly like a baseline safety shock: at unchanged policy
rates, Foreigners seek to borrow in non-safe dollar bonds to hold safe dollar
bonds (capturing the convenience benefit without changing their currency
exposure), which raises the other-dollar-bond rate $\ioth_t$ in the absence of a
decline in the policy rate $\inom_t$, making saving more attractive for Home
households and inducing the same effects on quantities and relative prices as in
the baseline model.

\paragraph{The international risk-sharing conditions are unchanged.} KL emphasize
(their app.~p.~80) that despite only Foreigners having a special demand for safe
dollar bonds, the \emph{international risk-sharing conditions are unchanged}. The
three first-order risk-sharing relations remain
\begin{align}
\Et\,\pk_{t,t+1}\frac{1+\rr_{t+1}}{1-\cy_t}
&= \Et\,\fr{\pk}_{t,t+1}\frac{\rer_{t+1}}{\rer_t}\frac{1+\rr_{t+1}}{1-\cy_t}, \label{eq:b4-rs-safe}\\
\Et\,\pk_{t,t+1}\frac{\rer_t}{\rer_{t+1}}(1+\fr{r}_{t+1})
&= \Et\,\fr{\pk}_{t,t+1}(1+\fr{r}_{t+1}), \label{eq:b4-rs-fbond}\\
\Et\,\pk_{t,t+1}(1+\rk_{t+1})
&= \Et\,\fr{\pk}_{t,t+1}\frac{\rer_{t+1}}{\rer_t}(1+\rk_{t+1}), \label{eq:b4-rs-cap}
\end{align}
i.e.\ the same Home/Foreign kernel-equality conditions (up to the
real-exchange-rate conversion) that underlie the local-completeness
relation~\eqref{eq:risk-sharing}. With one more asset than shocks, the simplified
environment remains \emph{locally complete} as in the prior subsections, so
Propositions~\ref{prop:prices}--\ref{prop:riskpremium} carry over (with
$b_{Ht,s}=0$ on the Home side and the convenience yield now reflecting Foreign
demand only).
\gapflag{Equations~\eqref{eq:b4-rs-safe}--\eqref{eq:b4-rs-cap} transcribe KL's
displayed ``international risk sharing conditions'' (their app.~p.~80), which they
note are unchanged from the baseline. We reproduce them and the conclusion that
local completeness (one more asset than shocks) is preserved. The content
load-bearing for the big\_cy reading is exactly KL's statement: Foreign-only
safe-dollar demand propagates like a global safety shock (up to seigniorage), and
the risk-premium/portfolio results are unchanged --- so the baseline analytical
results are robust to the more realistic Foreign-only demand structure. We did not
re-derive the propagation; KL assert it follows from the system reducing to the
baseline with $B_{Ht,s}=0$, and we transcribe that reduction.}

\subsection{Distinct capital stocks and sticky prices (App.~B.5)}
\label{sec:b5}

KL's online Appendix~B.5 studies an \emph{alternative} simplified environment with
(i) \emph{distinct} national capital stocks --- Home agents choose separate
positions $k_{Ht}$ (capital used at Home) and $k_{Ft}$ (capital used at Foreign),
which offer distinct payoffs and trade at distinct prices --- and (ii) \emph{sticky
prices} (one-period-in-advance \emph{price} setting by monopolistically
competitive retailers) in place of the baseline's sticky \emph{wages}. Because
there is now trade in one more asset (two distinct capital claims rather than a
single global capital stock), one more shock is needed to pin down portfolios; KL
add a fully transitory Foreign productivity shock $z_{Ft}$ (with $\rho^{F}=0$),
analogous to the safety shock. KL provide analogs of Propositions~1--5 in this
environment but \emph{exclude the proofs} (``available on request''). We summarize
the results compactly and flag the omitted proofs.

\gapflag{This entire subsection summarizes KL's App.~B.5, which KL state
\emph{without proof} (their app.~p.~83: ``We now provide analogs of
Propositions~1--5 in this environment. We exclude the proofs for brevity but they
are available on request.''). We therefore state the analog results as KL display
them and do \emph{not} reconstruct their proofs --- a full re-derivation is out of
scope for this faithful reference and is not available in the published material.
Where a result differs qualitatively from the baseline (most importantly the
ambiguous sign of the Home capital return under sticky prices), we note the
difference and KL's stated intuition. The big\_cy extension worker should treat
B.5 as a stated-but-unproved alternative environment.}

\paragraph{Environment (B.5.1).} Relative to the baseline (Section~II): Home
households hold distinct claims $k_{Ht},k_{Ft}$ to Home- and Foreign-used capital,
with fixed national stocks $\bar k,\bar k^{*}$. Consumption of each country's
goods is a CES aggregate across varieties with elasticity $\lel$. Each country has
monopolistically competitive intermediate-good producers and final-good retailers;
retailers set prices either flexibly or one period in advance (the sticky-price
analog of the baseline's sticky wages). Households own a share of all retailers
from a given country in proportion to their ownership of that country's capital, so
aggregate profits accrue to capital owners. Capital-market clearing is now
$k_{Ht}+\zs\fr k_{Ht}=\bar k$ and $k_{Ft}+\zs\fr k_{Ft}=\bar k^{*}$. KL retain the
maintained simplifications of the baseline (e.g.\ no disaster risk or depreciation
in the analytical treatment) but \emph{do not} impose the infinite Frisch
elasticity. Because markets remain \emph{locally complete}, the results are
independent of which capital owner's stochastic discount factor is used to price
firms' dynamic decisions.

\paragraph{Proposition~1 analog (prices and quantities).} With
$\Et\dev{\wsh}_{t+1}=0$ on impact of all shocks (the symmetric-portfolio benchmark
of Propositions~\ref{prop:prices}--\ref{prop:wealth}), the transmission of a
positive safety shock is \emph{largely unchanged} from
Proposition~\ref{prop:prices}. Under flexible prices: the Home CPI declines
($\D\dev P_t=-\frac{1}{\taylor}\dev{\cy}_t$), the Home real interest rate declines
($\Et\rr_{t+1}=-\dev{\cy}_t$), and the Home real exchange rate and employment in
each country are unchanged ($\dev{\rer}_t=\dev{\ell}_t=\fr{\dev{\ell}}_t=0$). Under
prices set one period in advance: the Home CPI is unchanged ($\D\dev P_t=0$), the
Home real interest rate is unchanged ($\Et\rr_{t+1}=0$), the Home real exchange
rate \emph{appreciates} ($\dev{\rer}_t=\dev{\cy}_t$), and global employment falls,
disproportionately so in Home
($\frac{1}{1+\zs}\dev{\ell}_t+\frac{\zs}{1+\zs}\fr{\dev{\ell}}_t
=-\frac{1}{1-\cshare}\dev{\cy}_t$ and
$\dev{\ell}_t-\fr{\dev{\ell}}_t=-\frac{1}{1-\cshare}\dev{\cy}_t$). KL conclude the
transmission of safety shocks to prices and quantities is largely unchanged from
Proposition~\ref{prop:prices}.

\paragraph{Proposition~2 analog (returns), with the ambiguous-sign result.}
Defining the real returns on the two distinct capital stocks
$1+\rk_t=(\divr_t+\qk_t)/\qk_{t-1}\cdot P_{t-1}/P_t$ and the Foreign analog
$1+r^{k*}_t$, the analog of Proposition~\ref{prop:returns} on impact of a
positive safety shock is: under flexible prices, the dollar-bond real return rises
($\rr_t=\frac{1}{\taylor}\dev{\cy}_t$) while the Foreign-bond, Home-capital, and
Foreign-capital returns are unaffected (each zero net of the real-exchange-rate
change); under prices set one period in advance, the dollar-bond return is
unchanged ($\rr_t=0$), the Foreign-bond return falls
($\fr{r}_t-\D\dev{\rer}_t=-\dev{\cy}_t$), the Foreign-capital return falls
($r^{k*}_t-\dev{\rer}_t=-\dev{\cy}_t$), and --- the key new feature --- the
\emph{Home}-capital return has an \emph{ambiguous} sign:
\begin{equation}
\rk_t = \left[\frac{(1-\disc)(1-\lwedge)\bigl(\tfrac{1}{\frisch}+1\bigr)}{1-(1-\cshare)(1-\lwedge)}-1\right]\dev{\cy}_t. \label{eq:b5-rk-ambiguous}
\end{equation}
KL explain (their app.~p.~84) that the ambiguous response of the Home capital
return owes to \emph{competing effects} of a safety shock under sticky prices: the
shock \emph{reduces Home output} (pushing the capital return down), but it
\emph{raises the Home markup} $P_{Ht}/P^{i}_t$ (pushing the capital return up,
since with sticky output prices and a falling marginal cost the price--cost margin
widens). What is \emph{definitive}, however, is that the real return on Home
capital \emph{exceeds} that on Foreign capital:
\begin{equation}
\rk_t - r^{k*}_t + \dev{\rer}_t \;\propto\; \dev{\cy}_t \;>\;0 \quad\text{(on a positive safety shock).} \label{eq:b5-rk-gap}
\end{equation}
This relative-capital-return result plays the role in B.5 that the
positive-excess-return results play in the baseline: Home capital outperforms
Foreign capital on a flight to safety. KL note that the markup-rise channel --- a
positive safety shock raising relative Home equity returns through higher markups
under sticky prices --- ``can help to account for the rise in relative US equity
returns in crises documented in Dahlquist et al.\ (2023).''

\paragraph{Proposition~3 analog (wealth share / portfolios at $\rra=\rraF=1$).}
With prices set one period in advance and the unitary-risk-aversion benchmark
$\rra=\rraF=1$, the equilibrium portfolios are
\begin{equation}
\frac{\qk k_H}{\sav}=1,\quad \frac{\qk k_F}{\sav}=0,\quad
\frac{b_F}{\sav}=-\disc\,\frac{\zs}{1+\zs}\,\frac{b^g_{H,s}}{\sav},\quad
\frac{b_H}{\sav}=\disc\,\frac{\zs}{1+\zs}\,\frac{b^g_{H,s}}{\sav}, \label{eq:b5-prop3-port}
\end{equation}
i.e.\ Home fully owns its domestic capital stock and holds none of Foreign capital,
and the dollar-bond/Foreign-bond positions are the seigniorage-hedging pair. These
imply that Home's \emph{financial} wealth share $\dev{\wsh}_t$ \emph{rises} on
impact of a positive safety shock (because, by~\eqref{eq:b5-rk-gap}, Home capital
outperforms Foreign capital), \emph{falls} on a positive Foreign-productivity
shock, and is unchanged on a global-productivity shock; while Home's net foreign
assets $\nfa_t$ and \emph{expected} future wealth share $\Et\dev{\wsh}_{t+1}$ are
unchanged on impact of all shocks. KL emphasize (their footnote~58 and p.~85) that
in this $\rra=\rraF=1$ benchmark, efficient risk sharing calls for Home's
\emph{overall} wealth share (inclusive of labor income) to be \emph{constant} in
response to all shocks: with full domestic capital ownership, financial returns
exactly offset labor-income changes in response to safety and Foreign-productivity
shocks, and wealth does not redistribute across countries in response to
global-productivity shocks (building on Engel and Matsumoto 2009 and Coeurdacier
and Gourinchas 2016). A notable implication: even though Home's overall wealth
share and NFA \emph{fall} on a safety shock (Home is insuring Foreign), its
\emph{financial} wealth share $\dev{\wsh}_t$ can \emph{rise} on impact, because the
Home capital return outperforms Foreign capital.

\paragraph{Proposition~4 analog (portfolios around the benchmark).} Around the
$\rra=\rraF=1$, $b^g_{H,s}<0$ benchmark with a common positive steady-state labor
wedge, the comparative statics mirror Proposition~\ref{prop:portfolios}: Home's
portfolio share in Home capital (Foreign capital, dollar bonds) is unaffected
(unaffected, falls) with $-b^g_{H,s}$; and rises (rises, falls) with $\rraF/\rra$
holding $\rra+\frac{1}{\zs}\rraF$ fixed. That is, efficient risk sharing calls for
Home's overall wealth share to \emph{fall} on a positive safety shock as Home gets
more risk tolerant than Foreign, implemented by Home owning a \emph{levered}
portfolio of Home \emph{and} Foreign capital financed by dollar bonds.

\paragraph{Proposition~5 analog (risk premium).} KL state that the final result
characterizes the currency risk premium around the $\rra=\rraF=1$, $b^g_{H,s}<0$
benchmark and is ``essentially identical to Proposition~\ref{prop:riskpremium}'':
the covariance $\Cov_t(-\dev{\pk}_{t,t+1},\fr{r}_{t+1}-\D\dev{\rer}_{t+1}-\rr_{t+1})$
is proportional to $\rra-\rraF$ when $\sigma^{\cy}=0$ and is rising in
$\sigma^{\cy}$. The reserve-currency-paradox logic of
Proposition~\ref{prop:riskpremium} therefore survives intact in the
distinct-capital, sticky-price environment.

\paragraph{Where B.5 differs from the baseline.} Two qualitative differences are
worth recording for the extension worker. First, under sticky \emph{prices} (B.5)
the Home capital return responds \emph{ambiguously} to a safety
shock~\eqref{eq:b5-rk-ambiguous} --- competing output-decline and markup-rise
effects --- whereas under sticky \emph{wages} (baseline)
Proposition~\ref{prop:returns} gives a clean realized capital-return response. The
markup-rise channel is genuinely new and is KL's structural account of rising
relative US equity returns in crises. Second, KL note (their p.~1669) that with
sticky wages and distinct capital stocks, ``the returns to capital at Home and
Foreign are still equated in response to safety shocks'' --- capital is not
reallocated across borders even if it could be, because any increase in Foreign
output relative to Home induces a Home real-exchange-rate appreciation that fully
offsets the increase in the Foreign return to capital --- whereas the
relative-return result $\rk_t>r^{k*}_t$~\eqref{eq:b5-rk-gap} that drives the
financial-wealth-share dynamics requires the sticky-\emph{price} environment (and
holds away from the limit of complete home bias provided $\tel\neq1$). KL present
these as showing that the baseline's assumptions on capital mobility are ``not
crucial'' for the results.


% ---- 11. synthesis: load-bearing assumptions + connection to empirical facts ----
% 11_synthesis — stub (to be filled by theorist)
\bigskip\noindent\textit{[pending section]}\par


\end{document}
